% !TEX program = pdflatex
%==============================================================================
%  paper_blueprint.tex
%  Blueprint-ready refactoring of "Phase Retrieval for the One-Dimensional Cubic NLS".
%
%  PURPOSE.  This file is a Lean-friendly restructuring of paper.tex intended to be
%  turned into a `leanblueprint` project and then autoformalized against Mathlib.
%  Every mathematical object is a numbered Definition; every step is a numbered
%  Lemma/Proposition/Theorem; the dependency graph is encoded with \uses{...}.
%  It contains NO Lean code: the Lean scaffold (\lean/\leanok tags, the actual
%  Lean files) is produced in a later step. See the final appendix for the
%  suggested Lean file split.
%
%  CONVENTIONS (read before using).
%   * \uses{a,b,...} on a STATEMENT lists the definitions/spaces its statement needs.
%   * \uses{a,b,...} as the first line of a PROOF lists the results the proof invokes.
%     Together these are the blueprint DAG. They are the ONLY forward/backward links.
%   * A proof given as a short in-line argument is complete. A proof given as a
%     STEP LIST cites named sub-lemmas (each separately stated here) and ends with a
%     comment  % FULL PROOF: paper.tex L<a>-L<b>  pointing at the epsilon-level text
%     in the original manuscript, from which the formalizer can transcribe details.
%   * Target: Lean 4 / Mathlib v4.32.0, commit 81a5d257c8e410db227a6665ed08f64fea08e997.
%==============================================================================
\documentclass[11pt,oneside]{book}

\usepackage[T1]{fontenc}
\usepackage{lmodern}
\usepackage{amsmath,amssymb,amsthm,mathtools}
\usepackage[margin=1.02in,headheight=14pt]{geometry}
\usepackage{enumitem}
\usepackage{booktabs,longtable,array}
\usepackage{xcolor}
\usepackage[colorlinks=true,linkcolor=blue!55!black,citecolor=blue!55!black,urlcolor=blue!55!black]{hyperref}
\usepackage[nameinlink,capitalise,noabbrev]{cleveref}

%------------------------------------------------------------------
%  Blueprint macros.  When compiled with the real `leanblueprint`/`blueprint`
%  package these are already defined; the \providecommand fallbacks let this file
%  compile as plain LaTeX and be swapped for the blueprint package later.
%------------------------------------------------------------------
\providecommand{\uses}[1]{}
\providecommand{\lean}[1]{}
\providecommand{\leanok}{}
\providecommand{\notready}{}
\providecommand{\discussion}[1]{}
\providecommand{\mathlibok}{}

\newtheorem{theorem}{Theorem}[chapter]
\newtheorem{proposition}[theorem]{Proposition}
\newtheorem{lemma}[theorem]{Lemma}
\newtheorem{corollary}[theorem]{Corollary}
\newtheorem{claim}[theorem]{Claim}
\theoremstyle{definition}
\newtheorem{definition}[theorem]{Definition}
\newtheorem{convention}[theorem]{Convention}
\theoremstyle{remark}
\newtheorem{remark}[theorem]{Remark}
\newtheorem{checkpoint}[theorem]{Dependency checkpoint}

\newcommand{\R}{\mathbb R}
\newcommand{\C}{\mathbb C}
\newcommand{\T}{\mathbb T}
\newcommand{\N}{\mathbb N}
\newcommand{\Q}{\mathbb Q}
\newcommand{\Z}{\mathbb Z}
\newcommand{\D}{\mathcal D}
\newcommand{\iu}{\mathrm i}
\newcommand{\Pfree}{i\partial_t+\Delta}
\newcommand{\NLS}{\mathrm{NLS}}
\newcommand{\Ssch}{\mathcal S}
\newcommand{\supp}{\operatorname{supp}}
\newcommand{\1}{\mathbf 1}
\newcommand{\loc}{\mathrm{loc}}
\newcommand{\sgn}{\operatorname{sgn}}
\newcommand{\esssup}{\operatorname*{ess\,sup}}
\newcommand{\Rea}{\operatorname{Re}}
\newcommand{\Ima}{\operatorname{Im}}
\newcommand{\dd}{\,\mathrm d}
\newcommand{\norm}[1]{\left\lVert #1\right\rVert}
\newcommand{\abs}[1]{\left\lvert #1\right\rvert}
\newcommand{\pair}[2]{\left\langle #1,#2\right\rangle}
\newcommand{\FT}{\mathcal F}
\newcommand{\e}{\mathrm e}
\newcommand{\eps}{\varepsilon}
\newcommand{\weakto}{\rightharpoonup}

\title{\textbf{Phase Retrieval for the One-Dimensional Cubic NLS}\\[0.4em]
\large Blueprint edition: numbered, split, and dependency-annotated for Lean formalization}
\author{}
\date{}

\begin{document}
\frontmatter
\maketitle

\chapter*{How to use this blueprint}
\addcontentsline{toc}{chapter}{How to use this blueprint}

This document is a Lean-oriented refactoring of the manuscript \emph{Phase Retrieval for the
One-Dimensional Cubic NLS}. It is not the manuscript's exposition; it is a \emph{dependency-annotated skeleton}
whose purpose is to be compiled into a \texttt{leanblueprint} and then autoformalized.

\paragraph{What was changed relative to the source.}
\begin{itemize}
\item Every object used in a statement is now a numbered \textbf{Definition} (the source left several implicit):
      the Sobolev--Hilbert spaces $H^{s}$ and the $H^{-1/2}$--$H^{1/2}$ pairing, the local spaces $H^s_\loc$,
      the fixed-point spaces $X(I)$, $X_k(J)$, the zero-set regularizer $T_\eps$, the current $\overline f f_x$,
      the Carleman spaces $X,X',Y$, the conjugated operator $P_\beta$ and its inverse $T_\beta$, and the
      diagonal traces.
\item Five monolithic proofs were \textbf{split} into named sub-lemmas: the Fourier--Gagliardo identity, the
      $(4,\infty)$ Strichartz endpoint, the critical Wronskian lemma, the Neumann trace, and the $T_\beta$ kernel
      estimate. The Carleman graph closure and the half-space propagation are likewise broken into their steps.
\item The scalar mixed-norm layer (Module~0) and the distributional utilities (Module~12) were given their own
      early modules, which the source's module order omitted.
\item \texttt{lem:smooth-multiplier} was upgraded to a \emph{quantitative}, $C^k$-controlled, parameter-uniform
      form, as its downstream use in the Neumann trace requires.
\item The auxiliary \texttt{cor:linear-common-potential} was isolated (it is off the main path) and given the
      missing variable-coefficient integral Gr\"onwall lemma it needs.
\end{itemize}

\paragraph{The dependency graph.}
Read \verb|\uses{...}| as the blueprint edges. A statement's \verb|\uses| lists the definitions it references;
a proof's \verb|\uses| lists the results it invokes. The graph is acyclic and every proof depends only on
earlier results and on the foundational layer below.

\paragraph{Foundational layer (imported from Mathlib v4.32.0, not re-proved).}
Scalar Lebesgue/Bochner integration, Fubini--Tonelli, dominated/monotone convergence, H\"older/Minkowski/Young;
complete normed and Hilbert spaces, Banach's fixed-point theorem, Banach--Alaoglu / weak compactness; density of
$C_c^\infty$/Schwartz in $L^p$; the one-dimensional $W^{1,1}$ fundamental theorem of calculus; the $L^1$/$L^2$
Fourier transform and Plancherel; the maximum-modulus principle. \emph{Verified present} at the pinned commit:
the $L^2$ Fourier isometry \texttt{MeasureTheory.Lp.fourierTransform}$_{\ell i}$ with
\texttt{norm\_fourier\_eq}/\texttt{inner\_fourier\_eq}, the tempered-distribution Sobolev predicate
\texttt{TemperedDistribution.MemSobolev} (with \texttt{lineDerivOp}, \texttt{laplacian} at $p=2$), the
Gaussian/oscillatory Fourier integral, \texttt{AbsMax}/\texttt{PhragmenLindelof}, and the du~Bois-Reymond lemma.
No Strichartz, smoothing, Sobolev-multiplication, trace, Carleman, interpolation, maximal-function, or
unique-continuation theorem is imported: those are all \emph{proved below}.

\paragraph{One caveat carried over from the readiness audit.}
Mathlib's \texttt{MemSobolev} is a \emph{predicate} on tempered distributions, not a bundled $H^{s}$ Hilbert
space with a Gagliardo inner product and an $H^{-1/2}$--$H^{1/2}$ pairing. Module~1 therefore \emph{constructs}
the $H^{\pm1/2}$ calculus (all proofs are supplied here, chiefly \cref{thm:fourier-gagliardo}); this is the
first piece of new infrastructure to build. See the final appendix.

\tableofcontents
\mainmatter


\chapter{The equation and the solution class}\label{chap:setup}

This chapter fixes the equation and its mild formulation. The free propagator $S(t)$ and the
fixed-point space $\mathcal X(I)$ are constructed in \cref{chap:free-group,chap:wellposed}; the
global solution class is recorded as \cref{def:global-mild-solution}. The target results are
\cref{thm:main} (spacetime-a.e.\ modulus) and \cref{cor:interval-version} (any open time
interval); in the Lean target file \texttt{Showcase.lean} these are
\texttt{phase\_retrieval\_cubic\_NLS} and \texttt{phase\_retrieval\_interval}, quantified over
the structure \texttt{GlobalSolution} that mirrors \cref{def:global-mild-solution}.

\begin{definition}[The one-dimensional cubic NLS]\label{def:cubic-nls}
Fix $\sigma\in\R$. The one-dimensional cubic nonlinear Schr\"odinger equation is
\begin{equation}
 i\partial_t u+\partial_x^2 u=\sigma\abs{u}^2u,\qquad(t,x)\in\R\times\R.
 \label{eq:NLS-main}
\end{equation}
\end{definition}

\chapter{Scalar mixed norms and representatives}\label{chap:mixed-norms}

The endpoint space $L^4_tL^\infty_x$ occurs in the nonlinear flow and again in the Carleman argument. It is not treated as a Bochner space with values in the nonseparable Banach space $L^\infty$. Instead every mixed norm is defined directly on jointly measurable scalar representatives, which removes all ambiguity about sections, essential suprema, restriction, and Minkowski integration.

\section{Joint representatives and section norms}

\begin{convention}[Jointly measurable spacetime representatives]\label{conv:joint-representatives}
A scalar spacetime function on $I\times\R^d$ (with $I$ an interval) means an equivalence class of jointly Lebesgue-measurable maps $f:I\times\R^d\to\C$, where two maps are identified when they agree almost everywhere for product measure. A representative is fixed only for the duration of a calculation. By Fubini's theorem two representatives determine the same spatial section $f(t,\cdot)$ for almost every $t$, so every section norm defined below is independent of the choice of representative.
\end{convention}

For a jointly measurable $f$ and $1\le q<\infty$ put
\[
 m_qf(t)=\left(\int_{\R^d}\abs{f(t,x)}^q\,\dd x\right)^{1/q},
 \qquad
 m_\infty f(t)=\esssup_{x\in\R^d}\abs{f(t,x)}.
\]

\begin{lemma}[Measurability of the endpoint section norm]\label{lem:measurable-esssup-section}
\uses{conv:joint-representatives}
If $f:I\times\R^d\to\C$ is jointly measurable, then $t\mapsto m_\infty f(t)$ is a measurable function on $I$.
\end{lemma}

\begin{proof}
\uses{conv:joint-representatives}
Fix a strictly positive integrable weight $\omega(x)=c_d\exp(-\abs{x})$ with $\int_{\R^d}\omega=1$. For every rational $a\ge0$, Tonelli's theorem shows that
\[
 t\longmapsto A_a(t):=\int_{\R^d}\omega(x)\,\1_{\{\abs{f(t,x)}>a\}}\,\dd x
\]
is measurable. Since $\omega>0$ almost everywhere,
\[
 m_\infty f(t)>a \quad\Longleftrightarrow\quad A_a(t)>0.
\]
Hence the inverse image of $(a,\infty]$ is measurable for every rational $a\ge0$, which proves that $t\mapsto m_\infty f(t)$ is measurable.
\end{proof}

\begin{definition}[Scalar mixed norm $L^p_tL^q_x$]\label{def:scalar-mixed-norm}
\uses{conv:joint-representatives, lem:measurable-esssup-section}
For $1\le p\le\infty$ and $1\le q\le\infty$, and $f$ jointly measurable on $I\times\R^d$, define
\[
 \norm{f}_{L^p_tL^q_x(I\times\R^d)}:=\norm{m_qf}_{L^p(I)}.
\]
For finite $q$ the map $t\mapsto m_qf(t)$ is measurable by Tonelli, and for $q=\infty$ by \cref{lem:measurable-esssup-section}, so the right-hand side is well defined. The space $L^p_tL^q_x(I\times\R^d)$ is the quotient of the jointly measurable functions of finite norm by the zero-norm relation, which by Fubini is exactly equality almost everywhere on the product space. If $E\subset I\times\R^d$ is measurable, the restricted norm of $f$ to $E$ means $\norm{\1_E f}_{L^p_tL^q_x}$.
\end{definition}

\begin{lemma}[Completeness of the scalar endpoint space]\label{lem:scalar-endpoint-complete}
\uses{def:scalar-mixed-norm}
For $1\le p\le\infty$ the scalar space $L^p_tL^\infty_x(I\times\R^d)$ of \cref{def:scalar-mixed-norm} is complete.
\end{lemma}

\begin{proof}
\uses{def:scalar-mixed-norm}
Strategy: pass to a rapidly convergent subsequence, sum the section norms by scalar Minkowski, and build the limit pointwise from an absolutely convergent telescoping series.

Let $(f_n)$ be Cauchy and choose a subsequence $(f_{n_k})$ with $\norm{f_{n_{k+1}}-f_{n_k}}_{L^p_tL^\infty_x}\le2^{-k}$. Set $a_k(t)=m_\infty(f_{n_{k+1}}-f_{n_k})(t)$. Scalar Minkowski in $L^p(I)$ (in the essential-supremum sense when $p=\infty$) gives $\norm{\sum_k a_k}_{L^p_t}\le\sum_k\norm{a_k}_{L^p_t}<\infty$, so $\sum_k a_k(t)<\infty$ for almost every $t$. For each $k$, $\abs{f_{n_{k+1}}(t,x)-f_{n_k}(t,x)}\le a_k(t)$ holds for almost every $(t,x)$; deleting the countable union of exceptional sets, the telescoping series converges absolutely for almost every $(t,x)$. Define $f$ as its sum there and zero elsewhere; it is jointly measurable, and $m_\infty(f-f_{n_k})(t)\le\sum_{j\ge k}a_j(t)$ for almost every $t$, so $f_{n_k}\to f$ in the mixed norm. Since the original sequence is Cauchy, the whole sequence converges to $f$.
% FULL PROOF: paper.tex L328-L358
\end{proof}

\begin{lemma}[Restriction and scalar Minkowski]\label{lem:scalar-mixed-minkowski}
\uses{def:scalar-mixed-norm}
Let $(S,\nu)$ be a $\sigma$-finite measure space and let $F:S\times I\times\R^d\to\C$ be jointly measurable with
\[
 \int_S\norm{F(s,\cdot,\cdot)}_{L^p_tL^\infty_x}\,\dd\nu(s)<\infty .
\]
Then $H(t,x)=\int_S F(s,t,x)\,\dd\nu(s)$ is defined for almost every $(t,x)$ and
\begin{equation}
 \norm{H}_{L^p_tL^\infty_x}\le\int_S\norm{F(s)}_{L^p_tL^\infty_x}\,\dd\nu(s).
 \label{eq:scalar-minkowski-endpoint}
\end{equation}
Moreover, multiplication by the indicator of a measurable subset has norm at most one in every scalar mixed norm $L^p_tL^q_x$.
\end{lemma}

\begin{proof}
\uses{def:scalar-mixed-norm, lem:measurable-esssup-section}
Put $B(t)=\int_S m_\infty(F(s))(t)\,\dd\nu(s)$; scalar Minkowski gives $B\in L^p(I)$. After removing one product-null set by Tonelli, $\abs{F(s,t,x)}\le m_\infty(F(s))(t)$ for almost every $(s,t,x)$, so the integral defining $H$ is absolutely convergent for almost every $(t,x)$. For almost every $t$, the definition of essential supremum gives $m_\infty H(t)\le\int_S m_\infty(F(s))(t)\,\dd\nu(s)=B(t)$; taking the $L^p(I)$ norm proves \eqref{eq:scalar-minkowski-endpoint}. The restriction statement is pointwise, since $\abs{\1_E f}\le\abs{f}$ everywhere.
\end{proof}

\section{A norming predual at the spatial endpoint}

We never identify the full Banach dual of $L^4_tL^\infty_x$; only a countable norming family is needed. Fix a countable dense subset $(g_n)_{n\ge1}$ of the closed unit ball of $L^1(\R^d)$, for instance compactly supported step functions on rational boxes with coefficients in $\Q+i\Q$, normalized.

\begin{lemma}[Countable norming of $L^\infty$]\label{lem:countable-norming-Linfty}
For every $h\in L^\infty(\R^d)$,
\begin{equation}
 \norm{h}_\infty=\sup_{n\ge1}\left|\int_{\R^d}h(x)g_n(x)\,\dd x\right|.
 \label{eq:countable-norming-Linfty}
\end{equation}
\end{lemma}

\begin{proof}
\uses{}
The right side is at most $\norm{h}_\infty$ since each $g_n$ is in the unit ball of $L^1$. Conversely let $M=\norm{h}_\infty$ and $\eps>0$. The set $E=\{\abs{h}>M-\eps\}$ has positive measure; intersecting with a large ball, take $E$ of finite positive measure. The function $g=\1_E\,\overline{h}/(\abs{h}\,\abs{E})$, set to zero where $h=0$, lies in the unit ball of $L^1$ and satisfies $\abs{\int hg}\ge M-\eps$. Approximating $g$ in $L^1$ by some $g_n$ and letting $\eps\downarrow0$ gives the reverse inequality.
\end{proof}

\begin{proposition}[Endpoint norming identity]\label{prop:endpoint-norming}
\uses{def:scalar-mixed-norm, lem:countable-norming-Linfty}
Let $1<p<\infty$ and $p'=p/(p-1)$. For every jointly measurable $H\in L^p_tL^\infty_x$,
\begin{equation}
 \norm{H}_{L^p_tL^\infty_x}
 =\sup\left\{\left|\iint H(t,x)G(t,x)\,\dd x\,\dd t\right|:\ \norm{G}_{L^{p'}_tL^1_x}\le1\right\}.
 \label{eq:endpoint-norming}
\end{equation}
It suffices in the supremum to use jointly measurable $G$ which, at each time $t$, are a scalar multiple of one of the countably many $g_n$.
\end{proposition}

\begin{proof}
\uses{def:scalar-mixed-norm, lem:countable-norming-Linfty}
Strategy: Hölder gives the upper bound; for the lower bound build explicit test functions from finite maxima of the countable pairings and pass to the limit by monotone convergence.

Put $a_n(t)=\int H(t,x)g_n(x)\,\dd x$ and $M_N(t)=\max_{1\le n\le N}\abs{a_n(t)}$. Each $a_n$ is measurable by Fubini, and by \cref{lem:countable-norming-Linfty}, $M_N(t)\uparrow m_\infty H(t)$ for almost every $t$. Let $n_N(t)$ be the least index realizing the maximum (measurable level sets), let $\theta_N(t)=\overline{a_{n_N(t)}(t)}/\abs{a_{n_N(t)}(t)}$ where $M_N(t)>0$ and $0$ otherwise, and where $M_N\not\equiv0$ set
\[
 b_N(t)=\frac{M_N(t)^{p-1}}{\norm{M_N}_{L^p}^{p-1}},\qquad G_N(t,x)=b_N(t)\theta_N(t)g_{n_N(t)}(x).
\]
Then $G_N$ is jointly measurable, $\norm{G_N}_{L^{p'}_tL^1_x}\le1$, and $\iint HG_N=\int b_N(t)M_N(t)\,\dd t=\norm{M_N}_{L^p}$. Monotone convergence gives $\norm{M_N}_{L^p}\uparrow\norm{H}_{L^p_tL^\infty_x}$, proving the reverse inequality.
% FULL PROOF: paper.tex L441-L472
\end{proof}

\begin{remark}
\Cref{prop:endpoint-norming} is a norming statement, not a claim that every continuous functional on $L^{p'}_tL^1_x$ is represented by a strongly measurable $L^\infty$-valued map. This distinction is exactly what is needed at the endpoint Strichartz step.
\end{remark}

\section{Almost-everywhere uniform convergence of continuous slices}

\begin{lemma}[Mixed-norm subsequence lemma]\label{lem:mixed-subsequence-uniform}
\uses{def:scalar-mixed-norm, lem:scalar-endpoint-complete}
Let $1\le p<\infty$ and suppose $f_n\to f$ in the scalar space $L^p_tL^\infty_x(I\times\R^d)$. There is a subsequence $(f_{n_k})$ such that for almost every $t$,
\[
 \sum_{k=1}^\infty\norm{f_{n_{k+1}}(t,\cdot)-f_{n_k}(t,\cdot)}_{L^\infty_x}<\infty.
\]
If every spatial slice $f_n(t,\cdot)$ is continuous, then the essential supremum in each term equals the ordinary supremum, and the subsequence is uniformly Cauchy in $x$ for almost every $t$.
\end{lemma}

\begin{proof}
\uses{lem:scalar-endpoint-complete}
Choose the subsequence so that the $k$th mixed-norm difference is at most $2^{-k}$. The proof of \cref{lem:scalar-endpoint-complete} shows that the sum of the section essential suprema is finite for almost every $t$. A continuous function whose absolute value exceeds a number at one point exceeds it on a set of positive measure, so its supremum and essential supremum agree.
\end{proof}

\begin{lemma}[Identification of the continuous limit]\label{lem:C0-limit-identification}
\uses{lem:mixed-subsequence-uniform}
Suppose, in addition, that $f_n(t,\cdot)\in C_0(\R^d)$, that $f_n\to f$ in $C(I;L^2_x)$, and that a subsequence converges uniformly in $x$ at a given time $t$. Then its uniform limit belongs to $C_0(\R^d)$ and represents the $L^2$ class $f(t)$.
\end{lemma}

\begin{proof}
\uses{lem:mixed-subsequence-uniform}
The space $C_0$ is closed in the uniform norm, so the uniform limit lies in $C_0$. The same subsequence converges to $f(t)$ in $L^2$, while uniform convergence implies convergence locally in $L^2$. Uniqueness of the $L^2$ limit on every bounded set identifies the two limits almost everywhere, so the continuous uniform limit represents $f(t)$.
\end{proof}

\section{Sum and intersection spaces}

\begin{definition}[Sum and intersection mixed-norm spaces]\label{def:sum-intersection-norm}
\uses{def:scalar-mixed-norm}
Let $A$ and $B$ be Banach spaces continuously embedded in a common Hausdorff topological vector space (in the applications, scalar mixed-norm spaces of \cref{def:scalar-mixed-norm} inside the jointly measurable functions on $I\times\R^d$). Define
\[
 A+B=\{a+b:a\in A,\ b\in B\},\qquad \norm{h}_{A+B}=\inf_{h=a+b}\left(\norm{a}_A+\norm{b}_B\right),
\]
and
\[
 A\cap B=\{h:h\in A\text{ and }h\in B\},\qquad \norm{h}_{A\cap B}=\max\left(\norm{h}_A,\norm{h}_B\right).
\]
\end{definition}

\begin{lemma}[Completeness of sums and intersections]\label{lem:sum-intersection-complete}
\uses{def:sum-intersection-norm}
With the hypotheses of \cref{def:sum-intersection-norm}, $A+B$ and $A\cap B$ are Banach spaces. Multiplication by an indicator is contractive on $A+B$ and on $A\cap B$ whenever it is contractive on each of $A$ and $B$.
\end{lemma}

\begin{proof}
\uses{def:sum-intersection-norm}
The sum $A+B$ is the quotient of the Banach space $A\oplus_1 B$ by the closed kernel of the continuous addition map $(a,b)\mapsto a+b$, hence a Banach space. The intersection $A\cap B$ is isometric to the closed diagonal $\{(a,b)\in A\oplus_\infty B:a=b\}$, hence a Banach space. The restriction statement follows by applying the componentwise contractions before taking the infimum, respectively the maximum.
\end{proof}

\section{Joint representative for a continuous \texorpdfstring{$L^2$}{L2} path}

\begin{lemma}[Joint representative for a continuous $L^2$ path]\label{lem:joint-representative-L2}
\uses{conv:joint-representatives, def:scalar-mixed-norm}
Let $I$ be an interval and $u:I\to L^2(\R^d)$ continuous. There is a jointly measurable scalar function $U:I\times\R^d\to\C$ such that, for almost every $t$, the section $U(t,\cdot)$ represents the $L^2$ class $u(t)$. On a bounded interval $I$ the representative can be obtained as an $L^2(I\times\R^d)$ limit of jointly measurable functions that are simple in time and whose sections converge to $u(t)$ uniformly in the $L^2$ norm.
\end{lemma}

\begin{proof}
\uses{conv:joint-representatives, def:scalar-mixed-norm}
Strategy: on a bounded interval, approximate $u$ uniformly by time-simple $L^2$-valued maps, realize each by a jointly measurable representative, and identify the product-$L^2$ limit sectionwise by uniqueness of the $L^2$ limit; then exhaust a general interval.

Let $I=[a,b]$. Uniform continuity of $u$ gives, for each $n$, a finite measurable partition $I=\bigcup_j E_{n,j}$ and points $t_{n,j}\in I$ with $u_n(t)=\sum_j\1_{E_{n,j}}(t)u(t_{n,j})$ satisfying $\sup_t\norm{u_n(t)-u(t)}_2\le2^{-n}$. Pick a measurable scalar representative $g_{n,j}$ of each class $u(t_{n,j})$ and set $U_n(t,x)=\sum_j\1_{E_{n,j}}(t)g_{n,j}(x)$, jointly measurable and representing $u_n(t)$ on every section. Then $\norm{U_n-U_m}_{L^2(I\times\R^d)}\le\abs{I}^{1/2}(2^{-n}+2^{-m})$, so $U_n\to U$ in product $L^2$ with $U$ jointly measurable. Passing to a subsequence with summable successive product-$L^2$ differences, scalar Minkowski shows the sum of the section $L^2_x$ norms of those differences is finite for almost every $t$, so $U_n(t,\cdot)\to U(t,\cdot)$ in $L^2_x$ for almost every $t$. Since the classes of $U_n(t,\cdot)$ converge to $u(t)$ in $L^2$ for every $t$, uniqueness of the $L^2$ limit identifies the section of $U$ with $u(t)$ off one null set of times. For an arbitrary interval, apply this on a countable increasing exhaustion by bounded intervals; on overlaps the representatives agree product-almost everywhere because their sections represent the same $L^2$ classes, so modifying on the countable union of overlap null sets and pasting along disjoint pieces gives a global $U$.
% FULL PROOF: paper.tex L567-L600
\end{proof}


\chapter{Fourier analysis and the $H^{s}$ calculus}\label{chap:fourier}

\begin{definition}[Fourier transform convention (unitary)]\label{def:fourier-convention}
For $f\in L^1(\R^d)$ the Fourier transform is
\begin{equation}
 \widehat f(\xi)=\FT f(\xi)=(2\pi)^{-d/2}\int_{\R^d}\e^{-ix\cdot\xi}f(x)\,\dd x,
 \label{eq:FT-convention}
\end{equation}
and the inverse transform $\FT^{-1}$ uses the character $\e^{ix\cdot\xi}$ with
the same normalization $(2\pi)^{-d/2}$. The map $\FT$ extends by density to a
unitary operator on $L^2(\R^d)$, so that Plancherel holds in the form
\[
 \norm{f}_{L^2}=\norm{\widehat f}_{L^2},
 \qquad
 \pair{f}{g}_{L^2}=\pair{\widehat f}{\widehat g}_{L^2}.
\]
For $f\in\Ssch(\R^d)$ one has $\widehat{\partial_jf}=i\xi_j\widehat f$ and
$\widehat{\Delta f}=-\abs{\xi}^2\widehat f$.
\end{definition}

\begin{remark}[Rescaling to Mathlib's $\e^{-2\pi ix\eta}$ character]
\label{rem:mathlib-fourier-normalization}
\uses{def:fourier-convention}
Mathlib's standard Fourier character uses the phase $\e^{-2\pi ix\cdot\eta}$.
If $\FT_{\mathrm M}$ denotes that transform and $\FT$ is
\eqref{eq:FT-convention}, then in dimension $d$
\begin{equation}
 (\FT f)(\xi)=(2\pi)^{-d/2}(\FT_{\mathrm M}f)\!\left(\frac{\xi}{2\pi}\right).
 \label{eq:mathlib-fourier-rescaling}
\end{equation}
The change of variables $\xi=2\pi\eta$ shows the right-hand side is unitary on
$L^2$. A Lean implementation may either define this rescaled transform or retain
Mathlib's convention and replace every occurrence of $\xi$ by $2\pi\eta$; no
analytic estimate below depends on the choice.
\end{remark}

\begin{definition}[Bessel and homogeneous multipliers]\label{def:homogeneous-multiplier}
\uses{def:fourier-convention}
Write $\langle\xi\rangle=(1+\abs{\xi}^2)^{1/2}$. For $s\in\R$ define the Bessel
and homogeneous Fourier multipliers
\[
 \langle D\rangle^sf=\FT^{-1}\!\big(\langle\xi\rangle^s\,\widehat f\big),
 \qquad
 \abs{D}^sf=\FT^{-1}\!\big(\abs{\xi}^s\,\widehat f\big),
\]
initially on $\Ssch'(\R^d)$ where the products lie in $\Ssch'$. Only the values
$s=\pm\tfrac12$ and $s=1$ are used below.
\end{definition}

\begin{definition}[The Sobolev--Hilbert space $H^{s}(\R^d)$]\label{def:Hs-space}
\uses{def:fourier-convention, def:homogeneous-multiplier}
For $s\in\R$ set
\[
 H^{s}(\R^d)=\{f\in\Ssch'(\R^d):\langle\xi\rangle^s\widehat f\in L^2(\R^d)\},
 \qquad
 \norm{f}_{H^{s}}=\norm{\langle D\rangle^sf}_{L^2}
 =\Big(\int_{\R^d}\langle\xi\rangle^{2s}\abs{\widehat f(\xi)}^2\,\dd\xi\Big)^{1/2}.
\]
Equipped with the inner product
\[
 \pair{f}{g}_{H^{s}}
 =\pair{\langle D\rangle^sf}{\langle D\rangle^sg}_{L^2}
 =\int_{\R^d}\langle\xi\rangle^{2s}\,\widehat f(\xi)\overline{\widehat g(\xi)}\,\dd\xi,
\]
$H^{s}(\R^d)$ is a Hilbert space, isometric via $\FT$ to the weighted space
$L^2(\langle\xi\rangle^{2s}\dd\xi)$. In particular $H^0=L^2$ and $\langle
D\rangle^s\colon H^{s}\to L^2$ is a unitary isomorphism.
\end{definition}

\begin{definition}[The complex-bilinear $H^{-1/2}$--$H^{1/2}$ pairing]\label{def:Hs-neg-pairing}
\uses{def:Hs-space}
For $T\in H^{-1/2}(\R)$ and $\phi\in H^{1/2}(\R)$ set
\begin{equation}
 \pair{T}{\phi}_{H^{-1/2},H^{1/2}}
 =\int_{\R}\big(\langle\xi\rangle^{-1/2}\widehat T(\xi)\big)
          \big(\langle\xi\rangle^{1/2}\widehat\phi(-\xi)\big)\,\dd\xi
 =\int_{\R}\widehat T(\xi)\,\widehat\phi(-\xi)\,\dd\xi.
 \label{eq:Hs-neg-pairing}
\end{equation}
This is the unique jointly continuous \emph{complex-bilinear} extension of
$(T,\phi)\mapsto\int_\R T\phi\,\dd x$ from $C_c^\infty(\R)\times C_c^\infty(\R)$
to $H^{-1/2}\times H^{1/2}$; it is \emph{not} the sesquilinear Hilbert inner
product. By Cauchy--Schwarz and $\langle-\xi\rangle=\langle\xi\rangle$,
\begin{equation}
 \abs{\pair{T}{\phi}_{H^{-1/2},H^{1/2}}}
 \le\norm{T}_{H^{-1/2}}\norm{\phi}_{H^{1/2}}.
 \label{eq:Hs-neg-pairing-bound}
\end{equation}
\end{definition}

\begin{definition}[Local Sobolev spaces $H^{s}_{\loc}$]\label{def:Hs-loc}
\uses{def:Hs-space}
A distribution $f\in\Ssch'(\R)$ belongs to $H^{s}_{\loc}(\R)$ if
$\chi f\in H^{s}(\R)$ for every $\chi\in C_c^\infty(\R)$. Membership is
invariant under the choice of cutoff: two cutoffs equal to $1$ near the support
of a given test function produce the same distributional pairing, so all
$H^{s}_{\loc}$ statements are independent of the globalizing cutoff.
\end{definition}

Set, for $0<s<1$,
\[
 c_s=\int_{\R}\frac{\abs{\e^{ih}-1}^2}{\abs{h}^{1+2s}}\,\dd h,
\]
which is finite and strictly positive: near $0$ the numerator is $O(h^2)$ and at
infinity it is bounded.

\begin{lemma}[Exact Fourier--Gagliardo identity for Schwartz functions]\label{lem:gagliardo-fourier-identity}
\uses{def:Hs-space}
Let $0<s<1$ and $f\in\Ssch(\R)$. Then
\begin{equation}
 \iint_{\R^2}\frac{\abs{f(x)-f(y)}^2}{\abs{x-y}^{1+2s}}\,\dd x\,\dd y
 =c_s\int_\R\abs{\xi}^{2s}\abs{\widehat f(\xi)}^2\,\dd\xi.
 \label{eq:exact-gagliardo-fourier}
\end{equation}
\end{lemma}

\begin{proof}
\uses{def:fourier-convention}
Substitute $h=y-x$. By Tonelli and Plancherel,
\[
 \iint\frac{\abs{f(x+h)-f(x)}^2}{\abs{h}^{1+2s}}\,\dd x\,\dd h
 =\int_\R\abs{\widehat f(\xi)}^2
   \Big(\int_\R\frac{\abs{\e^{ih\xi}-1}^2}{\abs{h}^{1+2s}}\,\dd h\Big)\dd\xi.
\]
For $\xi\ne0$ the substitution $k=h\xi$ gives inner value $c_s\abs{\xi}^{2s}$;
at $\xi=0$ the inner integral vanishes. This is
\eqref{eq:exact-gagliardo-fourier}.
\end{proof}

\begin{lemma}[Mollification step for Gagliardo]\label{lem:gagliardo-mollification}
\uses{def:Hs-space}
Let $0<s<1$ and $f\in L^2(\R)$ have finite Gagliardo seminorm
$[f]_{\dot H^s}^2=\iint\abs{f(x)-f(y)}^2\abs{x-y}^{-1-2s}\dd x\,\dd y$. Let
$\varphi\in C_c^\infty(\R)$ with $\int\varphi=1$, $\varphi_\eps(x)=\eps^{-1}\varphi(x/\eps)$,
and $f_\eps=\varphi_\eps*f$. Then
\begin{equation}
 [f_\eps]_{\dot H^s}\le[f]_{\dot H^s},
 \label{eq:mollified-gagliardo-monotone}
\end{equation}
and $\widehat{f_\eps}(\xi)=m_\eps(\xi)\widehat f(\xi)$ with
$m_\eps(\xi)=(2\pi)^{1/2}\widehat\varphi(\eps\xi)$, satisfying $m_\eps(\xi)\to1$
pointwise as $\eps\to0$.
\end{lemma}

\begin{proof}
\uses{lem:gagliardo-fourier-identity}
Writing $f_\eps(x)-f_\eps(y)=\int\varphi_\eps(z)\big(f(x-z)-f(y-z)\big)\dd z$,
Jensen's inequality and translation invariance of the Gagliardo measure give
\eqref{eq:mollified-gagliardo-monotone}. With the unitary convention
\eqref{eq:FT-convention} the convolution theorem carries the factor
$(2\pi)^{1/2}$, giving $\widehat{f_\eps}=m_\eps\widehat f$ with the stated
symbol; since $\int\varphi=1$, $m_\eps(\xi)=(2\pi)^{1/2}\widehat\varphi(\eps\xi)\to
(2\pi)^{1/2}\widehat\varphi(0)=1$ pointwise.
\end{proof}

\begin{lemma}[Cutoff error bound $C_{\chi,s}R^{-2s}$]\label{lem:gagliardo-cutoff-error}
\uses{def:Hs-space}
Let $0<s<1$ and $\chi\in C_c^\infty(\R)$ with $0\le\chi\le1$ and $\chi=1$ near
$0$; put $\chi_R(x)=\chi(x/R)$. Then there is $C_{\chi,s}<\infty$ such that,
uniformly in $y\in\R$ and $R\ge1$,
\begin{equation}
 \int_\R\frac{\abs{\chi_R(x)-\chi_R(y)}^2}{\abs{x-y}^{1+2s}}\,\dd x
 \le C_{\chi,s}\,R^{-2s}.
 \label{eq:cutoff-gagliardo-error}
\end{equation}
\end{lemma}

\begin{proof}
Substituting $z=(x-y)/R$,
\[
 \int_\R\frac{\abs{\chi_R(x)-\chi_R(y)}^2}{\abs{x-y}^{1+2s}}\,\dd x
 =R^{-2s}\int_\R\frac{\abs{\chi(y/R+z)-\chi(y/R)}^2}{\abs{z}^{1+2s}}\,\dd z.
\]
On $\abs{z}\le1$ the numerator is $\le\norm{\chi'}_\infty^2\abs{z}^2$, giving an
integrable factor $\abs{z}^{1-2s}$; on $\abs{z}>1$ the numerator is
$\le(2\norm{\chi}_\infty)^2$ against integrable $\abs{z}^{-1-2s}$. Both bounds
are uniform in $y$, yielding \eqref{eq:cutoff-gagliardo-error}.
\end{proof}

\begin{lemma}[Density/closure step for Gagliardo]\label{lem:gagliardo-density}
\uses{def:Hs-space}
Let $0<s<1$ and $f\in L^2(\R)$.
\begin{enumerate}
\item If $\abs{\xi}^s\widehat f\in L^2$, then approximating $\widehat f$ in the
weighted norm $\int(1+\abs{\xi}^{2s})\abs{\widehat f}^2\,\dd\xi$ by
$C_c^\infty$ functions and taking inverse transforms $f_n\in\Ssch$ gives
$f_n\to f$ in $L^2$, $\abs{D}^s(f_n-f)\to0$ in $L^2$, and
$(f_n(x)-f_n(y))/\abs{x-y}^{(1+2s)/2}$ Cauchy in $L^2(\R^2)$ with limit the same
difference quotient of $f$; hence \eqref{eq:exact-gagliardo-fourier} holds for
$f$ and $[f]_{\dot H^s}<\infty$.
\item Conversely, if $[f]_{\dot H^s}<\infty$ then $\abs{\xi}^s\widehat f\in L^2$
and \eqref{eq:exact-gagliardo-fourier} holds for $f$.
\end{enumerate}
\end{lemma}

\begin{proof}
\uses{lem:gagliardo-fourier-identity, lem:gagliardo-mollification, lem:gagliardo-cutoff-error}
For (1), the Schwartz identity \eqref{eq:exact-gagliardo-fourier} applied to
$f_n-f_m$ and Plancherel show the difference quotients are Cauchy in
$L^2(\R^2)$; passing to an a.e.\ convergent subsequence identifies the limit
with that of $f$, and passing to the limit in the identity gives the claim.
For (2), let $f_\eps=\varphi_\eps*f$ and, with $\chi_R$ as in
\cref{lem:gagliardo-cutoff-error}, $g_R=\chi_Rf_\eps\in C_c^\infty\subset\Ssch$.
Splitting $g_R(x)-g_R(y)=\chi_R(x)(f_\eps(x)-f_\eps(y))+f_\eps(y)(\chi_R(x)-\chi_R(y))$,
dominated convergence handles the first summand while
\eqref{eq:cutoff-gagliardo-error} bounds the second by
$C_{\chi,s}R^{-2s}\norm{f_\eps}_2^2\to0$; thus $g_R\to f_\eps$ in $L^2$ and in
the Gagliardo seminorm. Applying \cref{lem:gagliardo-fourier-identity} to
$g_R-g_{R'}$ shows $\abs{\xi}^s\widehat{g_R}$ is Cauchy in $L^2_\xi$ with limit
$\abs{\xi}^s\widehat{f_\eps}$, and by \cref{lem:gagliardo-mollification}
combined with $[f_\eps]_{\dot H^s}\le[f]_{\dot H^s}$ we obtain
$c_s\int\abs{\xi}^{2s}\abs{\widehat{f_\eps}}^2\le[f]_{\dot H^s}^2$. Since
$\widehat{f_\eps}=m_\eps\widehat f\to\widehat f$ pointwise, Fatou's lemma yields
$\abs{\xi}^s\widehat f\in L^2$; the identity for $f$ then follows from part (1).
\end{proof}

\begin{theorem}[Exact Fourier--Gagliardo equivalence]\label{thm:fourier-gagliardo}
\uses{def:Hs-space}
For $0<s<1$ and every $f\in L^2(\R)$, the following are equivalent:
\begin{enumerate}
\item $\abs{\xi}^s\widehat f(\xi)\in L^2_\xi$;
\item $\displaystyle\iint_{\R^2}\frac{\abs{f(x)-f(y)}^2}{\abs{x-y}^{1+2s}}\,\dd x\,\dd y<\infty$.
\end{enumerate}
Whenever either holds, the exact identity
\eqref{eq:exact-gagliardo-fourier} is valid, and the norm built from
$\norm{f}_{L^2}$ and the Gagliardo seminorm is equivalent to
$\norm{f}_{H^{s}}$.
\end{theorem}

\begin{proof}
\uses{lem:gagliardo-fourier-identity, lem:gagliardo-mollification, lem:gagliardo-cutoff-error, lem:gagliardo-density}
The identity is first established for $f\in\Ssch$ in
\cref{lem:gagliardo-fourier-identity}. The two implications between (1) and (2),
together with the identity for general $f$, are exactly parts (1)--(2) of
\cref{lem:gagliardo-density}, whose proof uses the mollification monotonicity
\cref{lem:gagliardo-mollification} and the cutoff error bound
\cref{lem:gagliardo-cutoff-error}. Norm equivalence follows from
\eqref{eq:exact-gagliardo-fourier} and
$\abs{\xi}^{2s}\le\langle\xi\rangle^{2s}\le C_s(1+\abs{\xi}^{2s})$.
% FULL PROOF: paper.tex L669-L772
\end{proof}

\begin{lemma}[Quantitative smooth multipliers on $H^{\pm1/2}$]\label{lem:smooth-multiplier-quantitative}
\uses{def:Hs-space}
There is an absolute constant $C$ such that for every $\chi\in C_c^\infty(\R)$
and every $f\in H^{1/2}(\R)$,
\begin{equation}
 \norm{\chi f}_{H^{1/2}}
 \le C\big(\norm{\chi}_\infty\norm{f}_{H^{1/2}}
          +\norm{\chi'}_\infty\norm{f}_{L^2}\big).
 \label{eq:smooth-multiplier-quantitative}
\end{equation}
The same bound holds on $H^{-1/2}(\R)$ by duality. In particular the operator
norm of multiplication by $\chi$ on $H^{\pm1/2}$ is controlled by
$\norm{\chi}_\infty$ and $\norm{\chi'}_\infty$ alone, with no dependence on the
size of $\supp\chi$; hence the bound is uniform over any family
$\{\chi_\lambda\}$ with $\sup_\lambda(\norm{\chi_\lambda}_\infty+\norm{\chi_\lambda'}_\infty)<\infty$.
\end{lemma}

\begin{proof}
\uses{thm:fourier-gagliardo}
By \cref{thm:fourier-gagliardo} (with $s=\tfrac12$) it suffices to bound
$\norm{\chi f}_{L^2}$ and the Gagliardo seminorm $[\chi f]_{\dot H^{1/2}}$. The
$L^2$ term is $\le\norm{\chi}_\infty\norm{f}_2$. From
$\chi(x)f(x)-\chi(y)f(y)=\chi(x)(f(x)-f(y))+f(y)(\chi(x)-\chi(y))$ and
$(a+b)^2\le2a^2+2b^2$,
\[
 [\chi f]_{\dot H^{1/2}}^2
 \le 2\norm{\chi}_\infty^2[f]_{\dot H^{1/2}}^2
    +2\iint\frac{\abs{\chi(x)-\chi(y)}^2}{\abs{x-y}^{2}}\abs{f(y)}^2\,\dd x\,\dd y.
\]
Splitting the last integral: on $\abs{x-y}\le1$ use
$\abs{\chi(x)-\chi(y)}\le\norm{\chi'}_\infty\abs{x-y}$ to get
$\le2\norm{\chi'}_\infty^2\norm{f}_2^2$; on $\abs{x-y}>1$ use
$\abs{\chi(x)-\chi(y)}\le2\norm{\chi}_\infty$ and $\int_{\abs{x-y}>1}\abs{x-y}^{-2}\dd x=2$
to get $\le8\norm{\chi}_\infty^2\norm{f}_2^2$. Both tails are finite
independently of $\supp\chi$. Collecting terms gives
\eqref{eq:smooth-multiplier-quantitative}; the $H^{-1/2}$ statement follows by
duality of the pairing \eqref{eq:Hs-neg-pairing} with the same constant.
\end{proof}

\begin{lemma}[Translation difference quotients on $H^{1/2}$]\label{lem:translation-dq}
\uses{def:Hs-space}
For $f\in H^{1/2}(\R)$,
\[
 \frac{f(\cdot+h)-f}{h}\longrightarrow f_x
 \quad\text{strongly in }H^{-1/2}(\R)\quad(h\to0),
\]
and the difference quotients are uniformly bounded as maps
$H^{1/2}(\R)\to H^{-1/2}(\R)$.
\end{lemma}

\begin{proof}
The Fourier multiplier of the difference quotient is $(\e^{ih\xi}-1)/h$. After
multiplication by $\langle\xi\rangle^{-1/2}$ its modulus is bounded by
$C\langle\xi\rangle^{1/2}$, uniformly in $h$, giving the uniform
$H^{1/2}\to H^{-1/2}$ bound. It converges pointwise to $i\xi$, so dominated
convergence in $L^2_\xi$ against $\langle\xi\rangle^{1/2}\widehat f\in L^2$
yields the strong convergence to $f_x$.
\end{proof}

\begin{lemma}[Difference quotients in $W^{1,p}$]\label{lem:Wp-difference-quotient}
Let $1\le p<\infty$ and $f\in W^{1,p}(\R)$. Then
\begin{equation}
 \frac{f(\cdot+h)-f}{h}\longrightarrow f'
 \quad\text{strongly in }L^p(\R)\quad(h\to0),
 \label{eq:Wp-difference-quotient}
\end{equation}
and the difference quotients have $L^p$ norm at most $\norm{f'}_p$. The same
holds after integration in an auxiliary variable, in the product $L^p$ norm.
\end{lemma}

\begin{proof}
Using the locally absolutely continuous representative, for a.e.\ $x$,
$\frac{f(x+h)-f(x)}{h}=\int_0^1 f'(x+\theta h)\,\dd\theta$, so Minkowski's
integral inequality gives the uniform bound $\norm{f'}_p$. Subtracting $f'(x)$
and applying Minkowski again,
\[
 \Big\|\tfrac{\tau_hf-f}{h}-f'\Big\|_p
 \le\int_0^1\norm{\tau_{\theta h}f'-f'}_p\,\dd\theta\longrightarrow0
\]
by strong $L^p$-continuity of translations and dominated convergence in
$\theta$. The auxiliary-variable case applies the identity to a.e.\ section with
the identical domination.
% FULL PROOF: paper.tex L825-L851
\end{proof}

\begin{lemma}[One-dimensional $H^1$ continuous representative and tail bounds]\label{lem:H1-C0}
\uses{lem:Wp-difference-quotient}
Every $f\in H^1(\R)$ has a unique continuous representative in $C_0(\R)$, which
is locally absolutely continuous and satisfies
\begin{equation}
 \norm{f}_\infty^2\le2\norm{f}_2\norm{f'}_2,
 \label{eq:H1-Linfty}
\end{equation}
and, for every $x\in\R$,
\begin{align}
 \abs{f(x)}^2&\le2\norm{f}_{L^2(x,\infty)}\norm{f'}_{L^2(x,\infty)},\label{eq:H1-tail-plus}\\
 \abs{f(x)}^2&\le2\norm{f}_{L^2(-\infty,x)}\norm{f'}_{L^2(-\infty,x)}.\label{eq:H1-tail-minus}
\end{align}
\end{lemma}

\begin{proof}
The one-dimensional weak fundamental theorem of calculus provides a locally
absolutely continuous representative. Since $f\in L^2$, choose Lebesgue points
$R_n\in[n,n+1]$ with $\abs{f(R_n)}^2\le\int_n^{n+1}\abs{f}^2\to0$ and, likewise,
$L_n\in[-n-1,-n]$ with $f(L_n)\to0$. For fixed $x$ and $R_n>x$ the product rule
gives $\abs{f(R_n)}^2-\abs{f(x)}^2=2\Rea\int_x^{R_n}f'(y)\overline{f(y)}\,\dd y$;
letting $n\to\infty$ and applying Cauchy--Schwarz yields
\eqref{eq:H1-tail-plus}, and integrating from $L_n$ to $x$ gives
\eqref{eq:H1-tail-minus}. Either bound implies \eqref{eq:H1-Linfty}. The tail
norms tend to zero at the respective infinities, so the representative lies in
$C_0(\R)$; two continuous representatives agreeing a.e.\ agree everywhere.
% FULL PROOF: paper.tex L853-L887
\end{proof}

\begin{lemma}[Real part of the logarithmic derivative]\label{lem:realpart-logform}
\uses{lem:H1-C0}
Let $f\in H^1_{\loc}(\R)$. Then $\Rea\bigl(\overline f f'\bigr)=\tfrac12\partial_x\abs f^2$ as
locally integrable functions, hence as distributions. More generally, for
$f\in H^{1/2}_{\loc}(\R)$ the identity holds in the $H^{-1/2}$--$H^{1/2}$ quadratic-form sense:
for real $\phi\in C_c^\infty(\R)$, $\Rea\pair{f'}{\phi\overline f}=-\tfrac12\int\phi'\abs f^2$.
In particular the real part of the current form depends only on the density $\abs f^2$.
\end{lemma}

\begin{proof}
\uses{lem:H1-C0}
For $f\in H^1_\loc$ the product rule for the locally absolutely continuous representative
(\cref{lem:H1-C0}) gives $\partial_x\abs f^2=f'\overline f+f\overline{f'}=2\Rea(\overline f f')$.
Pairing against $-\phi'$ and moving one derivative onto the test yields the quadratic-form
identity, which extends to $f\in H^{1/2}_\loc$ by density of $H^1_\loc$ in $H^{1/2}_\loc$.
\end{proof}

\begin{lemma}[Tensor product $L^2$ norm identity]\label{lem:tensor-norm}
For $f,g\in L^2(\R)$, the function $(f\otimes g)(x,y)=f(x)g(y)$ belongs to
$L^2(\R^2)$ and
\begin{equation}
 \norm{f\otimes g}_{L^2(\R^2)}=\norm{f}_2\norm{g}_2.
 \label{eq:tensor-norm}
\end{equation}
Consequently $(f,g)\mapsto f\otimes g$ is jointly continuous from
$L^2\times L^2$ to $L^2(\R^2)$ on bounded sets. Moreover the orthogonal change
of variables $r=(x+y)/\sqrt2$, $s=(x-y)/\sqrt2$ preserves Lebesgue measure and
satisfies $\partial_x^2+\partial_y^2=\partial_r^2+\partial_s^2$.
\end{lemma}

\begin{proof}
The identity \eqref{eq:tensor-norm} is Tonelli followed by factorization
$\iint\abs{f(x)}^2\abs{g(y)}^2\,\dd x\,\dd y=\norm{f}_2^2\norm{g}_2^2$. Joint
continuity on bounded sets follows by bilinearity from
$\norm{f_1\otimes g_1-f_2\otimes g_2}_{L^2}\le\norm{f_1}_2\norm{g_1-g_2}_2+\norm{f_1-f_2}_2\norm{g_2}_2$.
The coordinate change is a rotation, hence measure-preserving, and a direct
computation gives the stated Laplacian identity.
\end{proof}


\chapter{The free Schr\"odinger group}\label{chap:free-group}

\begin{definition}[Free Schr\"odinger group $S(t)$]
\label{def:propagator}
\uses{def:fourier-convention}
For $t\in\R$ and $f\in L^2(\R)$, define $S(t)f\in L^2(\R)$ by the Fourier
multiplier
\begin{equation}
 \widehat{S(t)f}(\xi)=\e^{-it\xi^2}\widehat f(\xi),
 \qquad\xi\in\R.
 \label{eq:free-group}
\end{equation}
\end{definition}

\begin{lemma}[Unitarity, group law, and strong continuity]
\label{lem:propagator-unitary}
\uses{def:propagator}
The family $\{S(t)\}_{t\in\R}$ defined by \eqref{eq:free-group} satisfies:
\begin{enumerate}
\item[(i)] $\norm{S(t)f}_{L^2}=\norm f_{L^2}$ for all $f\in L^2(\R)$ and
      $t\in\R$ (unitarity);
\item[(ii)] $S(t)S(s)=S(t+s)$ for all $s,t\in\R$, and $S(0)=\mathrm{Id}$
      (group law);
\item[(iii)] $t\mapsto S(t)f$ is continuous from $\R$ to $L^2(\R)$ for every
      $f\in L^2(\R)$ (strong continuity);
\item[(iv)] for $f\in\Ssch(\R)$ one has $(i\partial_t+\partial_x^2)S(t)f=0$.
\end{enumerate}
\end{lemma}

\begin{proof}
\uses{def:propagator}
By \eqref{eq:free-group} the symbol $\e^{-it\xi^2}$ is unimodular, so
Plancherel gives $\norm{S(t)f}_{L^2}=\norm{\e^{-it\xi^2}\widehat f}_{L^2}
=\norm{\widehat f}_{L^2}=\norm f_{L^2}$, which is (i).  On the Fourier side
$\e^{-it\xi^2}\e^{-is\xi^2}=\e^{-i(t+s)\xi^2}$ and $\e^{0}=1$, giving (ii).
For (iii), $\norm{S(t)f-S(t_0)f}_{L^2}^2=\int_\R
|\e^{-it\xi^2}-\e^{-it_0\xi^2}|^2\,|\widehat f(\xi)|^2\,\dd\xi\to0$ as
$t\to t_0$ by dominated convergence, the integrand being bounded by
$4|\widehat f|^2\in L^1$.  For (iv), $f\in\Ssch$ makes $\e^{-it\xi^2}\widehat
f(\xi)$ and its $t$--derivative dominated by an integrable function, so
differentiation under the Fourier integral is legitimate and gives, on the
Fourier side, $\partial_t\widehat{S(t)f}=-i\xi^2\widehat{S(t)f}$ and
$\widehat{\partial_x^2 S(t)f}=-\xi^2\widehat{S(t)f}$, whence
$i\partial_t\widehat{S(t)f}+\widehat{\partial_x^2 S(t)f}=0$.
\end{proof}

\begin{lemma}[Complex Gaussian Fourier transform]
\label{lem:complex-gaussian-kernel}
Let $z\in\C$ satisfy $\Rea z>0$, and choose the holomorphic square root on
the right half-plane which is positive on $(0,\infty)$.  For every $a\in\R$,
\begin{equation}
 \int_\R \e^{-z\xi^2+ia\xi}\,\dd\xi
 =\sqrt\pi\,z^{-1/2}\exp\!\left(-\frac{a^2}{4z}\right).
 \label{eq:complex-gaussian}
\end{equation}
\end{lemma}

\begin{proof}
For $\Rea z>0$, the integral $I(z,0)=\int_\R\e^{-z\xi^2}\,\dd\xi$ is
holomorphic in $z$: on every compact subset of the right half-plane, the
integrand and its $z$--derivative are dominated by a polynomial times
$\e^{-c\xi^2}$.  For positive real $z$ the substitution $y=\sqrt z\,\xi$ and
the real Gaussian integral give $I(z,0)=\sqrt\pi\,z^{-1/2}$; the identity
theorem extends this equality to the whole right half-plane.  Fix such a $z$
and set $I_z(a)=\int_\R\e^{-z\xi^2+ia\xi}\,\dd\xi$.  Differentiation under
the integral is legitimate, and integration by parts (with vanishing
boundary terms because $\Rea z>0$) gives
\[
 0=\int_\R(-2z\xi+ia)\e^{-z\xi^2+ia\xi}\,\dd\xi,
 \qquad
 I_z'(a)=i\int_\R\xi\,\e^{-z\xi^2+ia\xi}\,\dd\xi
        =-\frac a{2z}\,I_z(a).
\]
Solving this scalar ODE and using the value $I_z(0)=\sqrt\pi\,z^{-1/2}$ at
$a=0$ proves \eqref{eq:complex-gaussian}.
\end{proof}

\begin{lemma}[Explicit free kernel]
\label{lem:free-kernel-explicit}
\uses{lem:complex-gaussian-kernel, def:propagator}
For $t\ne0$ and $f\in\Ssch(\R)$,
\begin{equation}
 S(t)f(x)=(4\pi it)^{-1/2}\int_\R
 \exp\!\left(\frac{i|x-y|^2}{4t}\right)f(y)\,\dd y,
 \label{eq:free-kernel-explicit}
\end{equation}
where the square root is the boundary value from $\Rea z>0$.
\end{lemma}

\begin{proof}
\uses{lem:complex-gaussian-kernel, def:propagator, lem:propagator-unitary}
Take $z=\eps+it$ with $\eps>0$.  Insert the Fourier definition of
$\widehat f$ into the inverse-transform representation of $S(t)f$ from
\eqref{eq:free-group}, use Fubini (the Gaussian damping makes the double
integral absolutely convergent), and apply \eqref{eq:complex-gaussian} with
$a=x-y$:
\[
 (2\pi)^{-1}\iint_{}
  \e^{i(x-y)\xi-(\eps+it)\xi^2}f(y)\,\dd y\,\dd\xi
 =(4\pi(\eps+it))^{-1/2}\int_\R
  \exp\!\left(-\frac{|x-y|^2}{4(\eps+it)}\right)f(y)\,\dd y.
\]
On the Fourier side, dominated convergence against $|\widehat f|\in L^1$
gives, as $\eps\downarrow0$, the inverse transform defining $S(t)f$,
uniformly in $x$.  On the physical side, the exponential has modulus at most
one and the prefactor stays bounded as $\eps\downarrow0$, so dominated
convergence against $|f(y)|$ gives the right side of
\eqref{eq:free-kernel-explicit}.  The boundary value of the chosen square
root fixes the phase for both signs of $t$.
\end{proof}

\begin{lemma}[Dispersive estimate]\label{lem:dispersive-1d}
\uses{lem:free-kernel-explicit, def:propagator}
For $t\ne0$ and $f\in L^1(\R)$,
\[
  \norm{S(t)f}_{L^\infty}\le (4\pi|t|)^{-1/2}\norm f_{L^1}.
\]
\end{lemma}

\begin{proof}
\uses{lem:free-kernel-explicit}
For $f\in\Ssch(\R)$ the formula \eqref{eq:free-kernel-explicit} holds, and
its kernel has constant modulus $(4\pi|t|)^{-1/2}$, so
$|S(t)f(x)|\le(4\pi|t|)^{-1/2}\int_\R|f(y)|\,\dd y$ for every $x$, which is
the claimed bound.  For general $f\in L^1(\R)$, choose $f_n\in\Ssch(\R)$ with
$f_n\to f$ in $L^1$.  The kernel integrals converge uniformly in $x$, so
they define the unique $L^1\to L^\infty$ extension of $S(t)$ and preserve the
bound $(4\pi|t|)^{-1/2}\norm f_{L^1}$.
\end{proof}

\begin{lemma}[Compatibility of the $L^2$ and kernel propagators]
\label{lem:propagator-compatibility}
\uses{def:propagator, lem:free-kernel-explicit}
For $t\ne0$, let $S_2(t)$ be the $L^2$ Fourier multiplier
\eqref{eq:free-group}, and let $S_{1,\infty}(t)$ be the integral operator
with kernel
\[
 K_t(x-y)=(4\pi it)^{-1/2}\exp\!\left(\frac{i|x-y|^2}{4t}\right).
\]
If $f\in L^1(\R)\cap L^2(\R)$, then $S_2(t)f=S_{1,\infty}(t)f$ almost
everywhere.  Thus $S_2(t)$ and $S_{1,\infty}(t)$ are compatible endpoint
realizations of one linear operator.
\end{lemma}

\begin{proof}
\uses{lem:free-kernel-explicit, lem:complex-gaussian-kernel, lem:propagator-unitary}
The equality $S_2(t)f=S_{1,\infty}(t)f$ holds for Schwartz functions by
\cref{lem:free-kernel-explicit}.  Given $f\in L^1(\R)\cap L^2(\R)$, choose
$f_n\in\Ssch(\R)$ converging to $f$ simultaneously in $L^1$ and $L^2$
(convolution followed by a spatial cutoff produces such a sequence).  The
kernel formula gives uniform convergence $S_{1,\infty}(t)f_n\to
S_{1,\infty}(t)f$, while unitarity (\cref{lem:propagator-unitary}) gives
$S_2(t)f_n\to S_2(t)f$ in $L^2(\R)$; a subsequence of the latter converges
almost everywhere.  The two limits therefore agree almost everywhere.
\end{proof}


\chapter{Interpolation and fractional integration}\label{chap:frac-int}

\begin{lemma}[Hadamard three-lines estimate]\label{lem:three-lines-explicit}
Let $F$ be bounded and continuous on the closed strip
$\{z:0\le\Rea z\le1\}$ and holomorphic in its interior. Put
\[
 M_j=\sup_{y\in\R}\abs{F(j+iy)}\qquad(j=0,1).
\]
Then, for $0\le\theta\le1$,
\begin{equation}
 \abs{F(\theta)}\le M_0^{1-\theta}M_1^{\theta}. \label{eq:three-lines-explicit}
\end{equation}
\end{lemma}

\begin{proof}
\uses{}
Strategy: rescale $F$ by an exponential factor built from $\log M_0,\log M_1$ so
that the two boundary suprema become $\le1$, insert the regularizer
$\e^{\delta z^2}$ to gain decay on horizontal sides, apply the maximum-modulus
principle on a large rectangle $[0,1]\times[-R,R]$, and let $R\to\infty$ then
$\delta\downarrow0$. The degenerate case $M_0M_1=0$ follows by replacing $M_j$
with $M_j+\eps$ and letting $\eps\downarrow0$.
% FULL PROOF: paper.tex L1024-L1055
\end{proof}

\begin{lemma}[Riesz--Thorin interpolation in the form used here]\label{lem:riesz-thorin-used}
\uses{}
Let a linear operator $T$ be defined on simple functions and have compatible
extensions
\[
 T:L^{p_0}\to L^{q_0},\quad \norm{T}\le M_0,
 \qquad\text{and}\qquad
 T:L^{p_1}\to L^{q_1},\quad \norm{T}\le M_1,
\]
where $1\le p_j<\infty$ and $1\le q_j\le\infty$. If $0<\theta<1$ and
\[
 \frac1{p_\theta}=\frac{1-\theta}{p_0}+\frac{\theta}{p_1},
 \qquad
 \frac1{q_\theta}=\frac{1-\theta}{q_0}+\frac{\theta}{q_1},
\]
then
\[
 \norm{Tf}_{q_\theta}
 \le M_0^{1-\theta}M_1^{\theta}\norm{f}_{p_\theta}.
\]
\end{lemma}

\begin{proof}
\uses{lem:three-lines-explicit}
Strategy: reduce to simple $f,g$ normalized by
$\norm{f}_{p_\theta}=\norm{g}_{q_\theta'}=1$; form the analytic family
$f_z=\abs{f}^{p_\theta\alpha(z)}f/\abs{f}$ and
$g_z=\abs{g}^{q_\theta'\gamma(z)}g/\abs{g}$ with
$\alpha(z)=\tfrac{1-z}{p_0}+\tfrac{z}{p_1}$,
$\gamma(z)=\tfrac{1-z}{q_0'}+\tfrac{z}{q_1'}$, and consider
$\Phi(z)=\int (Tf_z)g_z$. On $\Rea z=j$ one has
$\norm{f_z}_{p_j}=\norm{g_z}_{q_j'}=1$, so $\abs{\Phi(j+iy)}\le M_j$; the
three-lines estimate \eqref{eq:three-lines-explicit} gives
$\abs{\Phi(\theta)}\le M_0^{1-\theta}M_1^{\theta}$, and $f_\theta=f$,
$g_\theta=g$. Taking the supremum over dual test functions yields the claim.
When $q_j=\infty$ the boundary bound is the $L^\infty$--$L^1$ H\"older
inequality with $1/q_j=0$; input exponents are finite in the stated form.
% FULL PROOF: paper.tex L1080-L1111
\end{proof}

\begin{definition}[Centered Hardy--Littlewood maximal operator on the line]\label{def:maximal-function}
\uses{}
For $F\in L^1_{\loc}(\R)$ the centered Hardy--Littlewood maximal function is
\[
 MF(t)=\sup_{R>0}\frac1{2R}\int_{t-R}^{t+R}\abs{F(s)}\,\dd s.
\]
\end{definition}

\begin{lemma}[Weak-$(1,1)$ bound for the maximal operator]\label{lem:maximal-weak11}
\uses{def:maximal-function}
For every $\lambda>0$ and every $F\in L^1(\R)$,
\begin{equation}
 \big|\{MF>\lambda\}\big|\le\frac{C}{\lambda}\norm{F}_{1},
 \label{eq:maximal-weak11}
\end{equation}
with an absolute constant $C$ (one may take $C=3$).
\end{lemma}

\begin{proof}
\uses{def:maximal-function}
Let $K\subset\{MF>\lambda\}$ be compact. For each $t\in K$ pick an interval
centered at $t$ on which the average of $\abs{F}$ exceeds $\lambda$; by
compactness finitely many of these cover $K$. Select from them greedily, in
decreasing order of length, a pairwise disjoint subfamily $\{I_j\}$. Every
unselected interval meets a longer selected one and is therefore contained in
its triple, so the triples cover $K$. Hence
\[
 \abs{K}\le3\sum_j\abs{I_j}
 <\frac3\lambda\sum_j\int_{I_j}\abs{F}
 \le\frac3\lambda\norm{F}_1,
\]
using disjointness of the $I_j$. Inner regularity gives
\eqref{eq:maximal-weak11}.
\end{proof}

\begin{lemma}[Strong-$(p,p)$ bound for the maximal operator]\label{lem:maximal-strongp}
\uses{def:maximal-function}
For every $1<p<\infty$ there is $C_p$ with
\begin{equation}
 \norm{MF}_{p}\le C_p\norm{F}_{p}\qquad(F\in L^p(\R)).
 \label{eq:maximal-strongp}
\end{equation}
\end{lemma}

\begin{proof}
\uses{lem:maximal-weak11}
Fix $\lambda>0$ and split
$F=F\1_{\{\abs{F}>\lambda/2\}}+F\1_{\{\abs{F}\le\lambda/2\}}$. The maximal
function of the second summand is at most $\lambda/2$, so
$\{MF>\lambda\}\subset\{M(F\1_{\{\abs F>\lambda/2\}})>\lambda/2\}$, and the
weak-$(1,1)$ bound \eqref{eq:maximal-weak11} yields
\[
 \abs{\{MF>\lambda\}}
 \le\frac{C}{\lambda}\int_{\{\abs{F}>\lambda/2\}}\abs{F(s)}\,\dd s.
\]
Multiply by $p\lambda^{p-1}$, integrate in $\lambda\in(0,\infty)$, and apply
Tonelli:
\begin{align*}
 \norm{MF}_p^p
 &\le Cp\int_0^\infty\lambda^{p-2}
       \int_{\{\abs{F}>\lambda/2\}}\abs{F(s)}\,\dd s\,\dd\lambda\\
 &=Cp\int_\R\abs{F(s)}
      \int_0^{2\abs{F(s)}}\lambda^{p-2}\,\dd\lambda\,\dd s
 \le C_p\norm{F}_p^p.
\end{align*}
\end{proof}

\begin{lemma}[Hedberg pointwise inequality]\label{lem:hedberg}
\uses{def:maximal-function}
Let $0<\alpha<1$ and $1<p<q<\infty$ with $1/q=1/p-(1-\alpha)$. Then for
$F\in L^p(\R)$ and every $t$,
\[
 I_{1-\alpha}F(t):=\int_\R\abs{t-s}^{-\alpha}\abs{F(s)}\,\dd s
 \le C\,(MF(t))^{p/q}\norm{F}_{p}^{1-p/q}.
\]
\end{lemma}

\begin{proof}
\uses{def:maximal-function}
Fix $R>0$ and split $I_{1-\alpha}F(t)$ at $\abs{t-s}=R$. Dyadic decomposition
of the near part gives
\[
 \int_{\abs{t-s}<R}\abs{t-s}^{-\alpha}\abs{F(s)}\,\dd s
 \le C R^{1-\alpha}MF(t),
\]
while H\"older's inequality on the far part gives
\[
 \int_{\abs{t-s}\ge R}\abs{t-s}^{-\alpha}\abs{F(s)}\,\dd s
 \le C R^{1-\alpha-1/p}\norm{F}_{p}.
\]
Choosing $R=(\norm{F}_p/MF(t))^{p}$ when $MF(t)>0$ (the cases $MF(t)=0$ and
$\norm{F}_p=0$ being immediate) and using $p(1-\alpha)=1-p/q$ balances the two
terms and yields the stated inequality.
\end{proof}

\begin{lemma}[Hardy--Littlewood--Sobolev on the time line]\label{lem:HLS-time}
\uses{}
Let $0<\alpha<1$, $1<p<q<\infty$, and $1/q=1/p-(1-\alpha)$. Then
\[
 \left\|\int_\R\abs{t-s}^{-\alpha}F(s)\,\dd s\right\|_{L^q_t}
 \le C_{p,\alpha}\norm{F}_{L^p_t}.
\]
\end{lemma}

\begin{proof}
\uses{lem:maximal-strongp, lem:hedberg}
By the Hedberg inequality \cref{lem:hedberg},
$I_{1-\alpha}F(t)\le C(MF(t))^{p/q}\norm{F}_p^{1-p/q}$ pointwise. Raising to the
$q$th power and integrating in $t$,
\[
 \norm{I_{1-\alpha}F}_{L^q_t}^q
 \le C^q\norm{F}_p^{q-p}\int_\R (MF(t))^p\,\dd t
 =C^q\norm{F}_p^{q-p}\norm{MF}_p^{p}.
\]
The strong-$(p,p)$ bound \cref{lem:maximal-strongp} gives
$\norm{MF}_p^p\le C_p^p\norm{F}_p^p$, whence
$\norm{I_{1-\alpha}F}_{L^q_t}^q\le C^q C_p^p\norm{F}_p^q$. Taking $q$th roots
proves the estimate.
% FULL PROOF: paper.tex L1139-L1204
\end{proof}


\chapter{Strichartz estimates}\label{chap:strichartz}

\section{Homogeneous estimates}

\begin{definition}[Admissible pairs]\label{def:admissible-pairs}
A pair $(q,r)$ is called one-dimensional Schr\"odinger \emph{admissible} if
\[
  2\le q,r\le\infty,
  \qquad \frac2q+\frac1r=\frac12.
\]
\end{definition}

We split the homogeneous estimate into the finite-exponent $TT^*$ bound, the scalar endpoint $TT^*$ bound, and the almost-everywhere identification at the spatial endpoint.

\begin{lemma}[$TT^*$ finite-exponent estimate]\label{lem:TTstar-finite}
\uses{def:admissible-pairs, def:propagator, def:scalar-mixed-norm}
Let $(q,r)$ be admissible with $2<q<\infty$ and $r<\infty$. For smooth compactly supported spacetime functions $F$ put
\[
 TT^*F(t)=\int_\R S(t-s)F(s)\,\dd s .
\]
Then
\[
 \norm{TT^*F}_{L^q_tL^r_x}
 \le C\norm F_{L^{q'}_tL^{r'}_x},
\]
where $q'$ and $r'$ are the H\"older conjugates of $q$ and $r$.
\end{lemma}

\begin{proof}
\uses{lem:dispersive-1d, lem:HLS-time, lem:propagator-compatibility}
By the dispersive estimate of \cref{lem:dispersive-1d} interpolated with the unitary $L^2$ bound, together with the scalar Minkowski inequality,
\[
 \norm{TT^*F(t)}_{L^r_x}
 \le C\int_\R|t-s|^{-(1/2-1/r)}
              \norm{F(s)}_{L^{r'}_x}\,\dd s .
\]
Admissibility gives $1/2-1/r=2/q$. The time-line Hardy--Littlewood--Sobolev estimate \cref{lem:HLS-time} with $p=q'$ and $\alpha=2/q$ then yields
\[
 \norm{TT^*F}_{L^q_tL^r_x}
 \le C\norm F_{L^{q'}_tL^{r'}_x}. \qedhere
\]
\end{proof}

\begin{lemma}[$TT^*$ scalar endpoint estimate]\label{lem:TTstar-endpoint}
\uses{def:propagator, def:scalar-mixed-norm}
For smooth compactly supported spacetime functions $F$, with
$TT^*F(t)=\int_\R S(t-s)F(s)\,\dd s$ and
$T^*F=\int_\R S(-s)F(s)\,\dd s$,
\[
 \norm{TT^*F}_{L^4_tL^\infty_x}
 \le C\norm F_{L^{4/3}_tL^1_x},
\]
and consequently
\[
 \norm{T^*F}_2\le C\norm F_{L^{4/3}_tL^1_x}.
\]
\end{lemma}

\begin{proof}
\uses{lem:dispersive-1d, lem:HLS-time, prop:endpoint-norming}
For smooth compactly supported $F$ the dispersive estimate \cref{lem:dispersive-1d} gives, with $m_\infty$ the endpoint section norm of \cref{prop:endpoint-norming},
\[
 m_\infty(TT^*F)(t)
 \le C\int_\R|t-s|^{-1/2}\norm{F(s)}_{L^1_x}\,\dd s .
\]
The time-line Hardy--Littlewood--Sobolev estimate \cref{lem:HLS-time} with $p=4/3$ and $\alpha=1/2$ yields the first bound. The quadratic identity
\[
 \norm{T^*F}_2^2
 =\left|\iint (TT^*F)(t,x)\overline{F(t,x)}\,\dd x\,\dd t\right|
 \le \norm{TT^*F}_{L^4_tL^\infty_x}\norm F_{L^{4/3}_tL^1_x}
 \le C\norm F_{L^{4/3}_tL^1_x}^2
\]
then gives the second bound.
\end{proof}

\begin{lemma}[Spatial-endpoint almost-everywhere identification]\label{lem:strichartz-endpoint-identification}
\uses{def:propagator, def:scalar-mixed-norm}
Let $f\in L^2(\R)$ and let $f_n$ be Schwartz with $f_n\to f$ in $L^2$. Suppose the endpoint bound $\norm{S(t)g}_{L^4_tL^\infty_x}\le C\norm g_2$ holds for Schwartz $g$, so that $S(t)f_n$ is Cauchy in the complete scalar space $L^4_tL^\infty_x$ and converges there to a jointly measurable $H$. Then $H=S(t)f$ for almost every $(t,x)$, and
\[
 \norm{S(t)f}_{L^4_tL^\infty_x}\le C\norm f_2 .
\]
\end{lemma}

\begin{proof}
\uses{lem:mixed-subsequence-uniform}
Fix a bounded time interval $I$. Unitarity gives
\[
 \int_I\norm{S(t)(f_n-f)}_2^2\,\dd t
 =|I|\,\norm{f_n-f}_2^2\longrightarrow0 .
\]
Choose a subsequence whose successive differences are summable both in $L^4_tL^\infty_x$ and in $L^2(I\times\R)$. By \cref{lem:mixed-subsequence-uniform}, for almost every $t$ this subsequence converges in the spatial essential-supremum norm to the section $H(t)$. By the ordinary summable-subsequence argument in $L^2(I\times\R)$, after a further subsequence it converges pointwise almost everywhere to the canonical spacetime representative of $S(t)f$; the further subsequence retains the sectionwise uniform convergence. Hence the two limits agree for almost every $(t,x)$, identifying $H$ with $S(t)f$ on $I$. A countable exhaustion of the time axis makes the identification global, and the norm bound passes to the limit.
\end{proof}

\begin{theorem}[Homogeneous Strichartz estimates]\label{thm:hom-strichartz}
\uses{def:admissible-pairs, def:propagator}
For every admissible pair $(q,r)$,
\[
  \norm{S(t)f}_{L^q_t(\R;L^r_x)}\le C_{q,r}\norm f_2 .
\]
In particular,
\[
 \norm{S(t)f}_{L^6_{t,x}}+\norm{S(t)f}_{L^4_tL^\infty_x}
 \le C\norm f_2 .
\]
\end{theorem}

\begin{proof}
\uses{lem:TTstar-finite, lem:TTstar-endpoint, lem:strichartz-endpoint-identification, lem:dispersive-1d, prop:endpoint-norming, lem:propagator-unitary}
The proof splits the admissible pairs into three regimes handled by the three preceding lemmas.

\emph{Finite pairs $2<q<\infty$, $r<\infty$.} On compactly supported smooth spacetime functions define $Tf(t)=S(t)f$ and $T^*F=\int_\R S(-s)F(s)\,\dd s$, an $L^2$ Bochner integral. By \cref{lem:TTstar-finite},
\[
 \norm{TT^*F}_{L^q_tL^r_x}\le C\norm F_{L^{q'}_tL^{r'}_x}.
\]
The quadratic identity $\norm{T^*F}_2^2=|\iint (TT^*F)\overline F\,\dd x\,\dd t|\le C\norm F_{L^{q'}L^{r'}}^2$ shows $T^*$ extends boundedly to $L^{q'}L^{r'}$, and finite-exponent duality yields the estimate.

\emph{Spatial endpoint $(q,r)=(4,\infty)$.} By \cref{lem:TTstar-endpoint}, $\norm{T^*F}_2\le C\norm F_{L^{4/3}_tL^1_x}$. For Schwartz $f$, Fourier inversion gives $\norm{S(t)f}_\infty\le(2\pi)^{-1/2}\norm{\widehat f}_1$ for $|t|\le1$, while \cref{lem:dispersive-1d} gives $\norm{S(t)f}_\infty\le C|t|^{-1/2}\norm f_1$ for $|t|>1$, whose fourth power is integrable at infinity; hence $S(t)f\in L^4_tL^\infty_x$ in the scalar sense. For smooth compactly supported $G$,
\[
 \left|\iint S(t)f(x)\overline{G(t,x)}\,\dd x\,\dd t\right|
 =|\pair f{T^*G}_{L^2}|
 \le C\norm f_2\norm G_{L^{4/3}_tL^1_x}.
\]
Smooth compactly supported functions are dense in the norming class of \cref{prop:endpoint-norming}; approximation in $L^{4/3}_tL^1_x$ and the norming identity give $\norm{S(t)f}_{L^4_tL^\infty_x}\le C\norm f_2$ for Schwartz $f$. \cref{lem:strichartz-endpoint-identification} extends this to all $f\in L^2$.

\emph{Endpoint $(q,r)=(\infty,2)$} is unitarity, \cref{lem:propagator-unitary}. These regimes exhaust the admissible pairs.
% FULL PROOF: paper.tex L1225-L1320
\end{proof}

\section{Retarded estimates}

\begin{lemma}[Christ--Kiselev time ordering]\label{lem:CK}
Let $1\le p<q\le\infty$ and let $K:L^p(I;A)\to L^q(I;B)$ be a bounded integral operator
\[
 KF(t)=\int_I K(t,s)F(s)\,\dd s
\]
between Banach-valued function spaces. Then the retarded operator
\[
 K_{<}F(t)=\int_{s<t}K(t,s)F(s)\,\dd s
\]
is bounded from $L^p(I;A)$ to $L^q(I;B)$, with norm depending only on $p,q$ and the norm of $K$.
\end{lemma}

\begin{proof}
Normalize $\norm F_{L^p}=1$ and set $\Phi(t)=\int_{I\cap(-\infty,t)}\norm{F(s)}_A^p\,\dd s$; up to null sets $\Phi$ is nondecreasing with values in $[0,1]$. At generation $n$ pair the sibling dyadic intervals
\[
 E^-_{n,j}=[2j2^{-n},(2j+1)2^{-n}),
 \qquad
 E^+_{n,j}=[(2j+1)2^{-n},(2j+2)2^{-n}),
\]
and put $I^\pm_{n,j}=\Phi^{-1}(E^\pm_{n,j})$. Every pair $(s,t)$ with $s<t$ and $\Phi(s)<\Phi(t)$ belongs to exactly one rectangle $I^-_{n,j}\times I^+_{n,j}$, namely at the first binary generation where the two values lie in different children; the remaining pairs lie on level sets of $\Phi$, where $F=0$ almost everywhere. Thus $K_<$ is the sum over $n,j$ of $\1_{I^+_{n,j}}(t)\int_{I^-_{n,j}}K(t,s)F(s)\,\dd s$. For fixed $n$ the output sets are disjoint and each input piece has $L^p$ norm at most $2^{-n/p}$, so if the full operator has norm $M$ the $L^q$ norm of the $n$th generation is at most
\[
 M\left(\sum_j 2^{-nq/p}\right)^{1/q}
 \le CM2^{n/q}2^{-n/p}.
\]
When $q=\infty$, disjointness is replaced by the maximum over $j$ and the bound reads $M2^{-n/p}$ since $2^{n/q}=1$. Summing in $n$ gives the convergent geometric series $\sum_n2^{-n(1/p-1/q)}$.
\end{proof}

\begin{theorem}[Inhomogeneous estimates]\label{thm:inhom-str}
\uses{def:propagator}
For every interval $I$ and every $F\in L^{6/5}(I\times\R)$,
\begin{align*}
 \left\|\int_{t_0}^tS(t-s)F(s)\,\dd s\right\|_{L^6(I\times\R)}
 &\le C\norm F_{L^{6/5}(I\times\R)},\\
 \left\|\int_{t_0}^tS(t-s)F(s)\,\dd s\right\|_{L^\infty_tL^2_x(I)}
 &\le C\norm F_{L^{6/5}(I\times\R)}.
\end{align*}
The constants do not depend on $I$ or $t_0\in I$.
\end{theorem}

\begin{proof}
\uses{thm:hom-strichartz, lem:CK}
The full convolution estimate from $L^{6/5}_{t,x}$ to $L^6_{t,x}$ is the $TT^*$ estimate underlying \cref{thm:hom-strichartz} (the admissible pair $(6,6)$). Apply \cref{lem:CK} with time exponents $6/5<6$ to restrict to $s<t$; reversing time handles $t<t_0$, giving the first estimate. For the second, duality and the homogeneous $(6,6)$ estimate give $\|\int_I S(-s)F(s)\,\dd s\|_2\le C\norm F_{6/5}$, so the full operator $F\mapsto S(t)\int_I S(-s)F(s)\,\dd s$ maps $L^{6/5}_{t,x}$ to $L^\infty_tL^2_x$. Apply \cref{lem:CK} with input exponent $6/5$ and output exponent $\infty$ to restrict to $s<t$; reversing time handles the other orientation around $t_0$.
\end{proof}

\begin{lemma}[Continuous retarded operator, base time, Bochner compatibility]\label{lem:retarded-continuous}
\uses{thm:inhom-str, def:sum-intersection-norm, def:propagator}
Let $I$ be a bounded interval and $t_0\in I$. For smooth compactly supported $F$ put
\[
 (\mathcal R_{t_0}F)(t)=\int_{t_0}^tS(t-s)F(s)\,\dd s.
\]
This operator extends uniquely to a bounded map
\[
 \mathcal R_{t_0}:L^{6/5}(I\times\R)
 \longrightarrow C(I;L^2(\R))\cap L^6(I\times\R),
\]
with the estimates of \cref{thm:inhom-str}. If $t_1\in I$, then
\[
 \mathcal R_{t_0}F(t)
 =\mathcal R_{t_1}F(t)+S(t-t_1)\mathcal R_{t_0}F(t_1)
\]
for every $t\in I$. If in addition $F\in L^1(I;L^2)\cap L^{6/5}(I\times\R)$, the completed operator is the ordinary $L^2$-valued Bochner integral displayed above.
\end{lemma}

\begin{proof}
\uses{thm:inhom-str, lem:propagator-unitary, lem:scalar-mixed-minkowski}
For smooth compactly supported $F$ the displayed integral is an ordinary Bochner integral in $L^2$; strong continuity of $S(t)$ and dominated convergence make it continuous in $t$. Let $F_n$ be smooth compactly supported and Cauchy in $L^{6/5}$. By \cref{thm:inhom-str}, $\mathcal R_{t_0}F_n$ is Cauchy in both $L^6_{t,x}$ and $L^\infty_tL^2_x$; since the difference of two outputs is a continuous $L^2$-valued function, the essential supremum of its $L^2$ norm equals the ordinary supremum, so the sequence is uniformly Cauchy in $C(I;L^2)$ and converges there. Density and the two inhomogeneous estimates give existence, uniqueness, and the mapping property.

For smooth forcing, splitting the oriented integral at $t_1$ and using the group law gives
\[
 \int_{t_0}^tS(t-s)F(s)\,\dd s
 =\int_{t_1}^tS(t-s)F(s)\,\dd s
   +S(t-t_1)\int_{t_0}^{t_1}S(t_1-s)F(s)\,\dd s .
\]
Approximation in $L^{6/5}$, continuity of evaluation on $C(I;L^2)$, and strong continuity of $S(t-t_1)$ extend the base-time identity to the completed domain.

Finally suppose $F\in L^1_tL^2_x$ as well. Approximate it by smooth compactly supported functions in the intersection norm $\norm\cdot_{L^{6/5}}+\norm\cdot_{L^1L^2}$ of \cref{def:sum-intersection-norm}. Unitarity (\cref{lem:propagator-unitary}) and scalar Minkowski (\cref{lem:scalar-mixed-minkowski}) give
\[
 \sup_{t\in I}\left\|\int_{t_0}^tS(t-s)(F_n-F)(s)\,\dd s\right\|_2
 \le\norm{F_n-F}_{L^1_tL^2_x},
\]
so the Bochner integrals converge uniformly in $L^2$ to the same continuous limit as the Strichartz completion.
\end{proof}


\chapter{Global $L^2$ well-posedness}\label{chap:wellposed}

\begin{definition}[Nonlinear fixed-point space $\mathcal X(I)$ and the map $\Phi_{u_0}$]
\label{def:fixed-point-space}
\uses{def:scalar-mixed-norm}
Let $I=[t_0-T,t_0+T]$ be a bounded interval. Define
\[
  \mathcal X(I)=C(I;L^2(\R))\cap L^6(I\times\R),
  \qquad
  \norm u_{\mathcal X(I)}
  =\norm u_{L^\infty_tL^2_x(I)}+\norm u_{L^6_{t,x}(I)}.
\]
This is a Banach space: it is the closed diagonal subspace of
$C(I;L^2)\oplus_1L^6(I\times\R)$ under the identification of the two
components as distributions, equality of the distributional limits following
from testing against compactly supported smooth functions.
For initial data $u_0\in L^2(\R)$ and a fixed coupling $\sigma\in\R$, the
associated nonlinear map is
\begin{equation}
 \Phi_{u_0}(u)(t)
 =S(t-t_0)u_0-i\sigma\int_{t_0}^tS(t-s)(|u|^2u)(s)\,\dd s,  \label{eq:Phi}
\end{equation}
where the integral denotes the completed retarded operator
$\mathcal R_{t_0}$ of \cref{lem:retarded-continuous}.
\end{definition}

\begin{definition}[Mild solution and the Duhamel formula]\label{def:mild-solution}
\uses{def:cubic-nls, def:propagator}
A function $u\in C(I;L^2(\R))$ is a \emph{mild solution} of \eqref{eq:NLS-main} on an interval
$I$ if, for one and hence every $t_0\in I$,
\begin{equation}
 u(t)=S(t-t_0)u(t_0)-i\sigma\int_{t_0}^t S(t-s)\bigl(\abs{u}^2u\bigr)(s)\,\dd s
 \quad\text{in }L^2(\R).
 \label{eq:mild-main}
\end{equation}
\end{definition}

\begin{definition}[Global mild solution class]\label{def:global-mild-solution}
\uses{def:mild-solution, def:fixed-point-space, def:scalar-mixed-norm, lem:joint-representative-L2}
A \emph{global mild solution} of \eqref{eq:NLS-main} is a curve $u\in C(\R;L^2(\R))$ that, on
every bounded interval $I$, lies in $\mathcal X(I)$ (in particular $u\in L^6_{\loc}(\R^2)$), has
cubic forcing $\abs{u}^2u\in L^1_{\loc}(\R;L^2(\R))$, and satisfies the Duhamel identity
\eqref{eq:mild-main} for all $t,t_0\in\R$. This is the class formalized as the structure
\texttt{GlobalSolution} in \texttt{Showcase.lean}; \cref{thm:GWP} shows that every $u_0\in L^2$
generates a unique such solution.
\end{definition}

\begin{lemma}[Cubic estimate $L^{6/5}\to$]\label{lem:cubic-65}
\uses{def:fixed-point-space}
For $u\in\mathcal X(I)$,
\begin{equation}
 \norm{|u|^2u}_{L^{6/5}_{t,x}(I)}
 \le |I|^{1/2}\,\norm u_{L^\infty_tL^2_x(I)}\,
                    \norm u_{L^6_{t,x}(I)}^2.       \label{eq:cubic65}
\end{equation}
\end{lemma}

\begin{proof}
\uses{def:scalar-mixed-norm}
At each fixed time, H\"older in space with $\tfrac56=\tfrac12+\tfrac16+\tfrac16$
gives $\norm{|u|^2u}_{L^{6/5}_x}\le\norm{u}_{2}\,\norm{u}_{6}^2$. The function
$t\mapsto\norm{u(t)}_6^2$ belongs to $L^3(I)$, and the embedding
$L^3(I)\hookrightarrow L^{6/5}(I)$ costs the factor
$|I|^{5/6-1/3}=|I|^{1/2}$. This proves \eqref{eq:cubic65}.
\end{proof}

\begin{lemma}[Lipschitz cubic difference estimate]\label{lem:cubic-lip}
\uses{def:fixed-point-space, lem:cubic-65}
There is an absolute constant $C$ such that for all $u,v\in\mathcal X(I)$,
\begin{equation}
 \norm{|u|^2u-|v|^2v}_{L^{6/5}_{t,x}(I)}
 \le C|I|^{1/2}
  \bigl(\norm u_{\mathcal X(I)}^2+\norm v_{\mathcal X(I)}^2\bigr)
  \norm{u-v}_{\mathcal X(I)}.                      \label{eq:cubic-lip}
\end{equation}
\end{lemma}

\begin{proof}
\uses{lem:cubic-65}
Use the pointwise bound $\bigl||u|^2u-|v|^2v\bigr|\le C(|u|^2+|v|^2)|u-v|$.
At each time assign one of the three factors to $L^2_x$ and the remaining two
to $L^6_x$, so the spatial exponents again add to $5/6$; the two $L^6_t$
factors form an $L^3_t$ product, and $L^3(I)\hookrightarrow L^{6/5}(I)$ costs
$|I|^{1/2}$, exactly as in \cref{lem:cubic-65}.
\end{proof}

\begin{theorem}[Local existence, uniqueness, and Lipschitz dependence]
\label{thm:LWP}
\uses{def:fixed-point-space, thm:inhom-str, lem:retarded-continuous, lem:cubic-65, lem:cubic-lip}
For every $M>0$ there is $T_M>0$ such that whenever $\norm{u_0}_2\le M$, the
cubic equation $iu_t+u_{xx}=\sigma|u|^2u$ has a unique mild solution
$u=\Phi_{u_0}(u)$ on $I=[t_0-T_M,t_0+T_M]$. It belongs to $\mathcal X(I)$, and
the data-to-solution map is locally Lipschitz from the $L^2$ ball of radius
$M$ into $\mathcal X(I)$. One may take $T_M=c(1+|\sigma|M^2)^{-2}$ with a
sufficiently small absolute constant $c$. Uniqueness holds unconditionally in
the whole class of mild solutions in $\mathcal X(I)$, not only among fixed
points in a ball.
\end{theorem}
\begin{proof}
\uses{thm:hom-strichartz, thm:inhom-str, lem:cubic-65, lem:cubic-lip, lem:retarded-continuous}
Strategy: run a contraction-mapping argument for $\Phi_{u_0}$ on a closed ball
of $\mathcal X(I)$, using the Strichartz bounds to control the linear term and
the cubic estimates for the nonlinear term; smallness comes from $|I|^{1/2}$,
not from the mass.

By \cref{thm:hom-strichartz,thm:inhom-str},
$\norm{\Phi_{u_0}(u)}_{\mathcal X(I)}\le C\norm{u_0}_2+C|\sigma|\norm{|u|^2u}_{L^{6/5}_{t,x}}$.
Fix $R=2CM$ and the closed ball $B_R\subset\mathcal X(I)$. By
\cref{lem:cubic-65}, $\norm{\Phi_{u_0}(u)}_{\mathcal X(I)}\le CM+C|\sigma||I|^{1/2}R^3$;
choosing $|I|$ so that $C|\sigma||I|^{1/2}R^2\le\tfrac14$ keeps $\Phi_{u_0}$
mapping $B_R$ into itself, and \cref{lem:cubic-lip} makes it a strict
contraction. Banach's fixed-point theorem yields a unique fixed point in
$B_R$; the two-data difference estimate gives local Lipschitz dependence, and
\cref{lem:retarded-continuous} places the fixed point in $C(I;L^2)$.
For unconditional uniqueness, given two solutions $u,v\in\mathcal X(I)$ equal
at some $s\in I$, choose a subinterval $J\ni s$ short enough that
$C|\sigma||J|^{1/2}(\norm u_{\mathcal X(I)}^2+\norm v_{\mathcal X(I)}^2)<\tfrac12$;
the base-time identity of \cref{lem:retarded-continuous} rewrites both mild
formulas with base time $s$, and \cref{lem:cubic-lip} gives
$\norm{u-v}_{\mathcal X(J)}\le\tfrac12\norm{u-v}_{\mathcal X(J)}$, so $u=v$ on
$J$; chaining finitely many overlapping short intervals covers $I$.
% FULL PROOF: paper.tex L1536-L1576
\end{proof}

\begin{lemma}[Derivative estimate for the cubic map]
\label{lem:cubic-derivative-map}
\uses{def:fixed-point-space}
Let $J$ be bounded and $u,w\in\mathcal X(J)$. Define the real-linear expression
\[
 DN(u)[w]=2\Rea(\overline u\,w)\,u+|u|^2w.
\]
Then, with an absolute constant $C$,
\begin{equation}
 \norm{DN(u)[w]}_{L^{6/5}_{t,x}(J)}
 \le C|J|^{1/2}\,\norm u_{\mathcal X(J)}^2\,\norm w_{\mathcal X(J)}.
                                                    \label{eq:DN-bound}
\end{equation}
Moreover, for $u,\widetilde u,w,\widetilde w\in\mathcal X(J)$,
\begin{align}
 \norm{DN(u)[w]-DN(\widetilde u)[\widetilde w]}_{L^{6/5}_{t,x}(J)}
 \le C|J|^{1/2}\Big[&
 (\norm u_{\mathcal X}^2+\norm{\widetilde u}_{\mathcal X}^2)
 \norm{w-\widetilde w}_{\mathcal X}\notag\\
 &+(\norm u_{\mathcal X}+\norm{\widetilde u}_{\mathcal X})
 (\norm w_{\mathcal X}+\norm{\widetilde w}_{\mathcal X})
 \norm{u-\widetilde u}_{\mathcal X}\Big].          \label{eq:DN-Lipschitz}
\end{align}
\end{lemma}

\begin{proof}
\uses{lem:cubic-65}
Pointwise $|DN(u)[w]|\le 3|u|^2|w|$. At each time assign one $u$ factor to
$L^2_x$ and the other $u$ factor and $w$ to $L^6_x$ (spatial exponents summing
to $5/6$); the two $L^6_t$ factors give an $L^3_t$ product and
$L^3(J)\hookrightarrow L^{6/5}(J)$ costs $|J|^{1/2}$, proving
\eqref{eq:DN-bound}. Expanding the difference and using
$|DN(u)[w]-DN(\widetilde u)[\widetilde w]|
 \le C(|u|^2+|\widetilde u|^2)|w-\widetilde w|
   +C(|u|+|\widetilde u|)(|w|+|\widetilde w|)|u-\widetilde u|$
gives \eqref{eq:DN-Lipschitz} by the same allocation.
% FULL PROOF: paper.tex L1608-L1622
\end{proof}

\begin{lemma}[Local $H^1$ persistence]\label{lem:H1-persistence}
\uses{lem:cubic-derivative-map, lem:Wp-difference-quotient, def:fixed-point-space}
Let $J$ be a local fixed-point interval centered at $t_0$ on which
$u\in\mathcal X(J)$ satisfies
\begin{equation}
 C|\sigma||J|^{1/2}\norm u_{\mathcal X(J)}^2\le\tfrac14. \label{eq:H1-smallness}
\end{equation}
If $u(t_0)\in H^1(\R)$, then
$u\in C(J;H^1)\cap L^6(J;W^{1,6})$, and the equation
$u_t=iu_{xx}-i\sigma|u|^2u$ holds in $C(J;H^{-1})$. The size of $J$ required
depends only on an upper bound for the $L^2$ mass, not on the $H^1$ norm.
\end{lemma}

\begin{proof}
\uses{thm:inhom-str, lem:cubic-derivative-map, lem:Wp-difference-quotient, lem:H1-C0}
Strategy: solve the linearized (derivative) equation by contraction, then
identify its solution with the distributional derivative $u_x$ through spatial
translation difference quotients, avoiding any assumption that smooth data
generate smooth solutions.

For $z\in\mathcal X(J)$ set
\begin{equation}
 \Psi_u(z)(t)=S(t-t_0)u(t_0)'
 -i\sigma\int_{t_0}^tS(t-s)DN(u(s))[z(s)]\,\dd s.  \label{eq:derivative-fixed-point}
\end{equation}
By \cref{thm:inhom-str,lem:cubic-derivative-map} and \eqref{eq:H1-smallness},
$\Psi_u$ is a strict contraction on $\mathcal X(J)$, with unique fixed point
$z$ obeying $\norm z_{\mathcal X(J)}\le 2C\norm{u(t_0)'}_2$. Writing
$\tau_hf(x)=f(x+h)$, $q_h=(\tau_hu-u)/h$, and $q_{0,h}=(\tau_hu(t_0)-u(t_0))/h$,
translation invariance and uniqueness (\cref{thm:LWP}) give that $\tau_hu$ is
the solution with datum $\tau_hu(t_0)$, and strong $L^2$-continuity of
translations yields $\tau_hu\to u$ in $\mathcal X(J)$. The fundamental theorem
of calculus for $N(a)=|a|^2a$ gives
$(N(\tau_hu)-N(u))/h=A_h[q_h]$ with
$A_h[w]=\int_0^1DN(u+\theta(\tau_hu-u))[w]\,\dd\theta$, so $q_h$ solves the
same equation \eqref{eq:derivative-fixed-point} with datum $q_{0,h}$ and
operator $A_h$. The bounds \eqref{eq:DN-bound}--\eqref{eq:DN-Lipschitz} make
this a contraction with constant $\le\tfrac14$ uniformly in $h$, and since
$q_{0,h}\to u(t_0)'$ in $L^2$, subtracting \eqref{eq:derivative-fixed-point}
gives $q_h\to z$ in $\mathcal X(J)$. By
\cref{lem:Wp-difference-quotient}, translation difference quotients converge
to $u_x$ distributionally, hence $z=u_x$ and
$u\in C(J;H^1)\cap L^6(J;W^{1,6})$. Finally
$u_t=iu_{xx}-i\sigma|u|^2u$ with $u_{xx}\in C(J;H^{-1})$ and, by the
one-dimensional embedding $H^1\hookrightarrow L^\infty$ of \cref{lem:H1-C0},
$|u|^2u\in C(J;L^2)\subset C(J;H^{-1})$.
% FULL PROOF: paper.tex L1639-L1701
\end{proof}

\begin{proposition}[Mass conservation in $L^2$]\label{prop:mass-L2}
\uses{thm:LWP}
Every $L^2$ mild solution of $iu_t+u_{xx}=\sigma|u|^2u$ satisfies
\begin{equation}
  \norm{u(t)}_2=\norm{u(t_0)}_2\qquad(t,t_0\text{ in its lifespan}).
  \label{eq:mass-conservation}
\end{equation}
\end{proposition}

\begin{proof}
\uses{thm:LWP, lem:H1-persistence, lem:mass-H1}
Fix a mass-controlled local interval centered at $t_0$ and choose
$u_{0,n}\in H^1$ with $u_{0,n}\to u(t_0)$ in $L^2$ and uniformly bounded $L^2$
norms. The local construction of \cref{thm:LWP} produces all approximating
solutions on the same interval; by \cref{lem:H1-persistence} they lie in
$C(J;H^1)$, and by \cref{lem:mass-H1} their masses are constant. Local flow
stability gives $u_n\to u$ in $C_tL^2_x$, so \eqref{eq:mass-conservation}
passes to the limit on that interval. Restarting at any point of the lifespan
and chaining overlapping intervals proves the identity throughout.
\end{proof}

\begin{lemma}[Gelfand-triple mass chain rule]\label{lem:gelfand-chain-rule}
\uses{def:Hs-space}
Let $V\hookrightarrow H\hookrightarrow V^*$ be a dense continuous Gelfand
triple, with $H$ identified with its anti-dual. If $w\in C([a,b];V)$ with
distributional time derivative $w_t\in C([a,b];V^*)$, then
$t\mapsto\norm{w(t)}_H^2$ is continuously differentiable and
\begin{equation}
 \frac{\dd}{\dd t}\norm{w(t)}_H^2
 =2\Rea\pair{w_t(t)}{w(t)}_{V^*,V}.                \label{eq:gelfand-chain-rule}
\end{equation}
The same conclusion, with absolute continuity in place of continuous
differentiability, holds when $w\in L^p(a,b;V)$ and $w_t\in L^{p'}(a,b;V^*)$
for dual exponents $p,p'$, $w$ taken in its standard continuous $H$
representative.
\end{lemma}

\begin{proof}
\uses{def:Hs-space}
Strategy: mollify in time and pass to the limit in the classical
Hilbert-space product rule. Extend $w$ constantly beyond each endpoint and
$w_t$ by zero; continuity of the extension makes its distributional derivative
exactly the zero extension of $w_t$, with no endpoint Dirac mass. Convolving
in time with a smooth approximate identity, $w_\eps$ is smooth into $V$ with
classical derivative agreeing in $V^*$ with $(w_t)_\eps$, and the product rule
gives
$\norm{w_\eps(t)}_H^2-\norm{w_\eps(s)}_H^2
 =2\Rea\int_s^t\pair{(w_t)_\eps(\tau)}{w_\eps(\tau)}_{V^*,V}\,\dd\tau$.
Uniform convergence $w_\eps\to w$ in $V$ and $(w_t)_\eps\to w_t$ in $V^*$ on
compact subintervals yields the integrated form of
\eqref{eq:gelfand-chain-rule}, and continuity of its right side gives the
pointwise derivative. In the integrable case one works on compact
subintervals with the corresponding Bochner-space convergence and H\"older;
the integrated identity supplies the continuous $H$ representative and
absolute continuity.
% FULL PROOF: paper.tex L1722-L1743
\end{proof}

\begin{lemma}[$H^1$ mass identity]\label{lem:mass-H1}
\uses{lem:gelfand-chain-rule, lem:H1-persistence}
If $u\in C(J;H^1)$ solves $iu_t+u_{xx}=\sigma|u|^2u$ on $J$, then
\[
  \norm{u(t)}_2=\norm{u(t_0)}_2\qquad(t,t_0\in J).
\]
\end{lemma}

\begin{proof}
\uses{lem:gelfand-chain-rule}
Apply \cref{lem:gelfand-chain-rule} to the triple
$H^1\hookrightarrow L^2\hookrightarrow H^{-1}$. Since
$u_t=iu_{xx}-i\sigma|u|^2u$ in $H^{-1}$,
\[
 \frac{\dd}{\dd t}\norm u_2^2
 =2\Rea\pair{iu_{xx}-i\sigma|u|^2u}{u}_{H^{-1},H^1}
 =2\Rea\bigl(-i\norm{u_x}_2^2-i\sigma\norm u_4^4\bigr)=0,
\]
the integration by parts in the first pairing being the defining
weak-derivative identity and both parenthesized quantities being real.
\end{proof}

\begin{theorem}[Global $L^2$ well-posedness and gauge covariance]\label{thm:GWP}
\uses{thm:LWP, prop:mass-L2, def:fixed-point-space}
For every $u_0\in L^2(\R)$ and every $\sigma\in\R$ there is a unique global
mild solution
\[
 u\in C(\R;L^2)\cap L^6_{\loc}(\R^2).
\]
The flow map $\Phi_t:u_0\mapsto u(t)$ is continuous, obeys the group law, and
is gauge covariant:
\begin{equation}
  \Phi_t(\zeta u_0)=\zeta\,\Phi_t(u_0)\qquad(|\zeta|=1). \label{eq:gauge-cov}
\end{equation}
\end{theorem}

\begin{proof}
\uses{thm:LWP, prop:mass-L2}
The local lifespan in \cref{thm:LWP} depends only on an upper bound for the
$L^2$ norm, and by \cref{prop:mass-L2} that norm stays equal to
$\norm{u_0}_2$; hence the local construction iterates indefinitely in both
time directions. Uniqueness makes overlapping pieces agree and yields the
group law, and iterating the local Lipschitz estimate on any fixed compact
time interval gives continuity of the flow. Finally, if $u$ solves the
equation then $\zeta u$ solves it with datum $\zeta u_0$ for $|\zeta|=1$, so
uniqueness gives \eqref{eq:gauge-cov}.
\end{proof}


\chapter{Persistence and endpoint regularity}\label{chap:endpoint}

\section{Higher regularity spaces and persistence of derivatives}

\begin{definition}[Higher regularity space $\mathcal X_k(J)$]
\label{def:Xk-space}
\uses{def:fixed-point-space, def:Hs-space}
Let $k\ge0$ be an integer and let $J\subset\R$ be a bounded interval.  Recall
the fixed-point space $\mathcal X(J)=C(J;L^2(\R))\cap L^6(J\times\R)$.  Define
\[
 \mathcal X_k(J)=C(J;H^k(\R))\cap L^6(J;W^{k,6}(\R)),
 \qquad
 \norm u_{\mathcal X_k(J)}
 =\sum_{j=0}^k\norm{\partial_x^ju}_{\mathcal X(J)}.
\]
Since each weak derivative $\partial_x^j$ is a closed operator and
$\mathcal X(J)$ is a Banach space, $\bigl(\mathcal X_k(J),\norm\cdot_{\mathcal X_k(J)}\bigr)$
is a Banach space.
\end{definition}

\begin{lemma}[Cubic Leibniz decomposition and lower-order bound]
\label{lem:cubic-Hk-leibniz}
\uses{def:Xk-space, def:scalar-mixed-norm, lem:cubic-derivative-map}
Let $k\ge1$.  For a smooth complex-valued $u$, with
$DN(u)[w]=2\Rea(\overline u\,w)u+|u|^2w$,
\begin{equation}
 \partial_x^k(|u|^2u)=DN(u)[\partial_x^ku]+R_k(u),
 \label{eq:top-derivative-split}
\end{equation}
where
\begin{equation}
 R_k(u)=
 \sum_{\substack{a+b+c=k\\0\le a,b,c\le k-1}}
 \frac{k!}{a!\,b!\,c!}
 (\partial_x^au)(\partial_x^b\overline u)(\partial_x^cu).
 \label{eq:Rk-explicit}
\end{equation}
If $u\in\mathcal X_{k-1}(J)$, then $R_k(u)$ is well-defined in
$L^{6/5}(J\times\R)$ and
\begin{equation}
 \norm{R_k(u)}_{L^{6/5}_{t,x}(J)}
 \le C_k\,|J|^{1/2}\,\norm u_{\mathcal X_{k-1}(J)}^3.
 \label{eq:Rk-bound}
\end{equation}
Moreover $R_k:\mathcal X_{k-1}(J)\to L^{6/5}(J\times\R)$ is locally Lipschitz.
\end{lemma}

\begin{proof}
\uses{lem:cubic-derivative-map}
The multinomial Leibniz rule applied to $\partial_x^k(u\overline u u)$ gives
\[
 \partial_x^k(u\overline u u)
 =\sum_{a+b+c=k}\frac{k!}{a!\,b!\,c!}
   (\partial_x^au)(\partial_x^b\overline u)(\partial_x^cu).
\]
The three terms in which one of the indices equals $k$ are
$|u|^2\partial_x^ku+u^2\partial_x^k\overline u+|u|^2\partial_x^ku
=DN(u)[\partial_x^ku]$, and every remaining monomial has all indices at most
$k-1$; this proves \eqref{eq:top-derivative-split}--\eqref{eq:Rk-explicit}.
For a single remaining monomial $f_1f_2f_3$, assign one factor to
$L^\infty_tL^2_x$ and the other two to $L^6_tL^6_x$.  The spatial reciprocal
exponents add to $1/2+1/6+1/6=5/6$, the two $L^6$ factors form an $L^3_t$
function, and $L^3(J)\hookrightarrow L^{6/5}(J)$ has norm $|J|^{1/2}$, whence
\[
 \norm{f_1f_2f_3}_{L^{6/5}_{t,x}}
 \le |J|^{1/2}\norm{f_1}_{L^\infty_tL^2_x}
              \norm{f_2}_{L^6_tL^6_x}
              \norm{f_3}_{L^6_tL^6_x}.
\]
Summing the finitely many monomials gives \eqref{eq:Rk-bound}.  Expanding a
difference of two cubic monomials into three terms, each with one factor
replaced by its difference, gives local Lipschitz continuity by the same
allocation.
\end{proof}

\begin{proposition}[Persistence of all integer derivatives]
\label{prop:Hk-persistence}
\uses{def:Xk-space, lem:cubic-Hk-leibniz, thm:GWP}
Let $k\ge1$ be an integer.  If $u_0\in H^k(\R)$ and $u$ is the global $L^2$
solution with datum $u_0$, then on every compact time interval $J$
\[
 u\in\mathcal X_k(J)=C(J;H^k(\R))\cap L^6(J;W^{k,6}(\R)).
\]
\end{proposition}

\begin{proof}
\uses{lem:H1-persistence, thm:LWP, thm:inhom-str, lem:cubic-Hk-leibniz, lem:cubic-derivative-map, lem:Wp-difference-quotient, thm:GWP}
Strategy: induct on $k$, upgrading the local $H^1$ theory one derivative at a
time on mass-controlled fixed-point intervals and then covering every compact
interval by finitely many of them.

The base case $k=1$ is \cref{lem:H1-persistence}; by \cref{thm:LWP} and
\cref{thm:GWP} the fixed-point interval length and the bound for
$\norm u_{\mathcal X(J)}$ depend only on the conserved $L^2$ norm, so finitely
many such intervals cover any compact interval.  Assume persistence through
order $k-1$ with $k\ge2$, and work on a mass-controlled interval $J$ centered
at $t_0$ short enough that
\[
 C|\sigma|\,|J|^{1/2}\norm u_{\mathcal X(J)}^2\le\tfrac14.
\]
By the induction hypothesis and \cref{lem:cubic-Hk-leibniz},
$R_k(u)\in L^{6/5}_{t,x}(J)$, and the affine map
\[
 \Theta_k(z)(t)=S(t-t_0)\partial_x^ku(t_0)
 -i\sigma\int_{t_0}^tS(t-s)\bigl(DN(u(s))[z(s)]+R_k(u)(s)\bigr)\,\dd s
\]
is a contraction on $\mathcal X(J)$ by \cref{thm:inhom-str},
\cref{lem:cubic-derivative-map}, and the smallness of $|J|$; let $z_k$ be its
unique fixed point.  One identifies $z_k$ with $\partial_x^ku$ by a
translation difference-quotient argument: with $y=\partial_x^{k-1}u$ and
$q_h=(\tau_hy-y)/h$, differentiating the mild equation $k-1$ times,
translating, subtracting, and dividing by $h$ yields the same fixed-point
equation with coefficients $A_h[w]=DN(\tau_hu)[w]$ and an explicit remainder
$R_{k,h}$ converging to $R_k(u)$ in $L^{6/5}_{t,x}(J)$; here
\cref{lem:Wp-difference-quotient} supplies the $L^6_tL^6_x$ convergence of
lower-order difference quotients and the uniformly bounded multipliers
$(\e^{ih\xi}-1)/(ih\xi)\to1$ supply the $C_tL^2_x$ convergence.  The
contraction estimate then gives $q_h\to z_k$ in $\mathcal X(J)$, while
$q_h\to\partial_x^ku$ in $\mathcal{D}'(J\times\R)$, so $z_k=\partial_x^ku$ and
$u\in\mathcal X_k(J)$.  Finite iteration over the mass-controlled intervals
proves the claim on every compact interval.
% FULL PROOF: paper.tex L1891-L1994
\end{proof}

\begin{lemma}[Schwartz data yield a smooth solution]
\label{lem:schwartz-smooth}
\uses{prop:Hk-persistence}
If $u_0\in\Ssch(\R)$ and $u$ is the corresponding global $L^2$ solution, then
$u$ is smooth in $(t,x)$ on every compact time interval.
\end{lemma}

\begin{proof}
\uses{prop:Hk-persistence}
Since $u_0\in H^k(\R)$ for every $k$, \cref{prop:Hk-persistence} gives all
spatial derivatives locally in time.  The equation
$u_t=iu_{xx}-i\sigma|u|^2u$ then yields the first time derivative with
arbitrarily high spatial regularity, and differentiating the equation
repeatedly in time yields every mixed derivative; hence $u\in C^\infty$.
\end{proof}

\section{Endpoint spacetime regularity}

\begin{lemma}[Interpolation to $L^8_tL^4_x$]
\label{lem:interp-L8L4}
\uses{def:scalar-mixed-norm}
Let $I\subset\R$ be bounded and let $u$ be measurable on $I\times\R$.  Then
\[
 \norm u_{L^8_tL^4_x(I)}
 \le\norm u_{L^\infty_tL^2_x(I)}^{1/2}
      \norm u_{L^4_tL^\infty_x(I)}^{1/2}.
\]
\end{lemma}

\begin{proof}
\uses{}
Pointwise in time, $\norm{u(t)}_4^4\le\norm{u(t)}_2^2\norm{u(t)}_\infty^2$.
Squaring and integrating in time gives
\[
 \int_I\norm{u(t)}_4^8\,\dd t
 \le\norm u_{L^\infty_tL^2_x(I)}^4\int_I\norm{u(t)}_\infty^4\,\dd t
 =\norm u_{L^\infty_tL^2_x(I)}^4\,\norm u_{L^4_tL^\infty_x(I)}^4,
\]
and taking eighth roots gives the claim.
\end{proof}

\begin{proposition}[Endpoint bound and integrable forcing]
\label{prop:endpoint}
\uses{thm:hom-strichartz, thm:GWP, lem:scalar-mixed-minkowski, lem:interp-L8L4}
Let $I\subset\R$ be bounded and let $u$ be a global $L^2$ solution.  Then
\begin{align}
 u&\in L^4(I;L^\infty(\R)), \label{eq:L4Linfty}\\
 u&\in L^8(I;L^4(\R)),      \label{eq:L8L4}\\
 |u|^2u&\in L^1(I;L^2(\R)). \label{eq:forcing-L1L2}
\end{align}
Moreover, if $u_n$ are the solutions with data $u_{0,n}\to u_0$ in $L^2$, then
on every bounded $I$
\[
 u_n\to u\ \text{ in }L^4_tL^\infty_x,
 \qquad
 |u_n|^2u_n\to|u|^2u\ \text{ in }L^1_tL^2_x.
\]
\end{proposition}

\begin{proof}
\uses{thm:hom-strichartz, thm:GWP, lem:scalar-mixed-minkowski, lem:interp-L8L4, lem:retarded-continuous}
Strategy: bound the cubic forcing in $L^1_tL^2_x$, run the homogeneous
endpoint Strichartz estimate through Duhamel to reach $L^4_tL^\infty_x$, and
interpolate to $L^8_tL^4_x$; convergence follows from a cubic difference
estimate.

For \eqref{eq:forcing-L1L2},
\[
 \norm{|u|^2u}_{L^1_tL^2_x(I)}=\int_I\norm{u(t)}_6^3\,\dd t
 \le|I|^{1/2}\norm u_{L^6(I\times\R)}^3<\infty,
\]
the last norm being finite by \cref{thm:GWP}.  For \eqref{eq:L4Linfty}, use the
Duhamel formula and the homogeneous endpoint estimate $(q,r)=(4,\infty)$ of
\cref{thm:hom-strichartz}: restricting the time variable can only decrease the
scalar $L^4_tL^\infty_x$ norm, so by \cref{lem:scalar-mixed-minkowski}
\[
 \Bigl\|\int_{t_0}^tS(t-s)F(s)\,\dd s\Bigr\|_{L^4_tL^\infty_x(I)}
 \le\int_I\norm{\1_{\{s\text{ between }t_0,t\}}S(t-s)F(s)}_{L^4_tL^\infty_x(I)}\,\dd s
 \le C\int_I\norm{F(s)}_2\,\dd s,
\]
and with $F=|u|^2u$ the bound \eqref{eq:forcing-L1L2} gives
\eqref{eq:L4Linfty}.  Then \eqref{eq:L8L4} follows from \cref{lem:interp-L8L4}
together with the $L^\infty_tL^2_x$ and $L^4_tL^\infty_x$ bounds.  For
convergence, $\norm{|u_n|^2u_n-|u|^2u}_{L^2_x}
\le C(\norm{u_n}_6^2+\norm u_6^2)\norm{u_n-u}_6$; the first factor is bounded in
$L^3_t$, the last converges in $L^6_t$, and Hölder on the finite interval gives
$L^1_tL^2_x$ convergence, which the endpoint Duhamel estimate turns into
$L^4_tL^\infty_x$ convergence.  Since the forcing lies in both $L^{6/5}_{t,x}$
and $L^1_tL^2_x$, \cref{lem:retarded-continuous} shows every mild formula is an
ordinary $L^2$-valued Bochner integral.
% FULL PROOF: paper.tex L2024-L2063
\end{proof}

\begin{corollary}[Continuous spatial representatives at almost every time]
\label{cor:C0-slices}
\uses{prop:endpoint, lem:H1-C0, lem:C0-limit-identification}
For every bounded interval $I\subset\R$ there is a full-measure set
$E_I\subset I$ such that, for every $t\in E_I$, the $L^2$ class $u(t)$ has a
representative in $C_0(\R)$.
\end{corollary}

\begin{proof}
\uses{prop:Hk-persistence, prop:endpoint, lem:mixed-subsequence-uniform, lem:H1-C0, lem:C0-limit-identification}
Choose Schwartz data $u_{0,n}\to u_0$ in $L^2$ and let $u_n$ be the
corresponding solutions.  By \cref{prop:Hk-persistence}, $u_n(t)\in H^1(\R)$
for every $t$, so by \cref{lem:H1-C0} every $u_n(t,\cdot)$ has a representative
in $C_0(\R)$.  By \cref{prop:endpoint}, $u_n\to u$ in $L^4(I;L^\infty_x)$;
applying \cref{lem:mixed-subsequence-uniform}, for almost every $t$ a
subsequence is uniformly Cauchy in $x$ and hence converges in $C_0(\R)$.  At
every such time the same sequence converges to $u(t)$ in $L^2$, since the flow
convergence holds in $C_tL^2_x$, so \cref{lem:C0-limit-identification}
identifies the uniform limit with a representative of $u(t)$.  A countable
exhaustion of $I$ and intersection of the resulting full-measure sets yield a
single set $E_I$.
\end{proof}


\chapter{Local smoothing by one half derivative}\label{chap:smoothing}

The cubic flow gains one half spatial derivative locally in space and time.
The gain is read off the Fourier representation of the free propagator and
then transported to forced solutions by Duhamel.

\section{The homogeneous identity}

\begin{lemma}[Pointwise-in-space homogeneous smoothing identity]\label{lem:homogeneous-smoothing-identity}
\uses{def:propagator, def:homogeneous-multiplier, def:Hs-space}
Let $f\in\Ssch(\R)$ and $x\in\R$. Then
\[
 |D|^{1/2}S(t)f(x)
 =(2\pi)^{-1/2}\int_\R
    \e^{ix\xi-it\xi^2}|\xi|^{1/2}\widehat f(\xi)\,\dd\xi ,
\]
and, up to one fixed normalization constant, the map $t\mapsto|D|^{1/2}S(t)f(x)$
is the inverse time Fourier transform of
\begin{equation}
 B_xf(\lambda)=\1_{\{\lambda>0\}}\lambda^{-1/4}
 \left(\e^{ix\sqrt\lambda}\widehat f(\sqrt\lambda)
      +\e^{-ix\sqrt\lambda}\widehat f(-\sqrt\lambda)\right).
 \label{eq:smoothing-Bx}
\end{equation}
Moreover $B_xf\in L^2_\lambda(\R)$ with
\begin{equation}
 \norm{B_xf}_{L^2_\lambda}^2
 \le C\int_0^\infty \lambda^{-1/2}
 \left(|\widehat f(\sqrt\lambda)|^2
      +|\widehat f(-\sqrt\lambda)|^2\right)\,\dd\lambda
 =C\norm f_2^2,
 \label{eq:smoothing-Bx-bound}
\end{equation}
the constant $C$ being independent of $x$.
\end{lemma}

\begin{proof}
\uses{def:propagator, def:homogeneous-multiplier}
By the definition of $S(t)$ and of the homogeneous multiplier $|D|^{1/2}$,
Fourier inversion gives the displayed representation of $|D|^{1/2}S(t)f(x)$.
Split the integral into the positive and negative frequency half-lines and
substitute $\lambda=\xi^2$, so that $\xi=\pm\sqrt\lambda$ and
$\dd\xi=\tfrac12\lambda^{-1/2}\dd\lambda$; the Jacobian factor together with
$|\xi|^{1/2}=\lambda^{1/4}$ produces the weight $\lambda^{-1/4}$ and yields the
integrand $\e^{-it\lambda}B_xf(\lambda)$ up to one fixed constant, which is
\eqref{eq:smoothing-Bx}. Plancherel in $\lambda$ and the elementary inequality
$(a+b)^2\le 2a^2+2b^2$ give the first inequality of
\eqref{eq:smoothing-Bx-bound}; reversing the substitution $\lambda=\xi^2$ turns
the last integral into $C\norm f_2^2$. No step depends on $x$, since $x$ enters
only through the unimodular phases $\e^{\pm ix\sqrt\lambda}$.
\end{proof}

\begin{definition}[Canonical representative $A_xf$ of $|D|^{1/2}S(t)f(x)$]\label{def:Ax-representative}
\uses{def:homogeneous-multiplier, def:propagator, lem:homogeneous-smoothing-identity}
For $f\in L^2(\R)$ and $x\in\R$, the function $B_xf$ of \eqref{eq:smoothing-Bx}
still belongs to $L^2_\lambda(\R)$. Its inverse time Fourier transform is a
well-defined element
\[
 A_xf\in L^2_t(\R),
\]
called the \emph{canonical representative} of $t\mapsto|D|^{1/2}S(t)f(x)$. The
assignment $x\mapsto B_xf$ is continuous from $\R$ into $L^2_\lambda$, because the
phase factors $\e^{\pm ix\sqrt\lambda}$ converge pointwise as $x_n\to x$ and are
dominated by the $x$-independent integrable majorant of
\eqref{eq:smoothing-Bx-bound}; hence $x\mapsto A_xf$ is continuous from $\R$ into
$L^2_t$ and admits a jointly measurable representative $(t,x)\mapsto A_xf(t)$.
\end{definition}

\begin{lemma}[Pointwise-in-space smoothing bound]\label{lem:point-smoothing}
\uses{def:Ax-representative, lem:homogeneous-smoothing-identity, def:propagator, def:homogeneous-multiplier}
For every $f\in L^2(\R)$ and every $x\in\R$, the canonical representative
$A_xf$ of \cref{def:Ax-representative} agrees with $|D|^{1/2}S(t)f(x)$ and
satisfies
\begin{equation}
 \int_\R\big| |D|^{1/2}S(t)f(x)\big|^2\,\dd t
 =\norm{A_xf}_{L^2_t}^2
 \le C\norm f_2^2 ,
 \label{eq:point-smoothing}
\end{equation}
with $C$ independent of $x$.
\end{lemma}

\begin{proof}
\uses{lem:homogeneous-smoothing-identity, def:Ax-representative}
For $f\in\Ssch$ the identity \eqref{eq:smoothing-Bx} and the inverse Plancherel
theorem in $t$ give $A_xf=|D|^{1/2}S(t)f(x)$ and, by
\eqref{eq:smoothing-Bx-bound}, $\norm{A_xf}_{L^2_t}^2=\norm{B_xf}_{L^2_\lambda}^2
\le C\norm f_2^2$, uniformly in $x$; this is \eqref{eq:point-smoothing} on the
dense class $\Ssch$. For general $f\in L^2$ the same estimate holds for the
canonical representative $A_xf$ of \cref{def:Ax-representative}. It remains to
identify $A_xf$ with the Fourier-multiplier distribution $|D|^{1/2}S(t)f$. Pick
Schwartz $f_n\to f$ in $L^2$. On every compact spatial set $K$,
\[
 \int_K\norm{A_x(f_n-f)}_{L^2_t}^2\,\dd x
 \le C|K|\,\norm{f_n-f}_2^2\longrightarrow 0 ,
\]
so $A_xf_n\to A_xf$ in $L^2_{t,x}(K\times\R)$. Since $A_xf_n=|D|^{1/2}S(t)f_n(x)$
converges to $|D|^{1/2}S(t)f$ in distributions as well, the two limits coincide,
and the norm bound \eqref{eq:point-smoothing} follows.
% FULL PROOF: paper.tex L2109-L2156
\end{proof}

\begin{corollary}[Localized homogeneous smoothing]\label{cor:hom-smoothing}
\uses{lem:point-smoothing, def:propagator, def:Hs-space, def:scalar-mixed-norm}
For every $\chi\in C_c^\infty(\R)$ and every bounded interval $I$ there is a
constant $C_{\chi,I}$ such that, for all $f\in L^2(\R)$,
\begin{equation}
 \norm{\chi S(t)f}_{L^2(I;H^{1/2}_x)}
 \le C_{\chi,I}\norm f_2 .
 \label{eq:hom-local-smoothing}
\end{equation}
\end{corollary}

\begin{proof}
\uses{lem:point-smoothing}
Integrating \eqref{eq:point-smoothing} against $|\chi(x)|^2$ over $x$ gives
\[
 \norm{\chi\,|D|^{1/2}S(t)f}_{L^2_{t,x}(\R^2)}
 \le C\,\norm\chi_2\,\norm f_2 .
\]
To pass from $\chi|D|^{1/2}$ to $|D|^{1/2}\chi$, note that in Fourier variables
the commutator $[|D|^{1/2},\chi]$ has kernel
$(|\xi|^{1/2}-|\eta|^{1/2})\widehat\chi(\xi-\eta)$. Since
$\big||\xi|^{1/2}-|\eta|^{1/2}\big|\le|\xi-\eta|^{1/2}$ and
$|\theta|^{1/2}\widehat\chi(\theta)\in L^1$, Young's inequality yields
$\norm{[|D|^{1/2},\chi]g}_2\le C_\chi\norm g_2$. Finally the $L^2_x$ part of the
$H^{1/2}$ norm contributes $|I|^{1/2}\norm\chi_\infty\norm f_2$ by unitarity of
$S(t)$. Adding the two contributions gives \eqref{eq:hom-local-smoothing}.
\end{proof}

\section{Duhamel smoothing}

\begin{lemma}[Local smoothing for forced solutions]\label{lem:Duhamel-smoothing}
\uses{cor:hom-smoothing, thm:inhom-str, def:propagator, def:Hs-space, def:scalar-mixed-norm}
Let $I$ be a bounded interval, $t_0\in I$, $f\in L^2(\R)$ and
$F\in L^1(I;L^2(\R))$. Then
\[
 w(t)=S(t-t_0)f+\int_{t_0}^tS(t-s)F(s)\,\dd s
\]
satisfies, for every $\chi\in C_c^\infty(\R)$,
\begin{equation}
 \norm{\chi w}_{L^2(I;H^{1/2})}
 \le C_{\chi,I}\left(\norm f_2+\norm F_{L^1_tL^2_x(I)}\right).
 \label{eq:forced-smoothing}
\end{equation}
\end{lemma}

\begin{proof}
\uses{cor:hom-smoothing}
The homogeneous term $S(t-t_0)f$ is controlled by $C_{\chi,I}\norm f_2$ through
\cref{cor:hom-smoothing}. For the Duhamel term, Minkowski's integral inequality
and translation invariance in time give
\begin{align*}
 \Big\|\chi\!\int_{t_0}^tS(t-s)F(s)\,\dd s\Big\|_{L^2_tH^{1/2}_x(I)}
 &\le\int_I
 \norm{\chi\,\1_{\{s\text{ between }t_0\text{ and }t\}}
       S(t-s)F(s)}_{L^2_tH^{1/2}_x(I)}\,\dd s\\
 &\le C_{\chi,I}\int_I\norm{F(s)}_2\,\dd s ,
\end{align*}
where each inner norm is bounded by \cref{cor:hom-smoothing} applied to the
datum $F(s)$ at base time $s$, and restricting the time domain to $I$ only
decreases the norm. Summing the two contributions gives
\eqref{eq:forced-smoothing}.
\end{proof}

\begin{corollary}[Local smoothing for the cubic flow]\label{cor:NLS-smoothing}
\uses{lem:Duhamel-smoothing, prop:endpoint}
Every global $L^2$ solution $u$ of the cubic flow satisfies
\begin{equation}
  \chi u\in L^2(I;H^{1/2}(\R))
  \qquad(\chi\in C_c^\infty(\R),\ I\Subset\R).
  \label{eq:NLS-local-smoothing}
\end{equation}
If $u_n$ are smooth-data solutions converging to $u$ in the $L^2$ flow, then for
every such $\chi$ and $I$,
\begin{equation}
  \chi u_n\to\chi u
  \quad\text{in }L^2(I;H^{1/2}).
  \label{eq:smoothing-convergence}
\end{equation}
\end{corollary}

\begin{proof}
\uses{lem:Duhamel-smoothing, prop:endpoint}
Write $u$ through its mild (Duhamel) equation with initial datum $u(t_0)\in L^2$
and forcing $F=|u|^2u$. By \cref{prop:endpoint}, in particular
\eqref{eq:forcing-L1L2}, one has $F\in L^1(I;L^2)$, so
\cref{lem:Duhamel-smoothing} yields \eqref{eq:NLS-local-smoothing}. For the
differences $u_n-u$, the initial data converge in $L^2$ and, by the convergence
statement of \cref{prop:endpoint}, the cubic forcings $|u_n|^2u_n\to|u|^2u$ in
$L^1_tL^2_x$; applying \cref{lem:Duhamel-smoothing} to the difference gives
\eqref{eq:smoothing-convergence}.
\end{proof}

\section{A useful countable full-measure set}

\begin{lemma}[A countable full-measure set of good times]\label{lem:countable-full-measure}
\uses{cor:NLS-smoothing, cor:C0-slices, def:Hs-loc}
Let $u,v$ be two global $L^2$ solutions of the cubic flow. There is a set
$E\subset\R$ of full measure such that, for every $t\in E$,
\begin{equation}
 u(t),\,v(t)\in H^{1/2}_{\loc}(\R)\cap C_0(\R).
 \label{eq:good-slices}
\end{equation}
\end{lemma}

\begin{proof}
\uses{cor:NLS-smoothing, cor:C0-slices}
Fix an increasing sequence of bounded intervals $I_m\Subset\R$ and cutoffs
$\chi_m\in C_c^\infty(\R)$ with $\chi_m\equiv1$ on $[-m,m]$. By
\eqref{eq:NLS-local-smoothing} of \cref{cor:NLS-smoothing},
$\chi_m u,\chi_m v\in L^2(I_m;H^{1/2})$, so for almost every $t$ one has
$\chi_m u(t),\chi_m v(t)\in H^{1/2}$; since $\chi_m\equiv1$ on $[-m,m]$ this
gives $u(t),v(t)\in H^{1/2}_{\loc}$ after intersecting the countably many
full-measure sets over $m$. Intersecting further with the full-measure set of
\cref{cor:C0-slices} on which $u(t),v(t)$ have $C_0(\R)$ representatives yields a
single full-measure set $E$ satisfying \eqref{eq:good-slices}. The set $E$ may be
shrunk later by countably many distributional identities without changing its
measure.
\end{proof}


\chapter{Critical $H^{1/2}$ calculus and the current}\label{chap:critical-calculus}

The endpoint argument repeatedly multiplies an $H^{1/2}$ function by a bounded $H^{1/2}$ function, composes with Lipschitz maps, and pairs a derivative against such products. At the critical index these operations are valid only when the $L^\infty$ information is retained. This chapter records the exact estimates and defines the current $\overline f f_x$ as an $H^{-1/2}$ distribution.

\section{The product estimate}

\begin{lemma}[Critical $H^{1/2}$ product estimate: $H^{1/2}\cap L^\infty$ is an algebra]\label{lem:gagliardo-product}
\uses{def:Hs-space}
If $f,g\in H^{1/2}(\R)\cap L^\infty(\R)$, then $fg\in H^{1/2}(\R)$ and
\begin{equation}
 \norm{fg}_{H^{1/2}}
 \le C\left(\norm{f}_{\infty}\,\norm{g}_{H^{1/2}}
           +\norm{g}_{\infty}\,\norm{f}_{H^{1/2}}\right).
 \label{eq:Hhalf-product}
\end{equation}
The corresponding statement holds locally after multiplication by a smooth cutoff.
\end{lemma}

\begin{proof}
\uses{thm:fourier-gagliardo, lem:smooth-multiplier-quantitative}
The $L^2$ part is controlled by placing either factor in $L^\infty$:
$\norm{fg}_{L^2}\le\min\{\norm{f}_\infty\norm{g}_{L^2},\norm{g}_\infty\norm{f}_{L^2}\}$.
For the Gagliardo seminorm write
\[
 f(x)g(x)-f(y)g(y)
 =f(x)\bigl(g(x)-g(y)\bigr)+g(y)\bigl(f(x)-f(y)\bigr).
\]
Squaring and using $(a+b)^2\le 2a^2+2b^2$ gives
\[
 |f(x)g(x)-f(y)g(y)|^2
 \le 2\norm{f}_\infty^2\,|g(x)-g(y)|^2
   +2\norm{g}_\infty^2\,|f(x)-f(y)|^2 .
\]
Dividing by $|x-y|^2$, integrating over $\R^2$, and using the Fourier identification of the Gagliardo seminorm from \cref{thm:fourier-gagliardo} yields \eqref{eq:Hhalf-product}. The local statement follows by applying the global estimate to $\chi f,\chi g$ for a smooth cutoff $\chi$ and controlling the resulting cutoff factors by \cref{lem:smooth-multiplier-quantitative}.
\end{proof}

\begin{lemma}[Lipschitz composition in $H^{1/2}$]\label{lem:lipschitz-composition}
\uses{def:Hs-space}
Let $F:\C\to\C$ be Lipschitz with $F(0)=0$. If $f\in H^{1/2}(\R)$, then $F\circ f\in H^{1/2}(\R)$ and
\[
 \norm{F(f)}_{H^{1/2}}\le C\,\operatorname{Lip}(F)\,\norm{f}_{H^{1/2}}.
\]
If $F$ is merely Lipschitz on the range of a locally bounded function, the same conclusion holds locally after subtracting the constant $F(0)$.
\end{lemma}

\begin{proof}
\uses{thm:fourier-gagliardo, lem:gagliardo-product}
The pointwise inequalities
\[
 |F(f(x))|\le\operatorname{Lip}(F)\,|f(x)|,
 \qquad
 |F(f(x))-F(f(y))|\le\operatorname{Lip}(F)\,|f(x)-f(y)|,
\]
control the $L^2$ norm and, after dividing by $|x-y|^2$ and integrating, the Gagliardo seminorm identified in \cref{thm:fourier-gagliardo}. The local version follows by first localizing with a smooth cutoff (using the product estimate \cref{lem:gagliardo-product}) and applying the global bound to the truncated function.
\end{proof}

\section{Mollification in critical and uniform norms}

\begin{lemma}[Simultaneous approximation in $H^{1/2}$ and $L^\infty$]\label{lem:mollification-critical-uniform}
\uses{def:Hs-space, def:fourier-convention}
Let $f\in H^{1/2}(\R)\cap C_0(\R)$ and let $\varphi_\delta$ be a standard smooth approximate identity. Then
\[
 \varphi_\delta*f\longrightarrow f
 \qquad\text{in }H^{1/2}(\R)\ \text{and in }L^\infty(\R).
\]
The same is true locally after inserting a cutoff.
\end{lemma}

\begin{proof}
\uses{def:fourier-convention, thm:fourier-gagliardo}
With the unitary Fourier normalization, convolution by $\varphi_\delta$ has multiplier $m_\delta(\xi)=(2\pi)^{1/2}\,\widehat\varphi(\delta\xi)$. It is bounded uniformly in $\delta$ and converges pointwise to $\widehat\varphi(0)\cdot(2\pi)^{1/2}=1$ as $\delta\downarrow0$ because $\int\varphi=1$. Since $|m_\delta(\xi)-1|^2\langle\xi\rangle\,|\widehat f(\xi)|^2$ is dominated by an integrable function, dominated convergence in the Fourier representation of \cref{thm:fourier-gagliardo} gives strong convergence in $H^{1/2}$. A continuous function vanishing at infinity is uniformly continuous, so the standard approximate-identity argument gives uniform convergence. For the local version, first multiply by a cutoff identically one on the compact set of interest.
\end{proof}

\section{The zero-set regularizers}

\begin{definition}[Zero-set regularizer $T_\eps$]\label{def:Teps-regularizer}
For $\eps>0$ define $T_\eps:\C\to\C$ by
\begin{equation}
 T_\eps(z)=\frac{\eps z}{|z|^2+\eps},
 \qquad z\in\C.
 \label{eq:Teps-def}
\end{equation}
\end{definition}

\begin{lemma}[$T_\eps$ properties and strong critical convergence]\label{lem:Teps}
\uses{def:Teps-regularizer, def:Hs-loc}
The maps $T_\eps$ are uniformly Lipschitz (uniformly in $\eps$), satisfy $T_\eps(0)=0$, and obey
\begin{equation}
 |T_\eps(z)|\le\frac{\sqrt\eps}{2}
 \qquad\text{and}\qquad
 |T_\eps(z)|=\frac{\eps|z|}{|z|^2+\eps}\le|z|.
 \label{eq:Teps-bounds}
\end{equation}
If $f\in H^{1/2}_{\loc}(\R)\cap C(\R)$ is locally bounded, then
\begin{equation}
 T_\eps(f)\longrightarrow0
 \qquad\text{strongly in }H^{1/2}_{\loc}(\R)
 \text{ as }\eps\downarrow0.
 \label{eq:Teps-strong}
\end{equation}
\end{lemma}

\begin{proof}
\uses{def:Teps-regularizer, thm:fourier-gagliardo}
Strategy: verify the uniform derivative and sup bounds by treating $T_\eps$ as a radial map on $\R^2$, then obtain \eqref{eq:Teps-strong} by dominated convergence in both the $L^2$ and the Gagliardo pieces.

Viewed on $\R^2$, $T_\eps$ is radial. Its tangential derivative equals $\eps/(r^2+\eps)$ and its radial derivative equals $\eps(\eps-r^2)/(r^2+\eps)^2$; both have absolute value at most one, giving the uniform Lipschitz bound. The maximum of $\eps r/(r^2+\eps)$ is attained at $r=\sqrt\eps$ and equals $\sqrt\eps/2$, and $\eps|z|/(|z|^2+\eps)\le|z|$, proving \eqref{eq:Teps-bounds}. For the convergence, localize so that $f\in H^{1/2}(\R)\cap C_0(\R)$. For each fixed $z$, $T_\eps(z)\to0$; since $|T_\eps(f)|^2\le|f|^2$, dominated convergence gives $T_\eps(f)\to0$ in $L^2$. For the Gagliardo seminorm, the uniform Lipschitz bound gives
\[
 \frac{|T_\eps(f(x))-T_\eps(f(y))|^2}{|x-y|^2}
 \le\frac{|f(x)-f(y)|^2}{|x-y|^2},
\]
whose right side is integrable by \cref{thm:fourier-gagliardo}; pointwise convergence plus dominated convergence on $\R^2$ sends the seminorm to zero. This yields strong $H^{1/2}$ convergence, and the local statement follows by multiplying $f$ by a cutoff equal to one on the compact set considered.
% FULL PROOF: paper.tex L2348-L2372
\end{proof}

\section{Quadratic-form distributions}

\begin{definition}[Multiplier test algebra $\mathcal A(K)$]\label{def:multiplier-test-algebra}
\uses{def:Hs-space}
For a compact interval $K\subset\R$, let $\mathcal A(K)$ be the space of functions $\psi\in H^{1/2}(\R)\cap L^\infty(\R)$ supported in $K$, equipped with the norm
\[
 \norm{\psi}_{\mathcal A}=\norm{\psi}_{H^{1/2}}+\norm{\psi}_{\infty}.
\]
Smooth compactly supported functions are dense in $\mathcal A(K)$ for the simultaneous topology whenever $\psi$ has a continuous representative; the required mollification statement is \cref{lem:mollification-critical-uniform}.
\end{definition}

\begin{definition}[The current $\overline f f_x$ as an $H^{-1/2}$ distribution]\label{def:quadratic-form-current}
\uses{def:Hs-neg-pairing, def:multiplier-test-algebra, def:Hs-loc}
If $f\in H^{1/2}_{\loc}(\R)$, define the local distribution $\overline f f_x$ by
\begin{equation}
 \pair{\overline f f_x}{\phi}
 :=\pair{f_x}{\phi\overline f}_{H^{-1/2},H^{1/2}},
 \qquad \phi\in C_c^\infty(\R).
 \label{eq:quad-form}
\end{equation}
The pairing is complex bilinear, extending $\int T\phi$; it is not the sesquilinear Hilbert inner product. By \cref{lem:smooth-multiplier-quantitative} the test factor $\phi\overline f$ belongs to $H^{1/2}(\R)$, so the right-hand side is well defined through the $H^{-1/2}$–$H^{1/2}$ duality of \cref{def:Hs-neg-pairing}. Local boundedness of $f$ is not needed for this quadratic form; it enters only when zero-set regularizers are multiplied into test factors.
\end{definition}

\begin{lemma}[Cutoff-independence and locality of the current]\label{lem:current-locality}
\uses{def:quadratic-form-current, def:Hs-loc}
The distribution $\overline f f_x$ of \cref{def:quadratic-form-current} is independent of the smooth cutoff used to globalize $f$: any two cutoffs equal to one near $\operatorname{supp}\phi$ give the same value in \eqref{eq:quad-form}, because the distributional derivative is local. Moreover, if on a compact set $K$ one has $f_n\to f$ in $H^{1/2}$ after multiplication by one fixed cutoff equal to one near $K$, then
\[
 \overline{f_n}\,(f_n)_x\longrightarrow\overline f f_x
 \qquad\text{in }\D'(K).
\]
The same conclusion holds after integration in time when the localized $H^{1/2}$ convergence is in $L^2_t$.
\end{lemma}

\begin{proof}
\uses{def:quadratic-form-current, lem:smooth-multiplier-quantitative}
Cutoff-independence is immediate: for $\phi\in C_c^\infty$ the two globalizations of $f$ agree on a neighborhood of $\operatorname{supp}\phi$, and locality of $\partial_x$ makes the pairings coincide. For stability, fix $\phi$ supported in $K$ and subtract the two pairings:
\[
 \pair{(f_n)_x-f_x}{\phi\overline{f_n}}
 +\pair{f_x}{\phi(\overline{f_n}-\overline f)}.
\]
Boundedness of multiplication by the fixed smooth $\phi$ from $H^{1/2}$ to $H^{1/2}$, provided by \cref{lem:smooth-multiplier-quantitative}, bounds the absolute value by
\[
 C_\phi\,\norm{f_n-f}_{H^{1/2}}
 \bigl(\norm{f_n}_{H^{1/2}}+\norm{f}_{H^{1/2}}\bigr)
\]
after localization, which tends to zero. In time, integrate the same bound and apply Cauchy–Schwarz to the two $H^{1/2}$ factors.
\end{proof}

\begin{lemma}[Quadratic-form continuity bound and extension]\label{lem:quadratic-form-extension}
\uses{def:quadratic-form-current, def:multiplier-test-algebra}
Let $f\in H^{1/2}(\R)\cap L^\infty(\R)$. The map initially defined on smooth $\psi$ by
\[
 \mathcal Q_f(\psi)=\pair{f_x}{\psi\overline f}_{H^{-1/2},H^{1/2}}
\]
extends continuously to $\mathcal A(K)$, with
\begin{equation}
 |\mathcal Q_f(\psi)|
 \le C\,\norm{f}_{H^{1/2}}
 \left(\norm{\psi}_{\infty}\,\norm{f}_{H^{1/2}}
       +\norm{f}_{\infty}\,\norm{\psi}_{H^{1/2}}\right).
 \label{eq:quadratic-extension-bound}
\end{equation}
If $f_n\to f$ in both $H^{1/2}$ and $L^\infty$, and $\psi_n\to\psi$ in both norms, then $\mathcal Q_{f_n}(\psi_n)\to\mathcal Q_f(\psi)$. The same assertions hold locally after one fixed cutoff.
\end{lemma}

\begin{proof}
\uses{lem:gagliardo-product, lem:smooth-multiplier-quantitative, def:Hs-neg-pairing}
Strategy: bound $\mathcal Q_f(\psi)$ by pairing $f_x\in H^{-1/2}$ against the product $\psi\overline f\in H^{1/2}$, estimated by the algebra bound; joint convergence then follows by a telescoping argument.

The derivative map is bounded from $H^{1/2}$ to $H^{-1/2}$ (a consequence of the multiplier bounds behind \cref{lem:smooth-multiplier-quantitative}), so $\norm{f_x}_{H^{-1/2}}\le C\norm{f}_{H^{1/2}}$. By the product estimate \cref{lem:gagliardo-product},
\[
 \norm{\psi\overline f}_{H^{1/2}}
 \le C\left(\norm{\psi}_{\infty}\,\norm{f}_{H^{1/2}}
           +\norm{f}_{\infty}\,\norm{\psi}_{H^{1/2}}\right).
\]
Pairing the two Sobolev spaces via \cref{def:Hs-neg-pairing} gives \eqref{eq:quadratic-extension-bound}, and since smooth $\psi$ are dense in $\mathcal A(K)$ by \cref{def:multiplier-test-algebra}, the map extends continuously. For joint convergence, add and subtract $\pair{(f_n)_x}{\psi_n\overline f}$ and $\pair{(f_n)_x}{\psi\overline f}$ and apply the same product estimate to each difference; uniform boundedness of the convergent sequences forces each term to zero. Localization is obtained by inserting a cutoff equal to one near $K$, two choices giving the same value by locality of the distributional derivative.
% FULL PROOF: paper.tex L2420-L2436
\end{proof}


\chapter{The density determines the current}\label{chap:current}

\section{Density equality at every time}

\begin{corollary}[Equal density at every time]\label{cor:all-time-density}
\uses{thm:GWP, cor:C0-slices}
Let $u,v\in C(\R;L^2(\R))$ be the global $L^2$ solutions of the cubic
equation with the equal-modulus hypothesis
\[
 |u(t,x)|=|v(t,x)|\qquad\text{for a.e. }(t,x)\in\R\times\R .
\]
Then
\begin{equation}
 |u(t)|^2=|v(t)|^2\quad\text{in }L^1(\R)
 \qquad\text{for \emph{every} }t\in\R. \label{eq:all-time-density}
\end{equation}
\end{corollary}

\begin{proof}
\uses{thm:GWP, cor:C0-slices}
The map $t\mapsto|f(t)|^2$ is continuous from $\R$ to $L^1(\R)$ whenever
$f\in C(\R;L^2(\R))$, by the elementary factorization
\[
 \norm{|f(t)|^2-|f(s)|^2}_{L^1}
 \le\norm{f(t)-f(s)}_2\bigl(\norm{f(t)}_2+\norm{f(s)}_2\bigr).
\]
Applying this to $u$ and to $v$ (which lie in $C(\R;L^2)$ by
\cref{thm:GWP,cor:C0-slices}), the difference $t\mapsto|u(t)|^2-|v(t)|^2$ is
continuous into $L^1(\R)$. Fubini's theorem and the a.e.\ equal-modulus
hypothesis give $|u(t)|^2=|v(t)|^2$ in $L^1(\R)$ for almost every $t$. A
continuous map that vanishes almost everywhere vanishes everywhere, which is
\eqref{eq:all-time-density}.
\end{proof}

\section{The smooth continuity equation}

For a solution set $\rho_u=|u|^2$ and $j_u=\Ima(\overline u\,u_x)$.

\begin{lemma}[Smooth continuity equation]\label{lem:smooth-continuity-equation}
\uses{def:propagator}
Every $H^1$ solution $u$ of $u_t=i u_{xx}-i\sigma|u|^2u$ satisfies
\begin{equation}
 \partial_t\rho_u+2\partial_xj_u=0
 \qquad\text{in }\mathcal{D}'(I\times\R). \label{eq:continuity-eq}
\end{equation}
Consequently, if $\phi\in C_c^\infty(\R)$ and
$a(x)=\int_{-\infty}^x\phi(y)\,\dd y$, then
\begin{equation}
 \frac{\dd}{\dd t}\int_\R a(x)\,|u(t,x)|^2\,\dd x
 =2\int_\R\phi(x)\,j_u(t,x)\,\dd x \label{eq:localized-mass-smooth}
\end{equation}
holds in the distributional sense in time.
\end{lemma}

\begin{proof}
\uses{lem:gelfand-chain-rule}
Strategy: apply the Gelfand-triple mass chain rule to the density weighted by
a real cutoff, substitute the equation, and pass from compactly supported
weights to the bounded primitive $a$.

For real $b\in C_c^\infty(\R)$, multiplication by $b$ is bounded on $H^1$, and
the time-mollification argument of \cref{lem:gelfand-chain-rule} applied to the
bounded sesquilinear form $B_b(f,g)=\int b\,f\overline g$ gives
$\frac{\dd}{\dd t}\int b|u|^2=2\Rea\pair{u_t}{bu}_{H^{-1},H^1}$ in
distributions in time. Substituting $u_t=i u_{xx}-i\sigma|u|^2u$, the
nonlinear pairing is purely imaginary, and weak integration by parts yields
$2\Rea\pair{i u_{xx}}{bu}=2\int b'\,\Ima(\overline u\,u_x)$; testing in time
proves \eqref{eq:continuity-eq}. Taking $b=b_R=a\chi_R$ with cutoffs $\chi_R$
of uniformly bounded derivative converging pointwise to zero, the term
$a\chi_R'$ vanishes in the limit (as $a$ is bounded and $u\,u_x\in L^1$) while
$a'\chi_R=\phi\chi_R$ converges to the right side of
\eqref{eq:localized-mass-smooth}.
% FULL PROOF: paper.tex L2545–L2573
\end{proof}

\section{Endpoint current as a spacetime distribution}

\begin{definition}[Endpoint current functional]\label{def:endpoint-current}
\uses{def:quadratic-form-current, def:Hs-neg-pairing, lem:current-locality}
For a global $L^2$ solution $u$ and almost every $t$, define the endpoint
current $\mathfrak j_u(t)$ acting on $\phi\in C_c^\infty(\R)$ by
\begin{equation}
 \pair{\mathfrak j_u(t)}{\phi}
 =\Ima\pair{u_x(t)}{\phi\,\overline{u(t)}}_{H^{-1/2},H^{1/2}}.
 \label{eq:endpoint-current-def}
\end{equation}
For any cutoff $\chi\in C_c^\infty(\R)$ equal to one on $\supp\phi$ one has the
bound
\[
 \bigl|\pair{\mathfrak j_u(t)}{\phi}\bigr|
 \le C_\phi\,\norm{\chi u(t)}_{H^{1/2}}^2 ,
\]
whose right-hand side is integrable in time; hence $t\mapsto\mathfrak j_u(t)$
is an $L^1_{\loc}$ family of distributions.
\end{definition}

\begin{theorem}[Endpoint recovery of the logarithmic derivative]
\label{thm:current-recovery}
\uses{def:endpoint-current, def:quadratic-form-current}
Let $u,v$ be the global $L^2$ solutions with the equal-modulus hypothesis of
\cref{cor:all-time-density}. Then
\begin{equation}
 \overline u\,u_x=\overline v\,v_x
 \qquad\text{as a local quadratic-form distribution on }\R_t\times\R_x .
 \label{eq:log-derivative-equality}
\end{equation}
Moreover there is a full-measure set of times on which the same equality holds
as a spatial distribution.
\end{theorem}

\begin{proof}
\uses{def:endpoint-current, lem:smooth-continuity-equation, lem:realpart-logform, lem:tensor-test-density, lem:distribution-slicing, lem:quadratic-form-extension, cor:all-time-density, prop:endpoint, cor:NLS-smoothing}
Strategy: pass the smooth localized mass identity to the $L^2$ limit to obtain
a weak localized mass identity, use equality of densities to match the two
currents against separated tests, upgrade to arbitrary spacetime tests by a
tensor-density argument, and finally slice in time.

\emph{Weak localized mass identity.} Take Schwartz data $u_{0,n}\to u_0$ with
smooth solutions $u_n$. For $\eta\in C_c^\infty(I)$, $\phi\in C_c^\infty(\R)$
and the bounded primitive $a(x)=\int_{-\infty}^x\phi$, identity
\eqref{eq:localized-mass-smooth} of \cref{lem:smooth-continuity-equation} gives
\begin{equation}
 -\int_I\eta'(t)\!\int_\R a(x)\,|u(t,x)|^2\,\dd x\,\dd t
 =2\int_I\eta(t)\,\pair{\mathfrak j_u(t)}{\phi}\,\dd t .
 \label{eq:weak-localized-mass}
\end{equation}
The left side converges since $u_n\to u$ in $C_tL^2_x$ and $a$ is bounded; the
right side converges by \cref{cor:NLS-smoothing,prop:endpoint}.

\emph{Matching the currents.} Applying \eqref{eq:weak-localized-mass} to $u$
and to $v$ and using $|u(t)|^2=|v(t)|^2$ from \cref{cor:all-time-density}, the
left sides agree, so $\pair{\mathfrak j_u(t)}{\phi}$ and
$\pair{\mathfrak j_v(t)}{\phi}$ agree against every tensor test
$\eta(t)\phi(x)$. For $\Psi\in C_c^\infty(I\times\R)$ with $\chi$ equal to one
on its spatial support, the multiplier algebra estimate gives
\begin{equation}
 \Bigl|\int_I\pair{u_x(t)}{\Psi(t)\overline{u(t)}}\,\dd t\Bigr|
 \le C_K\max_{|\alpha|\le1}\norm{\partial_x^\alpha\Psi}_\infty
       \norm{\chi u}_{L^2_tH^{1/2}_x}^2, \label{eq:current-test-continuity}
\end{equation}
and likewise for $v$. By \cref{lem:tensor-test-density} with $M=1$, finite
sums of separated tests converge to $\Psi$ in the seminorm on the right of
\eqref{eq:current-test-continuity}; hence the imaginary parts of the two forms
agree against every spacetime test. The real parts agree from equality of the
densities via \cref{lem:realpart-logform}, whose real part depends only on the common $\abs u^2$.
This gives \eqref{eq:log-derivative-equality} as a spacetime distribution.

\emph{Slicing.} Exhaust spacetime by compact $I_m\times K_m$ with cutoffs
$\chi_m$ equal to one near $K_m$. For $\phi\in C_c^\infty(K_m)$ the multiplier
proof gives, with
$\Lambda_t(\phi)=\pair{u_x(t)}{\phi\overline{u(t)}}
 -\pair{v_x(t)}{\phi\overline{v(t)}}$,
\begin{equation}
 |\Lambda_t(\phi)|
 \le C_m\bigl(\norm{\chi_m u(t)}_{H^{1/2}}^2
              +\norm{\chi_m v(t)}_{H^{1/2}}^2\bigr)
       \bigl(\norm\phi_\infty+\norm{\phi'}_\infty\bigr),
 \label{eq:sliced-current-test-bound}
\end{equation}
with integrable coefficient on $I_m$. Fixing a countable $C^1$-dense family
$\{\phi_{m,n}\}\subset C_c^\infty(K_m)$, the spacetime identity gives
$\int\eta\,\Lambda_t(\phi_{m,n})\,\dd t=0$ for all scalar $\eta$, so
$\Lambda_t(\phi_{m,n})=0$ off a null set; removing the countable union and
using \eqref{eq:sliced-current-test-bound} extends the equality to every test
supported in $K_m$, and then to all of spacetime. At these good times the
continuous extensions to the multiplier algebra agree by
\cref{lem:quadratic-form-extension}.
% FULL PROOF: paper.tex L2599–L2671
\end{proof}

\begin{corollary}[Equal density $\Rightarrow$ equal current]\label{cor:density-equal-current}
\uses{def:endpoint-current}
Under the equal-modulus hypothesis of \cref{cor:all-time-density}, the endpoint
currents agree:
\[
 \mathfrak j_u(t)=\mathfrak j_v(t)
 \qquad\text{as spatial distributions, for a.e.\ }t,
\]
and $\Ima(\overline u\,u_x)=\Ima(\overline v\,v_x)$ as a spacetime
distribution.
\end{corollary}

\begin{proof}
\uses{thm:current-recovery, cor:all-time-density}
By \cref{cor:all-time-density} the densities are equal, so the hypotheses of
\cref{thm:current-recovery} hold. Taking imaginary parts in
\eqref{eq:log-derivative-equality} gives
$\Ima(\overline u\,u_x)=\Ima(\overline v\,v_x)$ as a spacetime distribution,
i.e.\ $\pair{\mathfrak j_u(t)}{\phi}=\pair{\mathfrak j_v(t)}{\phi}$ for every
$\phi\in C_c^\infty(\R)$ on the full-measure set of good times produced by
\cref{thm:current-recovery}.
\end{proof}


\chapter{Wronskian and the exterior product}\label{chap:wronskian-exterior}

\begin{lemma}[Simultaneous smooth approximation for the Wronskian]\label{lem:wronskian-smooth-approx}
\uses{def:Hs-space, def:Hs-loc}
Let $f,g\in H^{1/2}_{\loc}(\R)\cap C(\R)$ be locally bounded and let
$\chi\in C_c^\infty(\R)$ be real. Then, as local
$H^{-1/2}$--$H^{1/2}$ quadratic-form distributions,
\begin{equation}
 \overline{\chi f}(\chi f)'-\overline{\chi g}(\chi g)'
 =\chi^2\bigl(\overline f f'-\overline g g'\bigr)
  +\chi\chi'\bigl(|f|^2-|g|^2\bigr),
\end{equation}
and the classical Wronskian factors as $W(\chi f,\chi g)=\chi^2\,W(f,g)$,
where $W(f,g)=gf'-fg'$. Consequently, if $|f|=|g|$ almost everywhere and
$\overline f f'-\overline g g'=0$, then $\chi f,\chi g\in H^{1/2}(\R)\cap
L^\infty(\R)$, they still satisfy $\overline{\chi f}(\chi f)'
-\overline{\chi g}(\chi g)'=0$, and their continuous representatives obey
$|\chi f|=|\chi g|$ at every point.
\end{lemma}

\begin{proof}
\uses{prop:endpoint, cor:NLS-smoothing}
The product rule for the quadratic form gives the two derivative terms
$\chi\chi'(|f|^2-|g|^2)$ and the factor $\chi^2$ on the intrinsic form; the
cross terms cancel because $\chi$ is real. Under $|f|=|g|$ the term
$\chi\chi'(|f|^2-|g|^2)$ vanishes, and under
$\overline f f'-\overline g g'=0$ the remaining $\chi^2$ term vanishes, so
the localized form is $0$. Multiplication by the compactly supported cutoff
$\chi$ places $\chi f,\chi g$ in $H^{1/2}\cap L^\infty$; the good-slice
regularity supplied by \cref{prop:endpoint} and \cref{cor:NLS-smoothing}
makes the chosen representatives continuous, so the almost-everywhere
identity $|f|=|g|$ holds pointwise, hence at every point after cutoff.
\end{proof}

\begin{lemma}[Zero-set-safe multipliers, no division]\label{lem:wronskian-zero-set-safe}
\uses{def:Hs-space, def:Teps-regularizer}
Let $f,g\in H^{1/2}(\R)\cap L^\infty(\R)$ be continuous and compactly
supported with $\rho:=|f|^2=|g|^2$ everywhere. For $\eps>0$ set
\[
 s_\eps=\frac{\rho}{\rho+\eps},
 \qquad
 h_\eps=\frac{fg}{\rho+\eps}.
\]
Then $s_\eps,h_\eps\in H^{1/2}(\R)\cap L^\infty(\R)$, and
$s_\eps f,\ s_\eps g\in H^{1/2}(\R)\cap L^\infty(\R)$. Moreover
\begin{equation}
 \overline f\,h_\eps=s_\eps g,
 \qquad
 \overline g\,h_\eps=s_\eps f .
\end{equation}
\end{lemma}

\begin{proof}
\uses{lem:lipschitz-composition, lem:gagliardo-product}
The map $r\mapsto r/(r+\eps)$ is bounded, Lipschitz on the bounded range of
$\rho$, and vanishes at $r=0$; by \cref{lem:lipschitz-composition}
$s_\eps\in H^{1/2}\cap L^\infty$, and the product estimate
\cref{lem:gagliardo-product} gives $s_\eps f,s_\eps g\in H^{1/2}\cap
L^\infty$. For $h_\eps$ the reciprocal $r\mapsto(r+\eps)^{-1}$ does not
vanish at $0$, so fix a cutoff $\kappa\in C_c^\infty$ equal to $1$ on
$\supp f\cup\supp g$. Since $r\mapsto(r+\eps)^{-1}-\eps^{-1}$ is Lipschitz
on the range of $\rho$ and vanishes at $0$,
\[
 \kappa(\rho+\eps)^{-1}
 =\frac{\kappa}{\eps}+\kappa\bigl((\rho+\eps)^{-1}-\eps^{-1}\bigr)
 \in H^{1/2}\cap L^\infty ,
\]
and because $fg$ is supported where $\kappa=1$,
$h_\eps=fg\,\kappa(\rho+\eps)^{-1}\in H^{1/2}\cap L^\infty$ by
\cref{lem:gagliardo-product}. Finally
$\overline f h_\eps=\overline f fg/(\rho+\eps)=\rho g/(\rho+\eps)=s_\eps g$,
and symmetrically $\overline g h_\eps=s_\eps f$.
\end{proof}

\begin{lemma}[Current form of the Wronskian pairing]\label{lem:wronskian-current-form}
\uses{def:quadratic-form-current, def:Teps-regularizer}
Under the hypotheses of \cref{lem:wronskian-zero-set-safe}, let
$D=gf'-fg'$ and $T_\eps(z)=\eps z/(|z|^2+\eps)$. Then for every
$\phi\in C_c^\infty(\R)$,
\begin{equation}
 \pair{f'}{s_\eps g\phi}-\pair{g'}{s_\eps f\phi}=0,
\end{equation}
and hence
\begin{equation}
 \pair{D}{\phi}
 =\pair{f'}{T_\eps(g)\phi}-\pair{g'}{T_\eps(f)\phi}. \label{eq:m10-W-Teps}
\end{equation}
Letting $\eps\to0$ gives $\pair{D}{\phi}=0$; that is, $D=gf'-fg'=0$ in
$\D'(\R)$.
\end{lemma}

\begin{proof}
\uses{lem:quadratic-form-extension, lem:mollification-critical-uniform, lem:Teps, cor:density-equal-current}
Test the intrinsic identity $\overline f f'-\overline g g'=0$ against
$h_\eps\phi$. The function $h_\eps\phi$ is continuous, compactly supported
and lies in $H^{1/2}\cap L^\infty$; by \cref{lem:quadratic-form-extension}
(whose convergence is furnished by the simultaneous $H^{1/2}$-and-$L^\infty$
mollification \cref{lem:mollification-critical-uniform}) the equality of
quadratic forms extends from smooth tests to $h_\eps\phi$. Using
$\overline f h_\eps=s_\eps g$ and $\overline g h_\eps=s_\eps f$ from
\cref{lem:wronskian-zero-set-safe} yields the first display. Subtracting it
from the defining pairing of $D$ and using $|f|=|g|$ to write
$1-s_\eps=\eps/(\rho+\eps)$ gives \eqref{eq:m10-W-Teps}, with
$(1-s_\eps)g=T_\eps(g)$ and $(1-s_\eps)f=T_\eps(f)$. By \cref{lem:Teps}
both $T_\eps(f),T_\eps(g)\to0$ strongly in local $H^{1/2}$; multiplication
by $\phi$ preserves this, and since $f',g'\in H^{-1/2}_{\loc}$ the right
side of \eqref{eq:m10-W-Teps} tends to $0$. The intrinsic-current input
$\overline f f'-\overline g g'=0$ used here is the sliced equal-density
current identity of \cref{cor:density-equal-current}.
\end{proof}

\begin{theorem}[Zero-set-safe Wronskian lemma]\label{thm:critical-wronskian}
\uses{def:Hs-space, def:Hs-loc, def:Hs-neg-pairing}
Let $f,g\in H^{1/2}_{\loc}(\R)\cap C(\R)$ be locally bounded. Assume
\begin{equation}
 |f|=|g|\quad\text{almost everywhere}                \label{eq:W-density}
\end{equation}
and
\begin{equation}
 \overline f f'-\overline g g'=0                     \label{eq:W-log}
\end{equation}
as a local $H^{-1/2}$--$H^{1/2}$ quadratic-form distribution. Then
\begin{equation}
 g f'-f g'=0\quad\text{in }\D'(\R).                  \label{eq:W-zero}
\end{equation}
\end{theorem}

\begin{proof}
\uses{lem:wronskian-smooth-approx, lem:wronskian-zero-set-safe, lem:wronskian-current-form}
Both hypotheses and the conclusion are local, so by
\cref{lem:wronskian-smooth-approx} we cut off and reduce to $f,g\in
H^{1/2}(\R)\cap L^\infty(\R)$ continuous and compactly supported with
$\rho:=|f|^2=|g|^2$ everywhere and $\overline f f'-\overline g g'=0$. For
$\eps>0$ the zero-set-safe multipliers $s_\eps,h_\eps$ of
\cref{lem:wronskian-zero-set-safe} are legal $H^{1/2}\cap L^\infty$ test
factors that use no quotient near the common zero set. Testing the intrinsic
identity against $h_\eps\phi$ and passing to the limit $\eps\to0$ via
\cref{lem:wronskian-current-form} yields $\pair{gf'-fg'}{\phi}=0$ for every
$\phi\in C_c^\infty(\R)$, which is \eqref{eq:W-zero}.
% FULL PROOF: paper.tex L2696-L2781
\end{proof}

\begin{corollary}[Diagonal Wronskian for the NLS pair]\label{cor:diagonal-W}
\uses{thm:critical-wronskian}
There is a full-measure set $E\subset\R$ such that, for every $t\in E$,
\begin{equation}
 u(t)v_x(t)-v(t)u_x(t)=0\quad\text{in }\D'(\R_x).    \label{eq:time-W}
\end{equation}
The set $E$ may be chosen so that the good-slice regularity and the sliced
logarithmic-derivative equality also hold at every $t\in E$.
\end{corollary}

\begin{proof}
\uses{thm:critical-wronskian, thm:current-recovery, lem:countable-full-measure, cor:density-equal-current}
On the countable full-measure set of good times of
\cref{lem:countable-full-measure}, the slices $u(t),v(t)$ lie in
$H^{1/2}_{\loc}\cap C$ and are locally bounded, satisfy $|u(t)|=|v(t)|$
almost everywhere, and by the sliced current identity of
\cref{thm:current-recovery} and \cref{cor:density-equal-current} obey
$\overline{u(t)}u_x(t)-\overline{v(t)}v_x(t)=0$ as a local quadratic form.
Applying \cref{thm:critical-wronskian} at each such time gives
\eqref{eq:time-W}. Intersecting finitely many full-measure sets preserves
full measure.
\end{proof}

\section{The two-particle exterior product}

\begin{definition}[Two-particle exterior product]\label{def:exterior-product}
\uses{def:fixed-point-space}
Let $u,v$ be the two global $L^2$ solutions of the cubic NLS. Define
\begin{equation}
 F(t,x,y)=u(t,x)v(t,y)-v(t,x)u(t,y).                 \label{eq:F-def}
\end{equation}
Then, by the tensor-norm identity and continuity of $u,v$ in $L^2$,
\begin{equation}
 F\in C(\R;L^2(\R^2)),                               \label{eq:F-C-L2}
\end{equation}
and $F$ is antisymmetric,
\begin{equation}
 F(t,y,x)=-F(t,x,y).                                 \label{eq:F-antisym}
\end{equation}
We write $\rho(t,x)=|u(t,x)|^2=|v(t,x)|^2$.
\end{definition}

\begin{proposition}[Exterior-product equation]\label{prop:F-equation}
\uses{def:exterior-product, cor:diagonal-W}
In distributions on $\R_t\times\R_x^2$, with $\sigma$ the NLS nonlinearity
sign,
\begin{equation}
 iF_t+F_{xx}+F_{yy}
 =\sigma\bigl(\rho(t,x)+\rho(t,y)\bigr)F.            \label{eq:F-equation}
\end{equation}
On every bounded interval $I$,
\begin{align}
 F&\in L^4(I\times\R^2),                             \label{eq:F-L4}\\
 \bigl(\rho(x)+\rho(y)\bigr)F&\in L^2(I\times\R^2).  \label{eq:QF-L2}
\end{align}
\end{proposition}

\begin{proof}
\uses{lem:tensor-norm, prop:endpoint, prop:mass-L2, thm:GWP}
\emph{Strategy.} Establish the unfactored tensor identity for smooth
approximating solutions, pass to the limit in $L^1_tL^2_{x,y}$, and only
then use the exact equal-moduli hypothesis to factor the forcing.

For a smooth pair, direct differentiation with $N(z)=|z|^2z$ gives the
unfactored identity
\begin{align}
 &(i\partial_t+\partial_x^2+\partial_y^2)\bigl(u(x)v(y)-v(x)u(y)\bigr)\notag\\
 &\quad=\sigma\bigl[N(u)(x)v(y)+u(x)N(v)(y)-N(v)(x)u(y)-v(x)N(u)(y)\bigr],
                                                     \label{eq:F-unfactored}
\end{align}
which is stable under arbitrary smooth approximation. Choosing smooth data
converging to the given data, the solutions satisfy $u_n\to u$, $v_n\to v$
in $C_tL^2_x$ and $N(u_n)\to N(u)$, $N(v_n)\to N(v)$ in $L^1_tL^2_x$, so each
tensor on the right of \eqref{eq:F-unfactored} converges in
$L^1_tL^2_{x,y}$ and the left side converges in distributions by
$C_tL^2_{x,y}$ convergence of the exterior products (\cref{lem:tensor-norm}).
Hence \eqref{eq:F-unfactored} holds for the rough solutions. Using
$|u|^2=|v|^2=\rho$, its four terms factor as
$\rho(x)F+\rho(y)F$, proving \eqref{eq:F-equation}. The bound
\eqref{eq:F-L4} follows from
$\norm{u(t,x)v(t,y)}_{L^4_{t,x,y}}^4=\int_I\norm{u(t)}_4^4\norm{v(t)}_4^4\dd t
\le\norm u_{L^8_tL^4_x(I)}^4\norm v_{L^8_tL^4_x(I)}^4$, finite by
\cref{prop:endpoint}. For \eqref{eq:QF-L2} a typical forcing term
$|u(x)|^2u(x)v(y)$ has squared $L^2$ norm
$\int_I\norm{u(t)}_6^6\norm{v(t)}_2^2\dd t$, finite by mass conservation
(\cref{prop:mass-L2}, \cref{thm:GWP}) and the $L^6_{t,x}$ bound.
% FULL PROOF: paper.tex L2846-L2898
\end{proof}

\section{Normal and tangential coordinates}

\begin{definition}[Normal/tangential coordinates to the diagonal]\label{def:normal-tangential-coords}
\uses{def:exterior-product, prop:F-equation}
Set
\begin{equation}
 r=\frac{x+y}{\sqrt2},\qquad s=\frac{x-y}{\sqrt2},  \label{eq:rs-coordinates}
\end{equation}
an orthogonal change of variables, and keep the letter $F$ for the
transformed function. Then
\begin{equation}
 iF_t+F_{rr}+F_{ss}=QF,                              \label{eq:F-rs}
\end{equation}
where the ridge potential is
\begin{equation}
 Q(t,r,s)=\sigma\left[
  \rho\!\left(t,\tfrac{r+s}{\sqrt2}\right)
 +\rho\!\left(t,\tfrac{r-s}{\sqrt2}\right)\right].   \label{eq:Q-ridge}
\end{equation}
Here $Q$ is real and even in $s$, while $F$ is odd in $s$.
\end{definition}

\begin{lemma}[Ridge-potential integrability bounds]\label{lem:Q-bounds}
\uses{def:normal-tangential-coords, prop:F-equation}
On every bounded time interval $I$,
\begin{align}
 Q&\in L^1(I;L^\infty(\R^2)),                        \label{eq:Q-L1Linfty}\\
 \norm Q_{L^1_tL^\infty_{r,s}(I)}
 &\le 2|\sigma|\,|I|^{1/2}\norm u_{L^4_tL^\infty_x(I)}^2.
                                                     \label{eq:Q-L1-bound}
\end{align}
Furthermore, for every strip $S_{a,h}=\{(r,s):a<s<a+h\}$,
\begin{equation}
 \norm{\1_{S_{a,h}}Q}_{L^2(I\times\R^2)}
 \le C|\sigma|\,h^{1/2}\norm u_{L^4_tL^4_x(I)}^2,    \label{eq:Q-strip}
\end{equation}
uniformly in $a$.
\end{lemma}

\begin{proof}
\uses{prop:endpoint}
At each time $\norm{Q(t)}_\infty\le2|\sigma|\norm{u(t)}_\infty^2$;
integrating over $I$ and applying Cauchy--Schwarz in time gives
\eqref{eq:Q-L1Linfty}--\eqref{eq:Q-L1-bound}, finite by the $L^4_tL^\infty_x$
endpoint bound (\cref{prop:endpoint}). For one ridge term, the change of
variables in $r$ gives
\[
 \int_I\!\int_a^{a+h}\!\int_\R
 \rho\!\left(t,\tfrac{r+s}{\sqrt2}\right)^2\dd r\,\dd s\,\dd t
 =\sqrt2\,h\int_I\norm{u(t)}_4^4\,\dd t .
\]
The second ridge is identical, and the triangle inequality yields
\eqref{eq:Q-strip}.
\end{proof}


\chapter{Distribution-valued diagonal traces}\label{chap:traces}

The equal-density condition gives the Dirichlet trace $F|_{s=0}=0$ formally, and
the Wronskian gives the Neumann trace. At $L^2$ regularity neither trace is
initially classical. We construct them by testing in the tangential variables and
using the equation as an ordinary differential equation in the normal variable.

\section{The one-dimensional trace theorem}

\begin{lemma}[Interior first-derivative estimate]\label{lem:interior-first-derivative}
\uses{def:Hs-space}
Let $J'\Subset J$ be bounded intervals. There is $C=C(J',J)$ such that every
smooth $h$ on $J$ satisfies
\begin{equation}
 \norm{h'}_{L^2(J')}\le C\left(\norm h_{L^2(J)}+\norm{h''}_{L^2(J)}\right).
 \label{eq:interior-H1-estimate}
\end{equation}
\end{lemma}

\begin{proof}
Choose $\chi\in C_c^\infty(J)$ with $\chi=1$ on $J'$. Integration by parts gives
\[
 \int\chi^2\abs{h'}^2
 =-\Rea\int\chi^2h''\overline h-2\Rea\int\chi\chi'h'\overline h.
\]
Cauchy--Schwarz and $2ab\le\tfrac12a^2+2b^2$ absorb half of the left side and
give $\norm{\chi h'}_2^2\le C_\chi(\norm h_2^2+\norm{h''}_2^2)$. Since $\chi=1$ on
$J'$, this is \eqref{eq:interior-H1-estimate}.
\end{proof}

\begin{lemma}[Mollified weak-limit $H^2$ step]\label{lem:H2-weak-limit}
\uses{lem:interior-first-derivative}
Let $h\in\D'(J)$ and suppose $h$ and its distributional second derivative are
represented by functions in $L^2_{\loc}(J)$. Then $h'\in L^2_{\loc}(J)$, hence
$h\in H^2_{\loc}(J)$, and for all $J'\Subset J''\Subset J$,
\begin{equation}
 \norm h_{H^2(J')}\le C_{J',J''}\left(\norm h_{L^2(J'')}+\norm{h''}_{L^2(J'')}\right).
 \label{eq:local-H2-bound}
\end{equation}
\end{lemma}

\begin{proof}
\uses{lem:interior-first-derivative}
Mollify $h$ inside $J''$. The mollifications $h_\eps$ and $h_\eps''$ converge in
$L^2$ on every smaller interval and are uniformly bounded there. Applying
\cref{lem:interior-first-derivative} on nested intervals yields a uniform $L^2$
bound for $h_\eps'$. Weak compactness gives a weakly convergent subsequence, and
the distributional identity identifies its weak limit with $h'$; hence
$h'\in L^2_{\loc}(J)$. Since $h''\in L^2_{\loc}(J)$ by hypothesis,
$h\in H^2_{\loc}(J)$, and passing to the limit in the uniform bound gives
\eqref{eq:local-H2-bound}.
\end{proof}

\begin{theorem}[Local $H^2$ characterization and $C^1$ representative]\label{thm:local-H2-characterization}
\uses{def:Hs-space}
Let $h\in\D'(J)$. If $h$ and its distributional second derivative are represented
by functions in $L^2_{\loc}(J)$, then $h\in H^2_{\loc}(J)$; it has a unique $C^1$
representative, and on every $J'\Subset J$ the bound \eqref{eq:local-H2-bound}
holds for any intermediate interval $J'\Subset J''\Subset J$.
\end{theorem}

\begin{proof}
\uses{lem:H2-weak-limit, lem:H1-C0}
The membership $h\in H^2_{\loc}(J)$ and the estimate \eqref{eq:local-H2-bound}
are exactly \cref{lem:H2-weak-limit}. For the representative, a one-dimensional
$H^1$ function $g$ has an absolutely continuous representative with
$\abs{g(x)-g(y)}\le\abs{x-y}^{1/2}\norm{g'}_{L^2([x,y])}$ by
\cref{lem:H1-C0}. Applying this first to $h'$ and then to $h$ shows that $h'$ and
$h$ are continuous, so $h$ has a $C^1$ representative; two continuous
representatives agreeing almost everywhere agree everywhere.
% FULL PROOF: paper.tex L3009--L3026
\end{proof}

\section{Scalarization in the normal variable}

\begin{lemma}[Scalarization in the normal variable]\label{lem:scalarization-normal}
\uses{def:normal-tangential-coords, def:exterior-product}
Fix a bounded interval $I$ and $\psi\in C_c^\infty(I\times\R_r)$. Define the
scalar distribution in $s$
\begin{equation}
 H_\psi(s)=\pair{F(\cdot,\cdot,s)}{\psi}_{t,r}. \label{eq:Hpsi-def}
\end{equation}
Then $H_\psi\in H^2_{\loc}(\R_s)$; in particular $H_\psi(0)$ and $H_\psi'(0)$ are
well-defined via the $C^1$ representative.
\end{lemma}

\begin{proof}
\uses{prop:F-equation, thm:local-H2-characterization}
Since $F\in L^2(I\times\R^2)$, Cauchy--Schwarz in $(t,r)$ gives
$H_\psi\in L^2_{\loc}(\R_s)$. Pairing the equation $iF_t+F_{rr}+F_{ss}=QF$ of
\cref{prop:F-equation} with $\psi$ and integrating by parts in $t$ and $r$ gives
\begin{equation}
 H_\psi''(s)=\pair{QF}{\psi}_{t,r}+\pair{F}{i\psi_t-\psi_{rr}}_{t,r}.
 \label{eq:Hpsi-second}
\end{equation}
Both right-hand terms belong to $L^2_{\loc}(\R_s)$ because $QF,F\in L^2(I\times\R^2)$
and the test functions are compactly supported, so $H_\psi''\in L^2_{\loc}(\R_s)$
distributionally. \cref{thm:local-H2-characterization} then gives
$H_\psi\in H^2_{\loc}(\R_s)$ with its $C^1$ representative.
\end{proof}

\section{The Dirichlet trace}

\begin{definition}[Dirichlet trace functional]\label{def:dirichlet-trace}
\uses{lem:scalarization-normal, thm:local-H2-characterization}
For $\psi\in C_c^\infty(I\times\R)$ the \emph{Dirichlet trace} of $F$ is the
functional
\[
 \psi\longmapsto H_\psi(0),
\]
where $H_\psi(0)$ is the value at $s=0$ of the $C^1$ representative supplied by
\cref{lem:scalarization-normal}. It is denoted $F|_{s=0}$ as a distribution in
$(t,r)$.
\end{definition}

\begin{proposition}[Vanishing Dirichlet trace]\label{prop:Dirichlet-trace}
\uses{def:dirichlet-trace}
For every $\psi\in C_c^\infty(I\times\R)$,
\[
 H_\psi(0)=0.
\]
Equivalently, $F|_{s=0}=0$ as a distribution in $(t,r)$.
\end{proposition}

\begin{proof}
\uses{def:exterior-product, lem:scalarization-normal}
The antisymmetry of $F$ (from \cref{def:exterior-product}) becomes
$F(t,r,-s)=-F(t,r,s)$ in the coordinates $(r,s)$. Hence $H_\psi$ is odd. Its
continuous representative, supplied by $H^2_{\loc}\hookrightarrow C^1_{\loc}$ in
one dimension through \cref{lem:scalarization-normal}, vanishes at $s=0$.
\end{proof}

\section{The Neumann trace}

\begin{definition}[Neumann trace functional]\label{def:neumann-trace}
\uses{lem:scalarization-normal, thm:local-H2-characterization}
For $\psi\in C_c^\infty(I\times\R)$ the \emph{Neumann trace} of $F$ is the
functional
\[
 \psi\longmapsto H_\psi'(0),
\]
where $H_\psi'(0)$ is the derivative at $s=0$ of the $C^1$ representative supplied
by \cref{lem:scalarization-normal}. It is denoted $F_s|_{s=0}$ as a distribution
in $(t,r)$.
\end{definition}

\begin{lemma}[Nested-cutoff localization for the Neumann trace]\label{lem:neumann-localization}
\uses{def:neumann-trace}
Fix $\psi\in C_c^\infty(I\times\R)$ and $s_0>0$. Writing $z=(r-s)/\sqrt2$, the set
of $z$ for which $\psi(t,\sqrt2z+s)$ is nonzero for some $t\in I$ and
$\abs s\le s_0$ is compact. There exist $\chi_0,\chi_1\in C_c^\infty(\R)$ with
$\chi_0=1$ on that set and $\chi_1=1$ on an open neighbourhood of
$\supp\chi_0+[-\sqrt2s_0,\sqrt2s_0]$, so that whenever $\abs s\le s_0$ and
$\psi(t,\sqrt2z+s)\ne0$ both $\chi_1(z)$ and $\chi_1(z+\sqrt2s)$ equal one.
Moreover
\[
 \chi_1u,\ \chi_1v,\ \chi_0u,\ \chi_0v\in L^2(I;H^{1/2}(\R)).
\]
\end{lemma}

\begin{proof}
\uses{cor:NLS-smoothing}
Compactness of the $z$-set is immediate from the compact support of $\psi$ and
the bounded range $\abs s\le s_0$. Choose $\chi_0$ equal to one on it and $\chi_1$
equal to one on a neighbourhood of $\supp\chi_0+[-\sqrt2s_0,\sqrt2s_0]$; the
coincidence property is then the defining support relation. Local smoothing for
the cubic flow, \cref{cor:NLS-smoothing}, gives the stated
$L^2(I;H^{1/2})$ membership after multiplication by the fixed compactly supported
cutoffs.
\end{proof}

\begin{lemma}[Quantitative uniform translated-test multiplier bound]\label{lem:translated-test-multiplier}
\uses{lem:smooth-multiplier-quantitative}
For $\psi\in C_c^\infty(I\times\R)$ and $s_0>0$ set
$\psi_s(t,z)=\psi(t,\sqrt2z+s)$ and $\psi_0(t,z)=\psi(t,\sqrt2z)$. Multiplication
by $\psi_s(t,\cdot)$ maps $H^{1/2}(\R)$ to itself, and there is $C_\psi$ with
\begin{equation}
 \norm{(\psi_s-\psi_0)f}_{H^{1/2}}\le C_\psi\abs s\,\norm f_{H^{1/2}}
 \qquad(0<\abs s\le s_0). \label{eq:translated-test-multiplier}
\end{equation}
\end{lemma}

\begin{proof}
\uses{lem:smooth-multiplier-quantitative}
The quantitative smooth-multiplier estimate \cref{lem:smooth-multiplier-quantitative}
bounds the $H^{1/2}\to H^{1/2}$ operator norm of multiplication by a smooth
function in terms of finitely many of its $C^k$ seminorms. Applying it to
$\psi_s-\psi_0$ and using the mean-value theorem for the first two spatial
derivatives of $\psi$ produces the factor $\abs s$ uniformly for
$0<\abs s\le s_0$, giving \eqref{eq:translated-test-multiplier}.
\end{proof}

\begin{lemma}[Strong $L^2_tH^{-1/2}$ difference-quotient convergence]\label{lem:neumann-strong-convergence}
\uses{lem:translated-test-multiplier, def:Hs-neg-pairing, lem:translation-dq}
With $\chi_1$ as in \cref{lem:neumann-localization}, define for $s\ne0$
\[
 \widetilde D_su=\frac{(\chi_1u)(\,\cdot+\sqrt2s)-\chi_1u}{s},\qquad
 \widetilde D_sv=\frac{(\chi_1v)(\,\cdot+\sqrt2s)-\chi_1v}{s}.
\]
Then, as $s\to0$,
\begin{equation}
 \widetilde D_su\longrightarrow\sqrt2(\chi_1u)_x,\qquad
 \widetilde D_sv\longrightarrow\sqrt2(\chi_1v)_x
 \qquad\text{in }L^2(I;H^{-1/2}). \label{eq:DQ-time-convergence}
\end{equation}
\end{lemma}

\begin{proof}
\uses{lem:translation-dq, lem:neumann-localization}
By \cref{lem:translation-dq}, for almost every $t$ the translation difference
quotients converge, $\widetilde D_su(t)\to\sqrt2(\chi_1u(t))_x$ and
$\widetilde D_sv(t)\to\sqrt2(\chi_1v(t))_x$ in $H^{-1/2}$, with the pointwise
bounds
\[
 \norm{\widetilde D_su(t)}_{H^{-1/2}}\le C\norm{\chi_1u(t)}_{H^{1/2}},\qquad
 \norm{\widetilde D_sv(t)}_{H^{-1/2}}\le C\norm{\chi_1v(t)}_{H^{1/2}},
\]
uniformly for $0<\abs s\le s_0$. Since $\chi_1u,\chi_1v\in L^2(I;H^{1/2})$ by
\cref{lem:neumann-localization}, squaring and applying dominated convergence in
$t$ yields \eqref{eq:DQ-time-convergence}.
\end{proof}

\begin{proposition}[Vanishing Neumann trace]\label{prop:Neumann-trace}
\uses{def:neumann-trace}
For every $\psi\in C_c^\infty(I\times\R)$,
\[
 H_\psi'(0)=0.
\]
Equivalently, $F_s|_{s=0}=0$ as a distribution in $(t,r)$.
\end{proposition}

\begin{proof}
\uses{lem:neumann-localization, lem:neumann-strong-convergence, lem:translated-test-multiplier, cor:diagonal-W}
\emph{Strategy.} Expand $H_\psi(s)/s$ as a difference-quotient pairing, pass to
the limit $s\to0$ using strong convergence of the localized quotients, and
identify the limit with the diagonal Wronskian, which vanishes.

Putting $z=(r-s)/\sqrt2$ so that $r=\sqrt2z+s$, $\dd r=\sqrt2\,\dd z$, the
definition \eqref{eq:Hpsi-def} becomes
\begin{align}
 H_\psi(s)=\sqrt2\iint
 &\big[u(t,z+\sqrt2s)v(t,z)-v(t,z+\sqrt2s)u(t,z)\big]\notag\\
 &\hspace{20mm}\times\psi(t,\sqrt2z+s)\,\dd z\,\dd t, \label{eq:Hpsi-expanded}
\end{align}
and the bracket vanishes at $s=0$.

The nested cutoffs $\chi_0,\chi_1$ of \cref{lem:neumann-localization} localize
$u,v$ to $L^2(I;H^{1/2})$ and make the localized difference quotients
$\widetilde D_su,\widetilde D_sv$ agree exactly with the unlocalized quotients on
the support of every test function used below. Dividing
\eqref{eq:Hpsi-expanded} by $s$ and interpreting products through the
$H^{-1/2}$--$H^{1/2}$ pairing gives the exact identity
\begin{align}
 \frac{H_\psi(s)}s
 =\sqrt2\int_I\Big(
   &\pair{\widetilde D_su(t)}{v(t)\psi_s(t)}
    -\pair{\widetilde D_sv(t)}{u(t)\psi_s(t)}\Big)\,\dd t. \label{eq:Hpsi-dq-exact}
\end{align}
Replacing $\psi_s$ by $\psi_0$ produces an error controlled by
\cref{lem:translated-test-multiplier}, namely at most
$C\abs s\big(\norm{\chi_1u}_{L^2_tH^{1/2}}\norm{\chi_0v}_{L^2_tH^{1/2}}
+\norm{\chi_1v}_{L^2_tH^{1/2}}\norm{\chi_0u}_{L^2_tH^{1/2}}\big)\to0$. By the
strong convergence \eqref{eq:DQ-time-convergence} of
\cref{lem:neumann-strong-convergence} we may pass to the limit in the two
pairings; since $\chi_1$ equals one near $\supp\psi_0$, distributional
differentiation replaces $(\chi_1u)_x,(\chi_1v)_x$ by $u_x,v_x$, giving
\begin{equation}
 H_\psi'(0)=2\int_I\pair{vu_x-uv_x}{\psi(t,\sqrt2\,\cdot)}\,\dd t.
 \label{eq:Hpsi-prime-W}
\end{equation}
The integrand belongs to $L^1(I)$ and vanishes for almost every $t$ by the
diagonal Wronskian \cref{cor:diagonal-W}; hence $H_\psi'(0)=0$.
% FULL PROOF: paper.tex L3087--L3205
\end{proof}

\section{The zero-extension jump formula}

\begin{proposition}[Zero-extension jump formula]\label{prop:zero-extension-jump}
\uses{prop:Dirichlet-trace, prop:Neumann-trace, def:exterior-product}
Let $G(t,r,s)=\1_{\{s>0\}}F(t,r,s)$, so that $G\in C(I;L^2_{r,s})\cap L^4(I\times\R^2)$.
Then, in distributions on $I\times\R^2$,
\begin{equation}
 iG_t+G_{rr}+G_{ss}=QG, \label{eq:G-equation}
\end{equation}
and $G=0$ on the open half-space $s<0$.
\end{proposition}

\begin{proof}
\uses{prop:Dirichlet-trace, prop:Neumann-trace, lem:tensor-test-density}
For $h\in H^2_{\loc}(\R)$ with $h_+=\1_{\{s>0\}}h$, integrating
$\int_0^\infty h\varphi''$ twice by parts against a test $\varphi$ gives the
one-dimensional jump identity
\[
 h_+''=\1_{\{s>0\}}h''+h'(0)\delta_0+h(0)\delta_0'\quad\text{in }\D'(\R),
\]
the boundary terms $h'(0)\varphi(0)-h(0)\varphi'(0)$ being exactly the pairings
with $h'(0)\delta_0+h(0)\delta_0'$.

Testing $G$ first in $(t,r)$ against $\psi$ produces the scalar functions
$H_\psi$ of \eqref{eq:Hpsi-def}. By the jump identity the only possible defect in
$\partial_s^2G$ is $H_\psi'(0)\delta_0+H_\psi(0)\delta_0'$, and both coefficients
vanish by \cref{prop:Dirichlet-trace,prop:Neumann-trace}. The remaining terms are
the restriction of the equation $iF_t+F_{rr}+F_{ss}=QF$ to $s>0$, so
\eqref{eq:G-equation} holds against tests $\psi(t,r)\varphi(s)$. Finite sums of
such tensor tests approximate an arbitrary compactly supported smooth test
together with two derivatives by \cref{lem:tensor-test-density}, and continuity
of all $L^2$ pairings extends the identity to general tests. Vanishing on $s<0$
is the definition of $G$.
% FULL PROOF: paper.tex L3223--L3248
\end{proof}


\begin{proposition}[The zero-extension solves a clean ridge equation]\label{prop:G-equation}
\uses{def:normal-tangential-coords, prop:F-equation, prop:zero-extension-jump, def:exterior-product}
Let $F$ be the exterior product expressed in the normal/tangential coordinates $(r,s)$ of
\cref{def:normal-tangential-coords}, solving the ridge equation of \cref{prop:F-equation}, and let
$G=\1_{\{s>0\}}F$ be its zero extension across the diagonal $s=0$. If the Dirichlet and Neumann
traces of $F$ on $\{s=0\}$ both vanish (\cref{prop:zero-extension-jump}), then on every bounded
open interval $I$
\[
 iG_t+G_{rr}+G_{ss}=QG,\qquad G=0\ \text{for }s<0,
\]
with $G\in C(I;L^2_{r,s})\cap L^4(I\times\R^2)$ and $Q\in L^1(I;L^\infty_{r,s})$.
\end{proposition}

\begin{proof}
\uses{prop:zero-extension-jump, prop:F-equation, lem:Q-bounds}
Differentiating the zero extension produces, a priori, a single- and a double-layer distribution
on $\{s=0\}$ carrying the Neumann and Dirichlet traces of $F$; by \cref{prop:zero-extension-jump}
both vanish, so no boundary distribution survives and $G$ satisfies the same ridge equation as $F$
throughout $I\times\R^2$. The regularity of $G$ and the bound on $Q$ are those of
\cref{prop:F-equation,lem:Q-bounds}.
% FULL PROOF: paper.tex L3207-L3257
\end{proof}

\chapter{Distributional utilities}\label{chap:utilities}

\section{Density of tensor-product tests}

\begin{lemma}[Finite-derivative density of separated tests]
\label{lem:tensor-test-density}
Let $K\Subset\R^{m+n}$ be compact, let $M\in\N$, and let
$\Psi\in C_c^\infty(\R^m\times\R^n)$ have support in $K$. Then there exist
finite sums
\[
 \Psi_N(t,x)=\sum_{j=1}^{J_N}\eta_{N,j}(t)\,\phi_{N,j}(x),
 \qquad
 \eta_{N,j}\in C_c^\infty(\R^m),\quad
 \phi_{N,j}\in C_c^\infty(\R^n),
\]
all supported in one fixed compact rectangle, such that
\begin{equation}
 \max_{\abs{\alpha}\le M}
 \norm{\partial^\alpha(\Psi_N-\Psi)}_\infty\longrightarrow0
 \qquad(N\to\infty).
 \label{eq:tensor-test-density}
\end{equation}
\end{lemma}

\begin{proof}
Strategy: periodize $\Psi$ on a rectangle strictly containing $K$, use decay
of its Fourier coefficients to get $C^M$-convergent rectangular partial sums,
and cut off by a product cutoff so the sum stays a finite combination of
separated smooth functions of $t$ and of $x$.

Choose open product rectangles $K\subset R_0\Subset R_1\Subset R_2$ and a
product cutoff $\kappa(t,x)=\kappa_t(t)\kappa_x(x)$ with $\kappa\equiv1$ on
$R_0$ and $\supp\kappa\subset R_1$. Since $\Psi$ vanishes near
$\partial R_2$, its restriction to $R_2$ extends to a smooth function
$\widetilde\Psi$ periodic on the torus obtained by identifying opposite faces
of $R_2$, and $\widetilde\Psi=\Psi$ on $R_2$. Writing the Fourier series
\[
 \widetilde\Psi(t,x)
 =\sum_{k\in\Z^m,\,\ell\in\Z^n}
 c_{k,\ell}\,\e^{2\pi i k\cdot t/L_t}\e^{2\pi i\ell\cdot x/L_x},
\]
repeated integration by parts $N_0$ times in a coordinate maximizing
$\abs{(k,\ell)}$ gives
\begin{equation}
 \abs{c_{k,\ell}}
 \le C_{N_0}\,(1+\abs{k}+\abs{\ell})^{-N_0}.
 \label{eq:tensor-fourier-decay}
\end{equation}
Fixing $N_0>M+m+n+1$, estimate \eqref{eq:tensor-fourier-decay} makes the
Fourier series absolutely and uniformly convergent after any derivative of
order $\le M$; hence the rectangular partial sums $P_N$ converge to
$\widetilde\Psi$ in the $C^M$ norm on $R_2$. Set $\Psi_N=\kappa P_N$,
extended by zero outside $R_2$: this is smooth because $\supp\kappa\Subset R_2$,
and each Fourier monomial in $P_N$ times $\kappa$ is a product of a function of
$t$ and a function of $x$, so $\Psi_N$ is a finite sum of separated compactly
supported smooth functions, all supported in the fixed rectangle $\supp\kappa$.
The product rule together with $C^M$-convergence of $P_N$ yields
\eqref{eq:tensor-test-density}, using $\kappa\widetilde\Psi=\kappa\Psi=\Psi$
since $\kappa=1$ on $\supp\Psi$.
% FULL PROOF: paper.tex L4313–L4356
\end{proof}

\section{Slicing an $L^1$ family of distributions}

\begin{lemma}[Slicing an $L^1$ family of distributions]
\label{lem:distribution-slicing}
\uses{lem:tensor-test-density}
Fix a compact spatial interval $K$ and a time interval $I$. Let
$t\mapsto\Lambda_t$ be a measurable family of distributions satisfying, for
every $\phi\in C_c^\infty$ supported in $K$,
\[
 \abs{\Lambda_t(\phi)}\le A(t)\bigl(\norm\phi_\infty+\norm{\phi'}_\infty\bigr),
 \qquad A\in L^1(I).
\]
If
\[
 \int_I\eta(t)\,\Lambda_t(\phi)\,\dd t=0
\]
for every $\eta\in C_c^\infty(I)$ and every $\phi\in C_c^\infty$ supported in
$K$, then $\Lambda_t=0$ as a distribution for almost every $t\in I$.
\end{lemma}

\begin{proof}
\uses{lem:tensor-test-density}
Choose a countable family $\{\phi_i\}$ that is dense, in the $C^1$ norm, among
the tests supported in $K$. For each fixed $\phi_i$, the scalar fundamental
lemma of the calculus of variations applied to
$t\mapsto\Lambda_t(\phi_i)\in L^1(I)$ produces a null set $Z_i\subset I$
outside of which $\Lambda_t(\phi_i)=0$. Remove the countable union
$Z=\bigcup_i Z_i$, still null. For $t\notin Z$ the linear functional
$\Lambda_t$ vanishes on the $C^1$-dense family $\{\phi_i\}$, and by the
displayed continuity bound (with $A(t)<\infty$ for a.e.\ $t$) it extends
continuously in the $C^1$ norm; hence $\Lambda_t(\phi)=0$ for every admissible
$\phi$, i.e.\ $\Lambda_t=0$. A countable compact exhaustion of the spatial
variable yields one common full-measure set of good times.
\end{proof}

\section{Distributional products in the exterior equation}

\begin{lemma}[Distributional products in the exterior equation]
\label{lem:distributional-products}
\uses{def:exterior-product}
Let $I$ be a compact time interval and let $u,v$ be $L^2$-solutions of the
cubic flow on $I$ with the endpoint regularity $u\in L^6_{t,x}(I)$ and
$v\in L^\infty_tL^2_x(I)$. Then the product $QF$ arising in the exterior
equation is a genuine $L^2$ function on $I$, not a formal distributional
multiplication. In particular, a typical term satisfies
\[
 \norm{\,\abs{u(t,x)}^2u(t,x)\,v(t,y)\,}_{L^2_{t,x,y}(I)}
 =\norm v_{L^\infty_tL^2_y}
 \left(\int_I\norm{u(t)}_6^6\,\dd t\right)^{1/2}<\infty .
\]
Consequently every passage to the limit in the exterior equation
\eqref{eq:F-equation} may be carried out in $L^2$, which is stronger than
distributional convergence.\end{lemma}

\begin{proof}
\uses{def:exterior-product}
The displayed norm factors exactly by Fubini: integrating first in $y$ gives
$\norm{v(t)}_{L^2_y}^2\le\norm v_{L^\infty_tL^2_y}^2$, and integrating the
remaining $\abs{u(t,x)}^6$ in $x$ and $t$ gives
$\int_I\norm{u(t)}_6^6\,\dd t$; taking the square root yields the stated
identity, finite by the endpoint hypotheses on $u$ and $v$. Since each factor
of $QF$ is thus in $L^2_{t,x,y}(I)$, the products are honest $L^2$ functions
and limits taken in \eqref{eq:F-equation} converge in $L^2$, hence a fortiori
in the sense of distributions.
\end{proof}


\chapter{A self-contained half-space Carleman theorem}\label{chap:carleman}

We prove the linear uniqueness statement needed for the phase-retrieval argument. The
proof has three parts: a uniform exponential-weight estimate obtained from an explicit
inverse of the conjugated free Schr\"odinger operator, a small-potential absorption lemma
converting that estimate into exponential separation, and a graph-closure step extending it
to the non-smooth exterior product.

\section{The solution and source spaces}

\begin{definition}[Source, solution, and multiplier spaces $X,X',Y$]\label{def:carleman-spaces}
\uses{def:scalar-mixed-norm, def:sum-intersection-norm}
On a time interval $J$ and in two spatial dimensions $z=(z_1,z_2)\in\R^2$, define
\begin{align}
 X(J)  &=L^1(J;L^2(\R^2))+L^{4/3}(J\times\R^2),      \label{eq:X-space}\\
 X'(J) &=L^\infty(J;L^2(\R^2))\cap L^4(J\times\R^2), \label{eq:Xprime-space}\\
 Y(J)  &=L^1(J;L^\infty(\R^2))+L^2(J\times\R^2).     \label{eq:Y-space}
\end{align}
Every component is a scalar mixed-norm space in the sense of
\cref{def:scalar-mixed-norm}. The sum spaces carry the infimum norm and the intersection
space the maximum norm, as in \cref{def:sum-intersection-norm}. When $J=\R$ the interval is
omitted from the notation.
\end{definition}

\section{The conjugated free operator}

\begin{definition}[Conjugated free Schr\"odinger operator $P_\beta$]\label{def:conjugated-operator}
\uses{def:propagator}
Let $P=i\partial_t+\Delta_z$ act on functions of $(t,z)\in\R\times\R^2$, so that $P$ is the
free Schr\"odinger operator whose $L^2$ evolution is generated by the propagator of
\cref{def:propagator}. For $\beta\in\R$ define the conjugated operator
\begin{equation}
 P_\beta=\e^{\beta z_1}P\,\e^{-\beta z_1}
        =i\partial_t+\Delta_z-2\beta\partial_{z_1}+\beta^2. \label{eq:Pbeta}
\end{equation}
\end{definition}

\begin{definition}[Explicit inverse $T_\beta$ on smooth compactly supported forcing]\label{def:Tbeta}
\uses{def:conjugated-operator}
For $\beta>0$ set $a(\xi)=\beta^2-|\xi|^2$. Taking the spatial Fourier transform of
$P_\beta w=f$ gives
\begin{equation}
 i\partial_t\widehat w+(\beta^2-|\xi|^2-2i\beta\xi_1)\widehat w=\widehat f. \label{eq:Pbeta-ODE}
\end{equation}
Solving this scalar ODE in the direction of exponential decay of the selected branch --- for
$\xi_1<0$ forward from $-\infty$, for $\xi_1>0$ backward from $+\infty$ --- defines $T_\beta$
by
\begin{equation}
 \widehat{T_\beta f}(t,\xi)=
 \begin{cases}
 -i\displaystyle\int_{-\infty}^t
   \e^{(ia(\xi)+2\beta\xi_1)(t-s)}\widehat f(s,\xi)\,\dd s, &\xi_1<0,\\[1.0em]
 \phantom{-}i\displaystyle\int_t^{\infty}
   \e^{(ia(\xi)+2\beta\xi_1)(t-s)}\widehat f(s,\xi)\,\dd s, &\xi_1>0.
 \end{cases}                                     \label{eq:Tbeta-def}
\end{equation}
In the selected integration direction the exponential has modulus at most one. We write
$K_\beta(\tau)$ for the spatial convolution operator with time increment $\tau=t-s$
appearing in \eqref{eq:Tbeta-def}, so that $T_\beta f(t)=\int_\R K_\beta(t-s)f(s)\,\dd s$.
\end{definition}

\begin{lemma}[Symbol and branch of the conjugated resolvent]\label{lem:Tbeta-symbol}
\uses{def:Tbeta, def:conjugated-operator}
For $\beta>0$ and $\tau\ne0$ the operator $K_\beta(\tau)$ of \cref{def:Tbeta} is the spatial
Fourier multiplier operator with symbol
\begin{equation}
 m_\tau(\xi)=\mp i\,\1_{\{\tau\xi_1<0\}}\,\e^{(ia(\xi)+2\beta\xi_1)\tau},
 \qquad a(\xi)=\beta^2-|\xi|^2,
\end{equation}
supported where $\tau$ and $\xi_1$ have opposite signs. It satisfies
$|m_\tau(\xi)|=\e^{-2\beta|\tau\xi_1|}\le1$, and consequently
\begin{equation}
 \norm{K_\beta(\tau)}_{L^2_z\to L^2_z}\le1. \label{eq:K-L2}
\end{equation}
\end{lemma}

\begin{proof}
\uses{def:Tbeta}
In \eqref{eq:Tbeta-def} the branch for $\xi_1<0$ integrates over $s<t$, i.e.\ $\tau=t-s>0$,
and the branch for $\xi_1>0$ over $\tau<0$; hence the support is $\{\tau\xi_1<0\}$. On that
support the real part of the exponent is $2\beta\xi_1\tau=-2\beta|\tau\xi_1|$, so the
modulus of $m_\tau$ equals $\e^{-2\beta|\tau\xi_1|}\le1$. Plancherel then gives
\eqref{eq:K-L2}.
\end{proof}

\begin{lemma}[Stationary-phase (van der Corput) kernel bound]\label{lem:vandercorput-kernel}
Let $\phi\in C^2([0,R];\R)$ satisfy $|\phi''|\ge\lambda>0$, and let $A$ be absolutely
continuous on $[0,R]$. Then
\begin{equation}
 \left|\int_0^R\e^{i\phi(y)}A(y)\,\dd y\right|
 \le C\lambda^{-1/2}\Big(\norm A_\infty+\int_0^R|A'(y)|\,\dd y\Big), \label{eq:vdc}
\end{equation}
with $C$ absolute. Consequently, for the phase $ay\pm y^2$ (for which $\lambda=2$),
\begin{equation}
 \left|\int_0^\infty\e^{i(ay\pm y^2)}\e^{-by}\,\dd y\right|\le C
 \qquad(a\in\R,\ b\ge0), \label{eq:uniform-Fresnel}
\end{equation}
the improper integral converging and the bound being uniform in $a,b$.
\end{lemma}

\begin{proof}
\emph{Strategy.} Split $[0,R]$ at the set where the phase derivative is small and integrate
by parts on the complement, exploiting monotonicity of $\phi'$.

Since $\phi'$ is monotone, $E=\{|\phi'|\le\sqrt\lambda\}$ is an interval of length at most
$2\lambda^{-1/2}$, contributing at most $2\lambda^{-1/2}\norm A_\infty$. The complement has
at most two components; on each, integration by parts with
$\e^{i\phi}=(i\phi')^{-1}(\e^{i\phi})'$ produces boundary terms, an $A'$ term bounded by
$\lambda^{-1/2}\int|A'|$, and a term controlled by
$\norm A_\infty\operatorname{Var}(1/\phi')\le2\lambda^{-1/2}\norm A_\infty$. Summing gives
\eqref{eq:vdc}. For \eqref{eq:uniform-Fresnel} with $b>0$ take $A=\e^{-by}$, so
$\norm A_\infty+\int_0^\infty|A'|\le2$; for $b=0$ a further integration by parts on
$[R_0,\infty)$ with $R_0>2(1+|a|)$ shows the partial integrals are Cauchy, giving the same
uniform bound.
% FULL PROOF: paper.tex L3373-3412
\end{proof}

\begin{lemma}[Homogeneous $L^4_{t,z}$ bounds for the kernel family]\label{lem:K-homogeneous}
\uses{def:conjugated-operator, def:scalar-mixed-norm}
For every $s\in\R$, $g\in L^2(\R^2)$, and $\beta\ge0$,
\begin{align}
 \norm{K_\beta(t-s)g}_{L^4_{t,z}(\R^{1+2})}   &\le C\norm g_2, \label{eq:K-homogeneous}\\
 \norm{K_\beta(t-s)^*g}_{L^4_{t,z}(\R^{1+2})} &\le C\norm g_2. \label{eq:K-adjoint-homogeneous}
\end{align}
The constants $C$ are absolute, independent of $s$ and of $\beta$.
\end{lemma}

\begin{proof}
\uses{lem:Tbeta-symbol, lem:vandercorput-kernel, lem:riesz-thorin-used, lem:HLS-time}
\emph{Strategy.} A $TT^*$ argument. Write $A_\beta g(t)=K_\beta(t)g$ and reduce to $s=0$ by
time translation. The multiplier of $K_\beta(t)K_\beta(s)^*$ vanishes when $t,s$ have
opposite signs (\cref{lem:Tbeta-symbol}) and otherwise carries oscillation
$\exp(i(\beta^2-|\xi|^2)(t-s))$ and damping $\exp(-2\beta(|t|+|s|)|\xi_1|)$, so its
$L^2_z$-norm is at most one. The factorization-and-half-Fresnel bound of
\cref{lem:vandercorput-kernel}, with oscillatory scale $|t-s|$, gives
$\norm{K_\beta(t)K_\beta(s)^*}_{L^1_z\to L^\infty_z}\le C|t-s|^{-1}$; the compatible-operator
Riesz--Thorin theorem \cref{lem:riesz-thorin-used} interpolates this to
$\norm{K_\beta(t)K_\beta(s)^*}_{L^{4/3}_z\to L^4_z}\le C|t-s|^{-1/2}$. The time-line
Hardy--Littlewood--Sobolev estimate \cref{lem:HLS-time} then yields
$\norm{A_\beta A_\beta^*F}_{L^4_{t,z}}\le C\norm F_{L^{4/3}_{t,z}}$, whence by duality
$\norm{A_\beta^*F}_2^2=|\pair{A_\beta A_\beta^*F}{F}|\le C\norm F_{4/3}^2$ and, dualizing
once more, $A_\beta:L^2\to L^4$, i.e.\ \eqref{eq:K-homogeneous}. Replacing each kernel by its
adjoint gives \eqref{eq:K-adjoint-homogeneous}.
% FULL PROOF: paper.tex L3447-3491
\end{proof}

\begin{lemma}[Kernel estimates for $T_\beta$ (assembled)]\label{lem:Tbeta-kernel}
\uses{def:Tbeta}
Let $K_\beta(\tau)$ be the spatial convolution operator of \cref{def:Tbeta}. For every
$\tau\ne0$,
\begin{align}
 \norm{K_\beta(\tau)}_{L^2_z\to L^2_z}      &\le1,          \label{eq:K-L2b}\\
 \norm{K_\beta(\tau)}_{L^1_z\to L^\infty_z} &\le C|\tau|^{-1},   \label{eq:K-disp}\\
 \norm{K_\beta(\tau)}_{L^{4/3}_z\to L^4_z}  &\le C|\tau|^{-1/2}. \label{eq:K-43-4}
\end{align}
The constants are absolute, independent of $\beta>0$.
\end{lemma}

\begin{proof}
\uses{lem:Tbeta-symbol, lem:vandercorput-kernel, lem:complex-gaussian-kernel, lem:riesz-thorin-used}
\emph{Strategy.} Estimate \eqref{eq:K-L2b} is \cref{lem:Tbeta-symbol}. For \eqref{eq:K-disp},
factor the inverse Fourier integral into the $\xi_2$ and $\xi_1$ variables: the full Gaussian
integral in $\xi_2$ has modulus $C|\tau|^{-1/2}$ by \cref{lem:complex-gaussian-kernel}
(boundary value from positive Gaussian damping), and after scaling $\eta=|\tau|^{-1/2}y$ the
$\xi_1$ factor is the half-line Fresnel integral \eqref{eq:uniform-Fresnel} of
\cref{lem:vandercorput-kernel}, of size $O(|\tau|^{-1/2})$; the two combine to
$C|\tau|^{-1}$. The oscillatory kernel is realized rigorously by inserting a Gaussian
multiplier cutoff $\e^{-\delta|\xi|^2}$, applying the same bounds uniformly in $\delta$, and
passing to the limit on Schwartz functions and then by $L^1$ density. Finally the
$L^1\to L^\infty$ kernel realization and the $L^2$ multiplier realization agree on Schwartz
functions, so the compatible-operator Riesz--Thorin theorem \cref{lem:riesz-thorin-used}
gives \eqref{eq:K-43-4}.
% FULL PROOF: paper.tex L3352-3433
\end{proof}

\section{The uniform linear-weight estimate}

\begin{lemma}[Critical multiplier estimate]\label{lem:XY-multiplier}
\uses{def:carleman-spaces}
There is an absolute $C$ such that, for every time interval $J$,
\begin{equation}
 \norm{WZ}_{X(J)}\le C\norm W_{Y(J)}\norm Z_{X'(J)}. \label{eq:XY-multiplier}
\end{equation}
\end{lemma}

\begin{proof}
\uses{def:carleman-spaces}
Write $W=W_1+W_2$ with $W_1\in L^1_tL^\infty_z$ and $W_2\in L^2_{t,z}$. Then
\[
 \norm{W_1Z}_{L^1_tL^2_z}\le\norm{W_1}_{L^1_tL^\infty_z}\norm Z_{L^\infty_tL^2_z},
 \qquad
 \norm{W_2Z}_{L^{4/3}_{t,z}}\le\norm{W_2}_{L^2_{t,z}}\norm Z_{L^4_{t,z}}.
\]
So $W_1Z\in L^1_tL^2_z$ and $W_2Z\in L^{4/3}_{t,z}$ give a decomposition of $WZ$ realizing
the $X$-norm. Taking the infimum over decompositions of $W$ yields \eqref{eq:XY-multiplier}.
\end{proof}

\begin{proposition}[Uniform linear-weight Carleman estimate]\label{prop:linear-carleman}
\uses{def:conjugated-operator, def:carleman-spaces}
For every $\beta\in\R$ and every $\varphi\in C_c^\infty(\R^{1+2})$,
\begin{equation}
 \norm{\e^{\beta z_1}\varphi}_{X'}\le C\norm{\e^{\beta z_1}P\varphi}_{X}, \label{eq:linear-carleman}
\end{equation}
where $C$ is absolute and independent of $\beta$. Equivalently, $T_\beta:X\to X'$ boundedly,
uniformly in $\beta$.
\end{proposition}

\begin{proof}
\uses{lem:Tbeta-kernel, lem:K-homogeneous, lem:HLS-time, def:Tbeta, lem:sum-intersection-complete}
\emph{Strategy.} By reflection take $\beta>0$. Set $w=\e^{\beta z_1}\varphi$ and
$f=\e^{\beta z_1}P\varphi=P_\beta w$; since $w$ has compact time support, solving
\eqref{eq:Pbeta-ODE} in the selected direction gives $w=T_\beta f$. We verify the four
mappings underlying the sum-to-intersection estimate for the decomposition $f=f_1+f_2$ with
$f_1\in L^1_tL^2_z$, $f_2\in L^{4/3}_{t,z}$ (\cref{def:carleman-spaces},
\cref{lem:sum-intersection-complete}).

\emph{Component $f_1$.} By \eqref{eq:K-L2b} and Minkowski,
$\norm{T_\beta f_1}_{L^\infty_tL^2_z}\le\norm{f_1}_{L^1_tL^2_z}$; and by
\cref{lem:K-homogeneous} with Minkowski,
$\norm{T_\beta f_1}_{L^4_{t,z}}\le\int_\R\norm{K_\beta(t-s)f_1(s)}_{L^4_{t,z}}\,\dd s
 \le C\norm{f_1}_{L^1_tL^2_z}$.

\emph{Component $f_2$.} By \eqref{eq:K-43-4} and the time-line HLS estimate
\cref{lem:HLS-time}, $\norm{T_\beta f_2}_{L^4_{t,z}}\le C\norm{f_2}_{L^{4/3}_{t,z}}$. For the
energy bound, fix $t$ and $g\in L^2_z$; spacetime H\"older and the adjoint estimate
\eqref{eq:K-adjoint-homogeneous} give
$|\pair{T_\beta f_2(t)}{g}|\le\norm{f_2}_{L^{4/3}_{s,z}}\norm{K_\beta(t-s)^*g}_{L^4_{s,z}}
 \le C\norm{f_2}_{4/3}\norm g_2$, so $\norm{T_\beta f_2}_{L^\infty_tL^2_z}\le C\norm{f_2}_{4/3}$.

Both components of the $X'$ norm are thus bounded by
$C(\norm{f_1}_{L^1L^2}+\norm{f_2}_{4/3})$; taking the infimum over decompositions proves
\eqref{eq:linear-carleman}.
% FULL PROOF: paper.tex L3503-3547
\end{proof}

\section{Extension to the graph class}

\begin{lemma}[Spatial mollification in the source spaces]\label{lem:spatial-mollification-source}
\uses{def:carleman-spaces}
Let $\rho_\eps$ be a standard approximate identity on $\R^2$ and let $J_0\Subset\R$ be
bounded.
\begin{enumerate}
\item If $f_1\in L^1_tL^2_z$ is supported in $J_0$ in time, then
      $\rho_\eps*_zf_1\to f_1$ in $L^1_tL^2_z$.
\item If $f_2\in L^{4/3}_{t,z}$ is supported in $J_0$ in time, then
      $\rho_\eps*_zf_2\to f_2$ in $L^{4/3}_{t,z}$ and, for each fixed $\eps>0$,
      $\rho_\eps*_zf_2\in L^1_tL^2_z$.
\end{enumerate}
Consequently spatial mollification converges strongly in $X$.
\end{lemma}

\begin{proof}
The two convergence statements are the approximate-identity theorem applied for a.e.\ time
and integrated (equivalently, approximate first by compactly supported simple functions).
For the extra mapping, Young's inequality in $z$ gives
$\norm{\rho_\eps*f_2(t)}_2\le\norm{\rho_\eps}_{4/3}\norm{f_2(t)}_{4/3}$, and since $J_0$ has
finite length $L^{4/3}(J_0)\hookrightarrow L^1(J_0)$, so $\rho_\eps*_zf_2\in L^1_tL^2_z$.
\end{proof}

\begin{lemma}[Scalar compact-support ODE]\label{lem:compact-support-ODE}
Let $\lambda\in\C$, let $y\in L^2(\R)$ have compact support, and let $g\in L^1(\R)$. If
$y'-\lambda y=g$ in distributions, then $y$ has a locally absolutely continuous
representative. Moreover, if $\Re\lambda<0$ then
\[
 y(t)=\int_{-\infty}^t\e^{\lambda(t-s)}g(s)\,\dd s,
\]
while if $\Re\lambda>0$ then
\[
 y(t)=-\int_t^\infty\e^{\lambda(t-s)}g(s)\,\dd s.
\]
\end{lemma}

\begin{proof}
On every bounded interval $\lambda y+g\in L^1$, so $y\in W^{1,1}_{\loc}$ and hence has an
absolutely continuous representative. Multiplying by $\e^{-\lambda t}$ and applying the
fundamental theorem of calculus gives
$\e^{-\lambda t}y(t)-\e^{-\lambda a}y(a)=\int_a^t\e^{-\lambda s}g(s)\,\dd s$. If
$\Re\lambda<0$, letting $a\to-\infty$ past the compact support of $y$ makes $y(a)=0$ and
yields the first formula; if $\Re\lambda>0$, integrating from $t$ to a point right of the
support yields the second.
\end{proof}

\begin{lemma}[Slicing a first-order distributional parameter family]\label{lem:parameterwise-distribution-ODE}
\uses{def:conjugated-operator}
Let $\lambda:\R^d\to\C$ be measurable, let $g\in L^2(\R_t\times\R^d_\xi)$, and let
$b(\cdot,\xi)\in L^1_{\loc}(\R_t)$ for a.e.\ $\xi$. If
\[
 \partial_t g-\lambda(\xi)\,g=b\quad\text{in }\D'(\R_t\times\R^d_\xi),
\]
then there is a full-measure set $\Xi\subset\R^d$ such that for every $\xi\in\Xi$ the section
$g(\cdot,\xi)$ satisfies the scalar distributional ODE
$\partial_t g(\cdot,\xi)-\lambda(\xi)\,g(\cdot,\xi)=b(\cdot,\xi)$ in $\D'(\R_t)$.
\end{lemma}

\begin{proof}
\uses{}
Test against $\eta(t)\phi(\xi)$ with $\eta\in C_c^\infty(\R_t)$, $\phi\in C_c^\infty(\R^d_\xi)$;
Fubini rewrites the pairing as $\int\phi(\xi)\,T_\eta(\xi)\,\dd\xi=0$, where $T_\eta(\xi)$ is the
scalar pairing of the $\xi$-section against $\eta$ (measurable in $\xi$ by Tonelli). Ranging
$\eta$ over a countable dense family of test functions forces $T_\eta(\xi)=0$ for every member of
that family off a $\xi$-null set; the countable intersection of the good sets is the required
full-measure $\Xi$, on which the section equation holds.
% FULL PROOF: paper.tex L3581-L3643
\end{proof}

\begin{lemma}[Fourier-section identification of the mollified inverse]\label{lem:fourier-section-identification}
\uses{def:conjugated-operator, def:carleman-spaces, def:Tbeta}
Let $w\in L^\infty_tL^2_z$ have compact support in a bounded time interval, and suppose
$P_\beta w=f$ in distributions with $f\in L^1_tL^2_z$. Then
\begin{equation}
 w=T_\beta f\quad\text{in distributions and in }X'. \label{eq:section-Tbeta}
\end{equation}
\end{lemma}

\begin{proof}
\uses{lem:parameterwise-distribution-ODE, lem:compact-support-ODE, def:Tbeta, prop:linear-carleman}
\emph{Strategy.} Take the spatial Fourier transform, slice the resulting first-order
parameter-family ODE to a full-measure set of frequencies, solve each scalar ODE by
\cref{lem:compact-support-ODE}, and recognize the formula as \eqref{eq:Tbeta-def}.

Bounded time support gives $w\in L^2_{t,z}$; spatial Plancherel and Fubini give a
full-measure frequency set on which $\widehat w(\cdot,\xi)\in L^2_t$, and Minkowski gives
$\widehat f(\cdot,\xi)\in L^1_t$ for a.e.\ $\xi$. The spatial Fourier transform of the
equation is
\[
 \partial_t\widehat w-(ia(\xi)+2\beta\xi_1)\widehat w=-i\widehat f
 \quad\text{in }\D'(\R_t\times\R^2_\xi).
\]
By \cref{lem:parameterwise-distribution-ODE}, slicing this to a single full-measure set of frequencies reduces it, for a.e.\ $\xi$, to the
scalar distributional ODE with $\lambda=ia(\xi)+2\beta\xi_1$, whose real part is negative for
$\xi_1<0$ and positive for $\xi_1>0$. \cref{lem:compact-support-ODE} then gives exactly the
formulas \eqref{eq:Tbeta-def}; the null hyperplane $\xi_1=0$ is irrelevant. Plancherel
identifies $w=T_\beta f$ in $L^2_{\loc}$ and hence in distributions, and
\cref{prop:linear-carleman} places $T_\beta f\in X'$, giving \eqref{eq:section-Tbeta} as an
equality of $X'$ classes.
% FULL PROOF: paper.tex L3689-3720
\end{proof}

\begin{lemma}[Graph closure of the Carleman estimate]\label{lem:carleman-closure}
\uses{def:conjugated-operator, def:carleman-spaces}
Suppose $U\in X'$ has compact support in time, $U(t,z)=0$ for $z_1>H$ for some $H$, and, for
a fixed $\beta>0$, the weighted distribution $\e^{\beta z_1}PU$ belongs to $X$. Then
\begin{equation}
 \norm{\e^{\beta z_1}U}_{X'}\le C\norm{\e^{\beta z_1}PU}_{X}, \label{eq:graph-carleman}
\end{equation}
with $C$ absolute and independent of $\beta$.
\end{lemma}

\begin{proof}
\uses{lem:spatial-mollification-source, lem:fourier-section-identification, prop:linear-carleman}
\emph{Strategy.} Set $w=\e^{\beta z_1}U$ and $f=\e^{\beta z_1}PU$; the upper support bound
$z_1\le H$ keeps the weight bounded on $\operatorname{supp}U$, so $w\in X'$ with the same
compact time support, and distributional conjugation gives $P_\beta w=f$. Mollify in space
and pass to the limit through the explicit inverse.

Fix a decomposition $f=f_1+f_2$ with $f_1\in L^1_tL^2_z$, $f_2\in L^{4/3}_{t,z}$; replacing
it by $\1_{J_0}f_1+\1_{J_0}f_2$ for a bounded interval $J_0$ containing the time support
keeps the sum equal to $f$ and does not increase either component norm, so both may be taken
supported in $J_0$. Convolving in $z$ with $\rho_\eps$ commutes with the constant-coefficient
$P_\beta$, giving $P_\beta w_\eps=f_\eps$; by \cref{lem:spatial-mollification-source},
$f_\eps\in L^1_tL^2_z$ for each $\eps>0$ and $f_\eps\to f$ in $X$, while
$w_\eps\in L^\infty_tL^2_z$ has compact time support. \cref{lem:fourier-section-identification}
then gives $w_\eps=T_\beta f_\eps$. The bounded map $T_\beta:X\to X'$ of
\cref{prop:linear-carleman} gives $T_\beta f_\eps\to T_\beta f$ in $X'$, while $w_\eps\to w$
in distributions; uniqueness of the distributional limit gives $w=T_\beta f$. Applying the
operator bound and taking the infimum over decompositions of $f$ proves
\eqref{eq:graph-carleman}.
% FULL PROOF: paper.tex L3732-3779
\end{proof}

\section{Exponential separation with a small potential}

\begin{lemma}[One-slab exponential separation]\label{lem:exponential-separation}
\uses{def:carleman-spaces, def:conjugated-operator}
There is an absolute $\eps_*>0$ with the following property. Suppose $U\in X'$ is compactly
supported in time, $U=0$ for $z_1>H$, and
\begin{equation}
 PU=VU+R \label{eq:one-slab-eq}
\end{equation}
in distributions, with $V\in Y$ and $\norm V_Y\le\eps_*$. Assume that for all sufficiently
large $\beta$ the weighted distribution $\e^{\beta z_1}R$ belongs to $X$ and
\begin{equation}
 \norm{\e^{\beta z_1}R}_X\longrightarrow0\quad(\beta\to+\infty). \label{eq:R-decay}
\end{equation}
This holds, for instance, when $R\in X$ is supported in $\{z_1\le-d\}$ for some $d>0$. Then
$U=0$ on $\{z_1>0\}$.
\end{lemma}

\begin{proof}
\uses{lem:carleman-closure, lem:XY-multiplier}
Since $PU=VU+R$, the graph closure estimate \cref{lem:carleman-closure} together with the
multiplier estimate \cref{lem:XY-multiplier} give
\begin{align*}
 \norm{\e^{\beta z_1}U}_{X'}
 &\le C\norm{V\e^{\beta z_1}U}_{X}+C\norm{\e^{\beta z_1}R}_{X}\\
 &\le C\norm V_Y\norm{\e^{\beta z_1}U}_{X'}+C\norm{\e^{\beta z_1}R}_{X}.
\end{align*}
Choose $\eps_*=(2C)^{-1}$, so that $C\norm V_Y\le\tfrac12$ absorbs the first term. On
$\{z_1>0\}$ the weight is at least one, hence
$\norm U_{X'(\{z_1>0\})}\le2C\norm{\e^{\beta z_1}R}_{X}$. Letting $\beta\to\infty$ and using
\eqref{eq:R-decay} gives $U=0$ on $\{z_1>0\}$.
\end{proof}


\chapter{Half-space propagation}\label{chap:halfspace}

We assemble the quantitative moving-interface machinery: a smallness scale $(\delta,h,N)$, a one-slab absorption estimate, a single interface-propagation step, and the resulting half-space uniqueness theorem.

\begin{definition}[Moving-interface smallness scale]\label{def:moving-interface-scale}
\uses{def:carleman-spaces}
Let $c_0>0$ be the absolute constant of \cref{lem:one-step-propagation} and let $J=(a,b)$ be a bounded interval of length $L=b-a$. For an integer $N\ge 1$ the \emph{moving-interface scale} is the triple $(\delta,h,N)$ with
\[
 \delta=\frac{L}{32N},
 \qquad
 h=\frac{c_0}{2}\sqrt\delta,
\]
so that $0<\delta<L/8$ and $h\le c_0\sqrt\delta$. Associated to a moving boundary $w(t)=B+h\,\omega(t)$ with $-1\le\omega\le2$, $|\omega'|\le C\delta^{-1}$, $|\omega''|\le C\delta^{-2}$ and $\omega'',\omega'$ supported in transition intervals of total length $O(\delta)$, one has the derivative bounds
\[
 |w'|\le C\,h\delta^{-1},
 \qquad
 |w''|\le C\,h\delta^{-2},
\]
the near-affine smallness
\[
 \int_J\bigl|w''(t)\,y_1\bigr|\,\1_{\{|y_1|\le h\}}\,\dd t
 \;\le\; C\,h^2\delta^{-1}
 \;=\; C\Bigl(\tfrac{c_0}2\Bigr)^{\!2},
\]
and, after $N$ successive steps that each cost at most $8\delta$ of time, the total propagated distance
\[
 Nh=\frac{c_0}{2}\sqrt{\frac{NL}{32}}\xrightarrow[N\to\infty]{}\infty
\]
while the central time interval retains length at least $3L/4$.
\end{definition}

\begin{lemma}[Moving translation and Galilean gauge in the rough class]\label{lem:moving-galilean-distribution}
\uses{def:conjugated-operator}
Let $P=i\partial_t+\Delta_z$ on $\R_t\times\R^2_z$, let $w\in C^\infty(J)$, and set
$a(t)=w'(t)/2$. If $Z\in C(J;L^2)\cap L^4_{\loc}(J\times\R^2)$ solves $PZ=WZ$, then, writing
$V(t,y)=Z(t,y_1+w(t),y_2)$ and
\[
 U(t,y)=\exp\!\Bigl(-ia(t)y_1-i\!\int_a^t a(s)^2\,\dd s\Bigr)V(t,y),
\]
one has
\[
 PU=\Bigl(W(t,y_1+w(t),y_2)+\tfrac{w''(t)}2\,y_1\Bigr)U
\]
in the same distributional class, with the $L^\infty_tL^2_y$ and $L^4_{t,y}$ norms of $U$ equal
to those of $Z$ and the support of $U$ obtained from that of $Z$ by the shift $y_1\mapsto y_1-w(t)$.
\end{lemma}

\begin{proof}
\uses{}
The translation $y_1\mapsto y_1+w(t)$ turns $\partial_t$ into $\partial_t-w'\partial_{y_1}$; the
Galilean phase $\exp(-ia y_1-i\int a^2)$ absorbs the resulting first-order term and produces the
linear potential $\tfrac{w''}2 y_1$, exactly as in the classical Galilean boost. Every step is
multiplication by a smooth function of $(t,y)$ or a smooth change of variables, so the identity is
valid in $\D'$; the norm and support statements are immediate.
% FULL PROOF: paper.tex L3838-L3866
\end{proof}

\begin{lemma}[Exponential suppression of the weighted error terms]\label{lem:moment-error}
\uses{def:carleman-spaces}
Let $A\in L^1(J)$ and $U\in L^\infty_tL^2_y$, and set
$R=\eta A_{\mathrm{far}}U+i\eta'U$, where $A_{\mathrm{far}}=\tfrac{w''}2 y_1\1_{\{y_1\le-h\}}$ is
supported in $\{y_1\le-h\}$ and $\eta'U$ is supported in $\{y_1\le-2h\}$. Then for every
$\beta>0$ one has $\e^{\beta y_1}R\in L^1_tL^2_y$ and
$\norm{\e^{\beta y_1}R}_{L^1_tL^2_y}\to0$ as $\beta\to\infty$.
\end{lemma}

\begin{proof}
\uses{}
On $\{y_1\le-h\}$, $\e^{\beta y_1}\le\e^{-\beta h}$, and on $\supp\eta'\subset\{y_1\le-2h\}$,
$\e^{\beta y_1}\le\e^{-2\beta h}$. Multiplying by $A\in L^1_t$ and $U\in L^\infty_tL^2_y$ and
integrating gives $\norm{\e^{\beta y_1}R}_{L^1_tL^2_y}\lesssim\e^{-\beta h}+\e^{-2\beta h}\to0$.
% FULL PROOF: paper.tex L3819-L3833
\end{proof}

\begin{lemma}[One-slab separation]\label{lem:one-slab}
\uses{def:carleman-spaces}
There is an absolute constant $\eps_*>0$ with the following property. Suppose $U\in X'$ is compactly supported in time, $U=0$ for $z_1>H$, and
\[
 PU=VU+R
 \qquad\text{in distributions,}
\]
with $V\in Y$ and $\norm{V}_Y\le\eps_*$. Assume that for all sufficiently large $\beta$ the weighted distribution $\e^{\beta z_1}R$ belongs to $X$ and
\[
 \norm{\e^{\beta z_1}R}_X\longrightarrow0
 \qquad(\beta\to+\infty).
\]
This holds, for example, when $R\in X$ is supported in $\{z_1\le-d\}$ for some $d>0$. Then $U=0$ on $\{z_1>0\}$.
\end{lemma}

\begin{proof}
\uses{lem:carleman-closure, lem:XY-multiplier, lem:exponential-separation}
Apply the graph-closure Carleman estimate \cref{lem:carleman-closure} to $\e^{\beta z_1}U$ and the multiplier estimate \cref{lem:XY-multiplier} to the potential term:
\[
 \norm{\e^{\beta z_1}U}_{X'}
 \le C\norm{V\,\e^{\beta z_1}U}_{X}+C\norm{\e^{\beta z_1}R}_{X}
 \le C\norm{V}_Y\norm{\e^{\beta z_1}U}_{X'}+C\norm{\e^{\beta z_1}R}_{X}.
\]
Choose $\eps_*=(2C)^{-1}$; then $C\norm{V}_Y\le\tfrac12$ and the first term is absorbed into the left side, giving $\norm{\e^{\beta z_1}U}_{X'}\le 2C\norm{\e^{\beta z_1}R}_{X}$. On $\{z_1>0\}$ the weight $\e^{\beta z_1}\ge1$, hence
\[
 \norm{U}_{X'(\{z_1>0\})}\le 2C\norm{\e^{\beta z_1}R}_{X}.
\]
Letting $\beta\to\infty$ makes the right side vanish, so $U=0$ on $\{z_1>0\}$.
\end{proof}

\begin{lemma}[Quantitative one-step propagation]\label{lem:one-step-propagation}
\uses{lem:one-slab, def:moving-interface-scale}
There is an absolute constant $c_0>0$ with the following property. Let $J=(a,b)$ have length $L$, let $0<\delta<L/8$, and let $0<h\le c_0\sqrt\delta$. Suppose
\[
 Z\in C(J;L^2)\cap L^4(J\times\R^2),
 \qquad PZ=WZ,
 \qquad \norm{W}_{L^1_tL^\infty_z(J)}\le c_0,
\]
and $Z=0$ on $J\times\{z_1>B\}$. Then
\[
 Z=0\quad\text{on }(a+4\delta,b-4\delta)\times\{z_1>B-h\}.
\]
\end{lemma}

\begin{proof}
\uses{lem:one-slab, lem:XY-multiplier, lem:moving-galilean-distribution, lem:moment-error, def:moving-interface-scale}
Strategy: translate the interface along a smooth path $w(t)=B+h\omega(t)$ chosen from the smallness scale of \cref{def:moving-interface-scale}, apply a Galilean gauge to reduce to the class of \cref{lem:one-slab}, verify the potential smallness through \cref{lem:XY-multiplier}, and check the exponential decay of the two error terms so that \cref{lem:one-slab} propagates the vanishing across a single slab.

\textbf{1. Compatible cutoffs.} Pick $\omega\in C^\infty([a,b])$ equal to $2$ on $[a,a+2\delta]\cup[b-2\delta,b]$, equal to $-1$ on $[a+3\delta,b-3\delta]$, with $-1\le\omega\le2$, $|\omega'|\le C\delta^{-1}$, $|\omega''|\le C\delta^{-2}$, and $\eta\in C_c^\infty(J)$ equal to $1$ on $[a+2\delta,b-2\delta]$, vanishing near the endpoints, $|\eta'|\le C\delta^{-1}$, with $\supp\eta'$ inside $\{\omega=2\}$. Set $w(t)=B+h\omega(t)$ and $V(t,y)=Z(t,y_1+w(t),y_2)$; the half-space vanishing gives $V=0$ for $y_1>-h\omega(t)$, hence $V=0$ for $y_1>h$ everywhere and $V=0$ for $y_1>-2h$ on $\supp\eta'$.

\textbf{2. Galilean gauge.} With $a(t)=w'(t)/2$ set $U(t,y)=\exp(-ia(t)y_1-i\int_a^t a(s)^2\,\dd s)V(t,y)$. The moving-translation and Galilean-gauge identity converts $PZ=WZ$ into
\[
 PU=\Bigl(W(t,y_1+w(t),y_2)+\tfrac{w''(t)}2 y_1\Bigr)U
\]
in the same rough distributional class, and preserves the $L^\infty_tL^2_y$ and $L^4_{t,y}$ norms and the support inequalities.
\textbf{3. Bounded near coefficient, weighted far error.} Put $Y=\eta U\in X'$, compactly time-supported, vanishing for $y_1>h$. Split $\tfrac{w''}2 y_1=A_{\mathrm{near}}+A_{\mathrm{far}}$ with $A_{\mathrm{near}}=\tfrac{w''}2 y_1\1_{\{-h<y_1\le h\}}$ and $A_{\mathrm{far}}=\tfrac{w''}2 y_1\1_{\{y_1\le-h\}}$. Since $w''=h\omega''$ has time support $O(\delta)$,
\[
 \norm{A_{\mathrm{near}}}_{L^1_tL^\infty_y}\le\frac h2\int_J|w''|\,\dd t\le C h^2\delta^{-1},
\]
which is the near-affine smallness of \cref{def:moving-interface-scale}. The equation for $Y$ is $PY=\widetilde V Y+R$ with $\widetilde V=W(t,y_1+w(t),y_2)+A_{\mathrm{near}}$ and $R=\eta A_{\mathrm{far}}U+i\eta'U$. Choosing $c_0$ so small that $c_0+Cc_0^2\le\eps_*$, the hypotheses $\norm{W}_{L^1L^\infty}\le c_0$ and $h^2/\delta\le c_0^2$ give $\norm{\widetilde V}_Y\le\eps_*$.

\textbf{4. Exponential decay of the errors.} For every $\beta>0$ the weighted far term $\e^{\beta y_1}\eta A_{\mathrm{far}}U$ lies in $L^1_tL^2_y$ and tends to $0$ as $\beta\to\infty$, since $A(t)=\eta(t)w''(t)/2\in L^1_t$, $U\in L^\infty_tL^2_y$, and it is supported in $\{y_1\le-h\}$. On $\supp\eta'$ one has $U=0$ for $y_1>-2h$, so $\norm{\e^{\beta y_1}\eta'U}_{L^1_tL^2_y}\le\e^{-2\beta h}\norm{\eta'}_{L^1_t}\norm{U}_{L^\infty_tL^2_y}\to0$. Thus $R$ satisfies the decay hypothesis of \cref{lem:one-slab}.
\textbf{5. One slab.} Apply \cref{lem:one-slab} with upper support bound $H=h$ to obtain $Y=0$ for $y_1>0$. On $(a+4\delta,b-4\delta)$ we have $\eta=1$ and $w=B-h$, so $\{y_1>0\}$ is exactly $\{z_1>B-h\}$, giving the conclusion.

% FULL PROOF: paper.tex L3920-L4033
\end{proof}

\begin{theorem}[Half-space uniqueness for integrably bounded potentials]\label{thm:halfspace-UC}
\uses{def:scalar-mixed-norm}
Let $J$ be an open interval and let
\[
 Z\in C(J;L^2(\R^2))\cap L^4_{\loc}(J\times\R^2),
 \qquad PZ=WZ,
\]
where $W\in L^1_{\loc}(J;L^\infty(\R^2))$. If for some $B\in\R$
\[
 Z=0\quad\text{almost everywhere on }J\times\{z_1>B\},
\]
then $Z=0$ on $J\times\R^2$.
\end{theorem}

\begin{proof}
\uses{lem:one-step-propagation, def:moving-interface-scale}
Strategy: fix a base time and iterate \cref{lem:one-step-propagation} along the moving-interface scale of \cref{def:moving-interface-scale}, whose per-step time loss is controlled while the propagated distance diverges.

Fix $t_0\in J$ and a target distance $D>0$. By absolute continuity of the $L^1_tL^\infty_z$ norm, choose a symmetric interval $I=(t_0-L/2,t_0+L/2)\Subset J$ with $\norm{W}_{L^1_tL^\infty_z(I)}\le c_0/2$. Take an integer $N$ and set $\delta=L/(32N)$, $h=(c_0/2)\sqrt\delta$ as in \cref{def:moving-interface-scale}; the constraint $h\le c_0\sqrt\delta$ holds. Each of the $N$ applications of \cref{lem:one-step-propagation} loses at most $8\delta$ of time, so the surviving central interval keeps length at least $3L/4$ and contains $t_0$; the potential norm on every subinterval is at most that on $I$. After $N$ steps $Z$ vanishes a.e. on a time neighborhood of $t_0$ and on $\{z_1>B-Nh\}$, and continuity into $L^2$ upgrades this to the exact time $t_0$. Since $Nh=(c_0/2)\sqrt{NL/32}\to\infty$, choose $N$ with $Nh>D$; then $Z(t_0,z)=0$ for $z_1>B-D$. As $D$ is arbitrary, $Z(t_0)=0$ in $L^2(\R^2)$, and as $t_0$ is arbitrary, $Z=0$ on $J\times\R^2$.

% FULL PROOF: paper.tex L4051-L4079
\end{proof}

\begin{corollary}[Half-space uniqueness in the $X,Y$ formulation]\label{cor:halfspace-XY}
\uses{def:carleman-spaces}
The conclusion of \cref{thm:halfspace-UC} remains valid if $W\in Y_{\loc}$ and, around every time, there is a sufficiently short interval on which $W$ admits a decomposition with $Y$ norm smaller than $c_0/2$. The case $W\in L^1_{\loc}L^\infty$ used in the main proof is the first summand of this formulation.
\end{corollary}

\begin{proof}
\uses{thm:halfspace-UC, lem:XY-multiplier, lem:one-slab}
The proof of \cref{lem:one-step-propagation} uses the potential only through the multiplier estimate \cref{lem:XY-multiplier} and the smallness required by \cref{lem:one-slab}. Replacing the $L^1_tL^\infty_z$ norm throughout by the corresponding small $Y$ norm leaves every step, and hence the iteration of \cref{thm:halfspace-UC}, unchanged.
\end{proof}


\chapter{Completion of the phase-retrieval proof}\label{chap:completion}

\section{Vanishing of the exterior product}

\begin{lemma}[Vanishing of the exterior product via reflection and half-space UC]
\label{lem:F-vanishing}
\uses{def:exterior-product, def:normal-tangential-coords, prop:G-equation, prop:F-equation, lem:Q-bounds}
Let $u,v\in C(\R;L^2(\R))$ be the global mild solutions of the cubic NLS with
equal densities, let
\[
 F(t,x,y)=u(t,x)v(t,y)-v(t,x)u(t,y)
\]
be their exterior product, and let $G(t,r,s)$ be $F$ expressed in the
normal/tangential coordinates $(r,s)$ of \cref{def:normal-tangential-coords},
so that on any bounded open interval $I$
\[
 iG_t+G_{rr}+G_{ss}=QG,\qquad G=0\ \text{for }s<0,
\]
with $G\in C(I;L^2_{r,s})\cap L^4(I\times\R^2)$ and, by \cref{lem:Q-bounds},
$Q\in L^1(I;L^\infty_{r,s})$.  Then
\[
 F(t)=0\quad\text{in }L^2(\R^2)\quad\text{for every }t\in\R.
\]
\end{lemma}

\begin{proof}
\uses{prop:G-equation, prop:zero-extension-jump, thm:halfspace-UC, lem:distributional-products}
Fix a bounded open interval $I$.  Reflect the normal variable by
\[
 z_1=-s,\qquad z_2=r,\qquad Z(t,z_1,z_2)=G(t,z_2,-z_1),
\]
and set $W(t,z_1,z_2)=Q(t,z_2,-z_1)$.  Then
\[
 (i\partial_t+\Delta_z)Z=WZ,\qquad Z=0\ \text{for }z_1>0 ,
\]
with $Z\in C(I;L^2)\cap L^4(I\times\R^2)$ and $W\in L^1(I;L^\infty)$.  All
hypotheses of \cref{thm:halfspace-UC} hold, so $Z=0$ on $I\times\R^2$, whence
$G=0$ for $s>0$.  Since $F$ is odd in $s$, it vanishes for almost every
$(t,r,s)\in I\times\R^2$.  The continuity $F\in C(I;L^2)$ upgrades this
almost-everywhere-in-time statement to $F=0$ in $C(I;L^2(\R^2))$.  As $I$ was
an arbitrary bounded open interval, $F(t)=0$ in $L^2(\R^2)$ for every
$t\in\R$.
% FULL PROOF: paper.tex L4106-L4142
\end{proof}

\section{The Hilbert-space exterior-product lemma}

\begin{lemma}[Zero exterior product means linear dependence, unimodular under equal norms]
\label{lem:Hilbert-wedge}
Let $f,g\in L^2(\R)$ and assume
\[
 f(x)g(y)-g(x)f(y)=0\quad\text{in }L^2(\R^2).
\]
If $f\ne0$, there is a unique $c\in\C$ with $g=cf$ in $L^2$.  If, in addition,
$\norm f_2=\norm g_2$, then either $f=g=0$ or $|c|=1$; equivalently, under the
equal-norm hypothesis there always exists $\zeta\in\C$ with $|\zeta|=1$ and
$g=\zeta f$.
\end{lemma}

\begin{proof}
Assume first $f\ne0$.  The contraction map
\[
 \mathcal C_f:L^2(\R^2)\to L^2(\R),\qquad
 (\mathcal C_fH)(y)=\int_\R\overline{f(x)}H(x,y)\,\dd x
\]
is bounded with norm at most $\norm f_2$ by Cauchy--Schwarz and Fubini.
Applying it to the hypothesis gives
\[
 \norm f_2^2\,g-\Bigl(\int_\R\overline f g\Bigr)f=0 ,
\]
so $g=cf$ with $c=\pair f g/\norm f_2^2$; uniqueness follows from $f\ne0$.  If
the norms are equal and nonzero, then $|c|\norm f_2=\norm g_2=\norm f_2$, hence
$|c|=1$.  If $f=0$ and the norms are equal, then $g=0$ as well and one may take
$\zeta=1$.
\end{proof}

\section{Proof of the main theorem}

\begin{theorem}[Main phase-retrieval theorem]\label{thm:main}
\uses{thm:GWP}
Let $u,v\in C(\R;L^2(\R))$ be the global mild solutions of the cubic NLS with
initial data $u_0,v_0\in L^2(\R)$.  Assume
\[
 |u(t,x)|=|v(t,x)|\quad\text{for almost every }(t,x)\in\R^2 .
\]
Then there is a scalar $\zeta\in\C$ with $|\zeta|=1$ such that
\[
 v(t)=\zeta u(t)\quad\text{in }L^2(\R)\quad\text{for every }t\in\R .
\]
Equivalently, after choosing measurable representatives,
$v(t,x)=\zeta u(t,x)$ for almost every $(t,x)$.
\end{theorem}

\begin{proof}
\uses{lem:F-vanishing, lem:Hilbert-wedge, thm:GWP, cor:all-time-density}
By \cref{lem:F-vanishing}, $F(t)=0$ in $L^2(\R^2)$ for every $t$; evaluate at
$t=0$.  If $u(0)=0$, then \cref{cor:all-time-density} gives $v(0)=0$, and
uniqueness of the global flow (\cref{thm:GWP}) shows $u=v=0$.  Assume
$u(0)\ne0$.  By \cref{lem:Hilbert-wedge} there is $\zeta\in\C$ with
$v(0)=\zeta u(0)$; equal densities imply equal $L^2$ norms
(\cref{cor:all-time-density}), so $|\zeta|=1$.  Gauge covariance and uniqueness
of the global mild flow (\cref{thm:GWP}) give
\[
 v(t)=\Phi_t(v(0))=\Phi_t(\zeta u(0))=\zeta\Phi_t(u(0))=\zeta u(t)
\]
for every $t\in\R$.
\end{proof}

\section{Immediate corollaries}

\begin{corollary}[Equality on any time interval]\label{cor:interval-version}
\uses{thm:GWP}
Let $u,v$ be global $L^2$ solutions.  If $|u|=|v|$ almost everywhere on
$J\times\R$ for one nonempty open time interval $J$, then there is $\zeta\in\C$,
$|\zeta|=1$, such that $v=\zeta u$ on all of $\R\times\R$.
\end{corollary}

\begin{proof}
\uses{thm:main, thm:GWP}
Repeating the proof of \cref{thm:main} on $J$ gives $F(t)=0$ for every $t\in J$,
hence $v(t_0)=\zeta u(t_0)$ at one $t_0\in J$.  Global uniqueness and gauge
covariance (\cref{thm:GWP}) propagate the relation both forward and backward.
\end{proof}

\begin{corollary}[The finite-energy case]\label{cor:H1-case}
The conclusion of \cref{thm:main} holds in particular for solutions in
$C(\R;H^1(\R))$; no separate one-dimensional lateral unique-continuation
theorem is required.
\end{corollary}

\begin{proof}
\uses{thm:main, lem:H1-persistence}
An $H^1$ solution is the $L^2$ mild solution with the same data, by persistence
(\cref{lem:H1-persistence}) and uniqueness, so \cref{thm:main} applies.
\end{proof}

\section{A variable-coefficient Gr\"onwall inequality and rigidity of a common potential}

\begin{lemma}[Variable-coefficient integral Gr\"onwall]\label{lem:integral-gronwall}
Let $a\in L^1_{\loc}(\R)$ with $a\ge0$, let $t_0\in\R$, and let
$y:\R\to[0,\infty)$ be locally bounded and measurable.  Suppose
\[
 y(t)\le\int_{t_0}^t a(s)\,y(s)\,\dd s\quad(t\ge t_0),
 \qquad
 y(t)\le\int_{t}^{t_0} a(s)\,y(s)\,\dd s\quad(t\le t_0).
\]
Then $y(t)=0$ for every $t\in\R$.
\end{lemma}

\begin{proof}
Consider $t\ge t_0$ (the case $t\le t_0$ is identical after reversing the
orientation).  Fix a compact interval $[t_0,T]$ and let
$M=\sup_{[t_0,T]}y<\infty$ and $A(t)=\int_{t_0}^t a$.  Iterating the hypothesis
$n$ times gives, for $t\in[t_0,T]$,
\[
 y(t)\le M\,\frac{A(t)^n}{n!}\le M\,\frac{A(T)^n}{n!}\xrightarrow[n\to\infty]{}0 ,
\]
so $y\equiv0$ on $[t_0,T]$.  Since $T$ was arbitrary, $y\equiv0$ on $[t_0,\infty)$,
and symmetrically on $(-\infty,t_0]$.
\end{proof}

\begin{corollary}[Rigidity of a common real potential]\label{cor:linear-common-potential}
Let $V\in L^1_{\loc}(\R;L^\infty(\R))$ be real, and suppose
$u,v\in C(\R;L^2)\cap L^4_{\loc}L^\infty$ are mild solutions of
\[
 iu_t+u_{xx}=Vu,\qquad iv_t+v_{xx}=Vv ,
\]
that have the same modulus almost everywhere and possess the local
half-derivative smoothing needed for the quadratic current.  Then $v=\zeta u$
for one constant unimodular $\zeta$.
\end{corollary}

\begin{proof}
\uses{thm:critical-wronskian, prop:zero-extension-jump, cor:halfspace-XY, lem:Hilbert-wedge, lem:integral-gronwall}
Because $V$ is real the continuity equation is unchanged, so the Wronskian
(\cref{thm:critical-wronskian}), exterior-product, trace
(\cref{prop:zero-extension-jump}), and half-space
(\cref{cor:halfspace-XY}) arguments apply verbatim and yield, at any fixed time
$t_0$, a scalar $\zeta\in\T$ with $v(t_0)=\zeta u(t_0)$ via
\cref{lem:Hilbert-wedge}.  The difference $w=v-\zeta u$ satisfies the mild
equation
\[
 w(t)=-i\int_{t_0}^t S(t-s)\,V(s)w(s)\,\dd s .
\]
Unitarity of $S(t)$ gives
\[
 \norm{w(t)}_2\le\Bigl|\int_{t_0}^t\norm{V(s)}_\infty\norm{w(s)}_2\,\dd s\Bigr| ,
\]
with the integration orientation adjusted when $t<t_0$.  Applying
\cref{lem:integral-gronwall} with $y(t)=\norm{w(t)}_2$ and
$a(s)=\norm{V(s)}_\infty\in L^1_{\loc}$ yields $w=0$ in both time directions,
i.e.\ $v=\zeta u$.
\end{proof}


\appendix

\chapter{Suggested split into Lean files}\label{app:lean-split}

This appendix proposes the Lean module layout for the eventual formalization. It records, for each planned
file, the blueprint results it should contain, its imports, and the Mathlib entry points it consumes. No Lean
code is given here; only the file-level organization. The order is a valid topological order of the
\verb|\uses| graph, so a file only imports files listed above it.

\section{Package layout}
A single library \texttt{CubicNLSPhaseRetrieval} with one Lean file per module below, plus a top-level
\texttt{CubicNLSPhaseRetrieval.lean} that imports them all. Pin \texttt{leanprover/lean4:v4.32.0} and Mathlib
\texttt{rev = "v4.32.0"} (commit \texttt{81a5d25}). Before writing any proof, run \verb|#check| on every Mathlib
entry point in \cref{sec:mathlib-entrypoints} to confirm the exact names at this revision.

\section{Files, contents, and dependencies}
\begin{longtable}{p{0.05\textwidth}p{0.24\textwidth}p{0.40\textwidth}p{0.22\textwidth}}
\toprule
\# & Lean file & Blueprint contents (chapter) & Imports (files) \\
\midrule\endhead
0  & \texttt{ScalarMixedNorms.lean} & \cref{chap:mixed-norms}: joint representatives, $L^p_tL^q_x$, completeness, endpoint norming, sum/intersection spaces & Mathlib \\
1  & \texttt{FourierSobolev.lean} & \cref{chap:fourier}: Fourier convention, $\langle D\rangle^s$/$|D|^s$, \textbf{$H^{s}$ Hilbert space, $H^{-1/2}$--$H^{1/2}$ pairing, $H^s_\loc$}, Fourier--Gagliardo, quantitative smooth multiplier, difference quotients, $H^1$ representative, tensor norm & Mathlib \\
2  & \texttt{FreeSchrodinger.lean} & \cref{chap:free-group}: $S(t)$, unitarity, complex Gaussian kernel, dispersive estimate, propagator compatibility & 1 \\
3  & \texttt{FractionalIntegration.lean} & \cref{chap:frac-int}: three-lines, Riesz--Thorin, Hardy--Littlewood maximal (weak/strong), Hedberg, HLS & 2 \\
4  & \texttt{Strichartz1D.lean} & \cref{chap:strichartz}: admissible pairs, $TT^*$, homogeneous \& inhomogeneous Strichartz, Christ--Kiselev, retarded operator & 0,2,3 \\
5  & \texttt{CubicFlow.lean} & \cref{chap:wellposed}: $X(I)$, cubic estimates, LWP, $H^1$ persistence, mass conservation, GWP + gauge covariance & 1,4 \\
6  & \texttt{EndpointRegularity.lean} & \cref{chap:endpoint}: $X_k(J)$, $H^k$ persistence, endpoint $L^4L^\infty$/$L^8L^4$, $|u|^2u\in L^1L^2$, $C^0$ slices & 1,4,5 \\
7  & \texttt{LocalSmoothing1D.lean} & \cref{chap:smoothing}: smoothing identity, $A_xf$, forced smoothing, Duhamel smoothing, good-time set & 1,2,4,6 \\
8  & \texttt{CriticalQuadraticForms.lean} & \cref{chap:critical-calculus}: product estimate, $T_\eps$, test algebra $A(K)$, current $\overline f f_x$, quadratic-form extension & 1,6 \\
9  & \texttt{DensityCurrent.lean} & \cref{chap:current}: all-time density, continuity equation, endpoint current, current recovery & 6,8,12 \\
10 & \texttt{WronskianExterior.lean} & \cref{chap:wronskian-exterior}: critical Wronskian, exterior product $F$, ridge equation, coordinates, $Q$-bounds & 7,8,9 \\
11 & \texttt{DiagonalTrace.lean} & \cref{chap:traces}: local $H^2\!\Rightarrow\!C^1$, scalarization, Dirichlet/Neumann traces, zero-extension jump & 1,10 \\
12 & \texttt{Utilities.lean} & \cref{chap:utilities}: tensor-test density, distribution slicing, distributional products & 1 \\
13 & \texttt{Carleman2D.lean} & \cref{chap:carleman}: $X,X',Y$, $P_\beta$, $T_\beta$, kernel bounds, linear Carleman estimate, graph closure, exponential separation & 0,2,3 \\
14 & \texttt{HalfspacePropagation.lean} & \cref{chap:halfspace}: interface scale, one-slab, one-step, half-space uniqueness & 13 \\
15 & \texttt{PhaseRetrieval.lean} & \cref{chap:completion}: vanishing $F$, Hilbert wedge, main theorem, interval/$H^1$ corollaries, integral Gr\"onwall, common-potential rigidity & 5,10,11,14 \\
\bottomrule
\end{longtable}

\noindent Build order (topological): $0,1,2,3,12,4,5,6,7,8,9,10,11,13,14,15$. Files $12$ and $13$ depend only on
$0$--$3$, so they may be built early and in parallel with $5$--$11$.

\section{Mathlib entry points to \texttt{\#check} first}\label{sec:mathlib-entrypoints}
All confirmed present at commit \texttt{81a5d25}. Confirm the exact spellings before use.
\begin{itemize}
\item \texttt{MeasureTheory.Lp.fourierTransform}$_{\ell i}$ (the $L^2$ Fourier isometry $\simeq_{\ell i}$),
      \texttt{MeasureTheory.Lp.norm\_fourier\_eq}, \texttt{MeasureTheory.Lp.inner\_fourier\_eq},
      \texttt{MeasureTheory.Lp.fourier\_toTemperedDistribution\_eq}  \quad(file \texttt{Analysis/Fourier/LpSpace.lean}).
\item \texttt{TemperedDistribution.MemSobolev} and \texttt{MemSobolev.lineDerivOp},
      \texttt{MemSobolev.laplacian} (both at $p=2$), \texttt{besselPotential}
      \quad(file \texttt{Analysis/Distribution/Sobolev.lean}).
      \textbf{Naming nuance:} the $p=2$ Fourier characterization is \emph{not} named
      \texttt{memSobolev\_two\_iff\_fourier} (that string is docstring-only); the real lemmas are
      \texttt{memSobolev\_iff\_exists\_smulLeftCLM\_fourier} and \texttt{memSobolev\_besselPotential\_iff}.
      \verb|#check| and alias whichever you use.
\item Schwartz/tempered-distribution stack: \texttt{Analysis/Distribution/\{SchwartzSpace,TemperedDistribution,
      FourierMultiplier,FourierSchwartz\}.lean}.
\item \texttt{Analysis/SpecialFunctions/Gaussian/FourierTransform.lean} (free kernel);
      \texttt{Analysis/Complex/\{AbsMax,PhragmenLindelof\}.lean} (max-modulus / three-lines input);
      \texttt{Analysis/Distribution/AEEqOfIntegralContDiff.lean} (du~Bois-Reymond);
      \texttt{MeasureTheory/Integral/IntervalIntegral/\{Basic,FundThmCalculus,AbsolutelyContinuousFun\}.lean}.
\end{itemize}

\section{The one genuinely missing foundational leaf}
Mathlib's \texttt{Gronwall.lean} provides only the constant-coefficient derivative form. The
variable-coefficient \emph{integral} Gr\"onwall inequality needed by \cref{cor:linear-common-potential} is stated
here as \cref{lem:integral-gronwall}; it must be proved from the fundamental theorem of calculus (a short lemma).
It is the only foundational fact used in this blueprint that is not already in Mathlib, and it lies off the main
path (\cref{thm:main} does not use it).

\section{Where to start, and where the risk is}
\begin{itemize}
\item \textbf{Start} with file~0 (\texttt{ScalarMixedNorms}) and file~1 (\texttt{FourierSobolev}). File~1 is the
      highest-leverage design task: fix the $H^{s}$ interface (bundled Hilbert space via
      $\langle\xi\rangle^s\widehat f\in L^2$, \cref{def:Hs-space}), the \emph{complex-bilinear}
      $H^{-1/2}$--$H^{1/2}$ pairing (\cref{def:Hs-neg-pairing}, extending $\int T\varphi$, not the sesquilinear
      inner product), and $H^s_\loc$ (\cref{def:Hs-loc}) before proving anything downstream. A good first proof
      to formalize as a warm-up is \cref{lem:Hilbert-wedge} (self-contained: Cauchy--Schwarz $+$ Fubini).
\item \textbf{Hardest / highest residual-risk} files, in order: \texttt{Carleman2D} (file~13), \texttt{DiagonalTrace}
      (file~11), \texttt{CriticalQuadraticForms} (file~8). Budget extra verification there; the blueprint splits
      their proofs into the sub-lemmas most likely to need care (kernel bounds, the quantitative multiplier
      \cref{lem:translated-test-multiplier}, and the $A(K)$/current definitions).
\end{itemize}

\backmatter
\end{document}
